%==============================================================================%
% PhD Dissertation: Formal Verification of the Number-Theoretic Transform
%                   in the HAETAE Signature Scheme
% Author: [Author Name]
% Institution: [Institution Name]
% Year: 2026
%==============================================================================%

\documentclass[12pt, a4paper, twoside, openright]{book}

%--- Encoding & Fonts ---------------------------------------------------------
\usepackage[T1]{fontenc}
\usepackage[utf8]{inputenc}
\usepackage{lmodern}
\usepackage{microtype}

%--- Mathematics --------------------------------------------------------------
\usepackage{amsmath, amssymb, amsthm, mathtools}
\usepackage{bm}
\usepackage{bbm}
\usepackage{stmaryrd}   % \llbracket, \rrbracket

%--- Layout & Typography ------------------------------------------------------
\usepackage[top=3cm, bottom=3cm, left=3.5cm, right=2.5cm]{geometry}
\usepackage{setspace}
\onehalfspacing
\usepackage{emptypage}
\usepackage{fancyhdr}

%--- Colors (must come before hyperref) ---------------------------------------
\usepackage[dvipsnames, table]{xcolor}
\definecolor{codebg}{RGB}{248,248,248}
\definecolor{keyword}{RGB}{0,0,180}
\definecolor{comment}{RGB}{60,120,60}
\definecolor{string}{RGB}{160,0,0}
\definecolor{number}{RGB}{100,60,0}

%--- References & Hyperlinks --------------------------------------------------
\usepackage[backend=bibtex, style=alphabetic, sorting=anyt, maxnames=4]{biblatex}
\addbibresource{bibliography/refs.bib}
\usepackage[colorlinks=true,
            linkcolor=NavyBlue,
            citecolor=OliveGreen,
            urlcolor=BrickRed]{hyperref}
\usepackage{cleveref}

%--- Listings & Code ----------------------------------------------------------
\usepackage{listings}
\usepackage{algorithm}
\usepackage{algpseudocode}

%--- Graphics & Tables --------------------------------------------------------
\usepackage{graphicx}
\usepackage{booktabs}
\usepackage{multirow}
\usepackage{array}
\usepackage{longtable}
\usepackage{tikz}
\usetikzlibrary{
  arrows.meta, positioning, shapes.geometric, shapes.multipart,
  fit, calc, matrix, backgrounds, patterns,
  decorations.pathreplacing, decorations.markings,
  shadows.blur,
}
\usepackage{ifthen}

%--- Theorem Environments -----------------------------------------------------
\theoremstyle{plain}
\newtheorem{theorem}{Theorem}[chapter]
\newtheorem{lemma}[theorem]{Lemma}
\newtheorem{proposition}[theorem]{Proposition}
\newtheorem{corollary}[theorem]{Corollary}

\theoremstyle{definition}
\newtheorem{definition}[theorem]{Definition}
\newtheorem{example}[theorem]{Example}
\newtheorem{remark}[theorem]{Remark}
\newtheorem{notation}[theorem]{Notation}

%--- Code Style ---------------------------------------------------------------
\lstdefinelanguage{EasyCrypt}{
  keywords={require, import, theory, end, op, type, abbrev, lemma, proof,
            done, qed, by, forall, exists, fun, if, then, else, match, with,
            let, in, return, proc, while, do, module, abstract, clone, as,
            include, export, local, global, print, axiom, hint, rewrite,
            apply, exact, split, left, right, intro, move, have, case,
            elim, congr, ring, field, smt, auto, trivial, assumption},
  morecomment=[l]{(*},
  morecomment=[s]{(*}{*)},
  morestring=[b]",
  sensitive=true,
}

\lstdefinelanguage{Jasmin}{
  keywords={fn, export, inline, reg, stack, ptr, u8, u16, u32, u64, u128,
            u256, bool, int, return, if, else, while, for, to, downto,
            param, global, mut},
  morecomment=[l]{//},
  morecomment=[s]{/*}{*/},
  sensitive=true,
}

\lstset{
  basicstyle=\ttfamily\small,
  backgroundcolor=\color{codebg},
  keywordstyle=\color{keyword}\bfseries,
  commentstyle=\color{comment}\itshape,
  stringstyle=\color{string},
  numberstyle=\tiny\color{gray},
  numbers=left,
  numbersep=8pt,
  frame=single,
  framesep=4pt,
  breaklines=true,
  breakatwhitespace=false,
  tabsize=2,
  captionpos=b,
  xleftmargin=15pt,
  xrightmargin=5pt,
}

%--- Custom Commands ----------------------------------------------------------
\newcommand{\Z}{\ensuremath{\mathbb{Z}}}
\newcommand{\Zq}{\ensuremath{\mathbb{Z}_q}}
\newcommand{\Rq}{\ensuremath{\mathbb{Z}_q[X]/(X^{256}+1)}}
\newcommand{\F}{\ensuremath{\mathbb{F}}}
\newcommand{\Fq}{\ensuremath{\mathbb{F}_q}}
\newcommand{\NTT}{\ensuremath{\mathsf{NTT}}}
\newcommand{\INTT}{\ensuremath{\mathsf{INTT}}}
\newcommand{\brev}{\ensuremath{\mathsf{brev}}}
\newcommand{\mont}{\ensuremath{\mathsf{mont}}}
\newcommand{\MONT}{\ensuremath{\mathsf{MONT}}}
\newcommand{\qinv}{\ensuremath{\mathsf{QINV}}}
\newcommand{\haetae}{\textsc{Haetae}}
\newcommand{\jasmin}{\textsf{Jasmin}}
\newcommand{\easycrypt}{\textsf{EasyCrypt}}
\newcommand{\ec}[1]{\texttt{#1}}
\newcommand{\jasm}[1]{\texttt{#1}}
\newcommand{\op}[1]{\ensuremath{\mathsf{#1}}}
\newcommand{\abs}[1]{\ensuremath{\left|#1\right|}}
\newcommand{\norm}[1]{\ensuremath{\left\|#1\right\|}}
\newcommand{\floor}[1]{\ensuremath{\left\lfloor #1 \right\rfloor}}
\newcommand{\ceil}[1]{\ensuremath{\left\lceil #1 \right\rceil}}
\newcommand{\modulo}{\ensuremath{\mathbin{\%}}}
\newcommand{\pRHL}{\texorpdfstring{\ensuremath{\mathsf{pRHL}}}{pRHL}}
\newcommand{\HL}{\texorpdfstring{\ensuremath{\mathsf{HL}}}{HL}}

%--- Header / Footer ----------------------------------------------------------
\pagestyle{fancy}
\fancyhf{}
\fancyhead[LE]{\small\itshape\leftmark}
\fancyhead[RO]{\small\itshape\rightmark}
\fancyfoot[C]{\thepage}
\renewcommand{\headrulewidth}{0.4pt}

%==============================================================================%
\begin{document}
%==============================================================================%
\setlength{\parskip}{0.5\baselineskip}
\setlength{\parindent}{0pt}

\frontmatter
%==============================================================================%
% Title Page
%==============================================================================%

\newpage
\thispagestyle{empty}

\null\vspace{2cm}

{\centering\LARGE\bfseries
Formal Verification of the Number-Theoretic Transform\\[0.5em]
in the \haetae{} Post-Quantum Signature Scheme\par}

\vspace{1.5cm}

{\centering\large
A Dissertation Submitted in Partial Fulfillment\\
of the Requirements for the Degree of\\[0.5em]
\textbf{Doctor of Philosophy}\par}

\vspace{1.5cm}

{\centering\Large\itshape {Author Name}\par}

\vspace{0.5cm}

{\centering\large
Department of Computer Science and Engineering\\
{Institution Name}\par}

\vspace{1.5cm}

\begin{center}\rule{0.6\textwidth}{0.5pt}\end{center}

\vspace{0.5cm}

{\centering\large
Supervised by:\\[0.3em]
\textbf{{Supervisor Name}}\\
Professor, Department of Computer Science and Engineering\par}

\vspace{1.5cm}

\begin{center}\rule{0.6\textwidth}{0.5pt}\end{center}

\vspace{0.5cm}

{\centering\large 2026\par}

\vfill

{\centering\small
This dissertation is submitted under the license of
\href{https://creativecommons.org/licenses/by/4.0/}{CC BY 4.0 International}\par}

\cleardoublepage

%==============================================================================%
% Abstract
%==============================================================================%

\chapter*{Abstract}
\addcontentsline{toc}{chapter}{Abstract}

\noindent
The transition to post-quantum cryptography demands not only mathematically
sound security arguments but also \emph{machine-verified} correctness proofs for
concrete implementations.  This dissertation presents a comprehensive formal
verification of the Number-Theoretic Transform (NTT) at the core of the
\haetae{} lattice-based digital signature scheme---a Korean national standard
candidate under the KPQC competition.

\medskip
\noindent
We establish functional correctness across three vertically composed
abstraction layers.  At the \emph{implementation layer}, the high-assurance
\jasmin{} programming language produces x86-64 assembly that admits direct
extraction into \easycrypt{} modules.  At the \emph{algorithmic layer}, an
imperative specification in \easycrypt{} faithfully models the in-place
Cooley--Tukey butterfly structure together with Montgomery modular arithmetic
over $\Zq$, $q = 64513$, including tight coefficient-bound invariants across
all eight butterfly stages.  At the \emph{mathematical layer}, a fully abstract
spectral specification expresses the NTT as a Vandermonde-style evaluation at
powers of a primitive $512$th root of unity in $\Zq$.

\medskip
\noindent
The key technical contributions are:
\begin{enumerate}
  \item A machine-checked \textbf{Montgomery reduction correctness} proof,
        showing that $\op{montgomery\_reduce}(a) \equiv a \cdot R^{-1}
        \pmod{q}$ for all $a$ in the $16$-bit to $32$-bit working range, with
        explicit arithmetic bit-width bounds.

  \item A \textbf{butterfly-stage induction} connecting the iterative
        forward NTT loop in \ec{NTT\_Fq.ec} to the closed-form evaluation
        formula in \ec{NTTFullSpec.ec}, parameterised over the 8-level
        Cooley--Tukey recursion for the negacyclic polynomial ring
        $\Z_q[X]/(X^{256}+1)$.

  \item An \textbf{inverse NTT correctness} proof that tracks the final
        normalisation by \(256^{-1}\) and the concrete Montgomery output
        convention, where the exported inverse transform returns coefficients
        represented with one factor of \(R\), as expected for
        \texttt{invntt\_tomont}.

  \item A \textbf{Jasmin-to-EasyCrypt refinement} proof in
        \ec{RefJasminNTT.ec} connecting the extracted Jasmin semantics to the
        imperative specifications, with a 16-bit-to-24-bit coefficient-bound
        analysis for the forward transform and a tighter 16-bit preservation
        proof for the inverse transform.

  \item An \textbf{end-to-end statement} in \ec{NTTEndToEnd.ec} asserting
        that the \jasmin{} exported functions \jasm{poly\_ntt\_jazz} and
        \jasm{poly\_invntt\_jazz} compute exactly the abstract spectral NTT and
        its inverse, through two top-level Hoare statements for the forward
        and inverse exported functions.
\end{enumerate}

\noindent
The verification identifies a critical structural difference between the
\haetae{} NTT and the MLKEM NTT: the former performs a \emph{complete}
$256$-point negacyclic transform with 8-bit bit-reversed twiddle indices,
whereas the latter applies an incomplete $128$-point transform with a
different exponent structure.
This difference invalidates a naive reuse of existing MLKEM algebra lemmas
and motivates new algebraic infrastructure developed in
\ec{NTTFullAlgebra.ec}.

\medskip
\noindent
The scalar NTT verification compiles in \easycrypt{} version~2024 on an
x86-64 Linux platform through the top-level \texttt{NTTEndToEnd.ec} file.
The checked development totals approximately 8\,000 lines of \easycrypt{}
and 245 lines of \jasmin{} support code.

\vspace{1em}
\noindent
\textbf{Keywords:} Number-Theoretic Transform, Formal Verification,
EasyCrypt, Jasmin, Post-Quantum Cryptography, HAETAE, Montgomery Arithmetic,
Lattice-Based Signatures, Functional Correctness, Negacyclic Convolution.

%==============================================================================%
% Acknowledgements
%==============================================================================%

\chapter*{Acknowledgements}
\addcontentsline{toc}{chapter}{Acknowledgements}

I am deeply grateful to my supervisor, \textbf{{Supervisor Name}}, whose
insight into formal methods and post-quantum cryptography shaped every aspect
of this work.  Their patience in guiding me through the intricacies of
\easycrypt{}'s pRHL calculus and the semantic model of \jasmin{} was
invaluable.

I thank the members of the \haetae{} design team---in particular the authors
of the KPQC submission---for open-sourcing the reference C implementation
and for helpful discussions about the mathematical structure of the negacyclic
NTT.

I am indebted to the \easycrypt{} development team for producing a tool that
makes large-scale cryptographic proofs tractable, and to the Jasmin compiler
team for the extraction infrastructure that forms the bridge between
implementation and specification.

Financial support was provided by {Funding Agency / Grant Number}.

\bigskip
\noindent
{City}, {Month} 2026
\hfill
\textit{{Author Name}}


\tableofcontents
\listoffigures
\listoftables
\lstlistoflistings

\mainmatter
%==============================================================================%
% Chapter 1 – Introduction
%==============================================================================%

\chapter{Introduction}
\label{ch:intro}

\section{Motivation}

The threat posed by large-scale quantum computers to today's public-key
infrastructure is no longer a theoretical curiosity.  Shor's algorithm reduces
the security of RSA, DSA, and elliptic-curve schemes to polynomial time once a
sufficiently powerful quantum computer is available~\cite{Shor1994}.  In
response, the U.S.\ National Institute of Standards and Technology (NIST) ran
a multi-year post-quantum cryptography standardisation project, selecting four
algorithms for standardisation in 2022--2024~\cite{NIST2022}, while the Korean
Cryptography Module Validation Programme (KPQC) is pursuing its own national
standards~\cite{KPQC2022}.

\haetae{}~\cite{HAETAE2022} is one of the KPQC signature finalists.  It is a
lattice-based scheme whose security reduces to the Decisional Module
Short Integer Solution (DMSIS) hardness assumption over module lattices.  Like
all modern lattice-based schemes, \haetae{} relies on efficient polynomial
multiplication modulo $X^{256}+1$ in the ring $\Z_q[X]/(X^{256}+1)$, which
in practice is implemented via the \emph{Number-Theoretic Transform} (NTT).
The NTT reduces an $O(n^2)$ schoolbook convolution to $O(n \log n)$ operations
and is the dominant computational primitive in the scheme.

For a cryptographic standard to be trustworthy, correctness proofs cannot stop
at pen-and-paper mathematics.  History has shown that hand-verified
implementations harbour subtle bugs: off-by-one errors in loop bounds,
incorrect twiddle factor tables, wrong modular arithmetic, or mishandled
Montgomery scaling can silently produce wrong outputs while passing standard
test vectors~\cite{BerBuc2017}.  This dissertation takes the position that
\emph{machine-checked} formal proofs, linking every line of executable code to
its abstract mathematical specification, are the appropriate gold standard for
cryptographic primitives in post-quantum standards.

\section{Problem Statement}

We address the following problem:

\begin{quote}
  \textit{Prove, in a machine-checkable formal system, that the \jasmin{}
  implementation of the NTT and inverse-NTT in \haetae{} is functionally
  correct with respect to the abstract mathematical specification of the
  negacyclic NTT over $\Z_q[X]/(X^{256}+1)$.}
\end{quote}

The challenge is to bridge three distinct abstraction levels:

\begin{enumerate}
  \item \textbf{Low-level word arithmetic.}  The \jasmin{} source is compiled
        to x86-64 assembly and works with 32-bit and 64-bit machine words.
        Correctness at this level means proving arithmetic identities modulo
        $2^{32}$ and establishing bounds (e.g., that an intermediate product
        fits in a 64-bit register without overflow).

  \item \textbf{Algorithmic (loop) level.}  The Cooley--Tukey butterfly
        algorithm is expressed as a nested three-loop structure with eight
        stages.  Correctness here means showing that after all loops terminate,
        the array holds the intended Fourier coefficients, using loop invariants
        expressed in terms of partial NTT computations.

  \item \textbf{Mathematical (spectral) level.}  The NTT is a linear map
        $\Z_q^{256} \to \Z_q^{256}$ defined by evaluation at specific powers
        of a primitive $512$th root of unity.  Correctness here means proving
        that the algorithmic output equals the evaluation formula, and that the
        inverse NTT is its exact algebraic inverse.
\end{enumerate}

Each layer requires different mathematical machinery: modular arithmetic
and bit-vector identities at level~1, loop invariants and array
relational assertions at level~2, and field theory with primitive roots at
level~3.  Composing all three in a single proof assistant is the central
technical challenge.

\section{Tools and Methodology}

We use two mature high-assurance tools:

\paragraph{\jasmin{}~\cite{Almeida2017}.}
A domain-specific programming language and compiler designed for
high-assurance cryptographic implementations.  \jasmin{} programs are compiled
to provably correct x86-64 assembly and simultaneously admit extraction into
the \easycrypt{} proof assistant as first-class program modules.  The language
provides explicit register and memory annotations that remove ambiguity from
the compiler's optimisation decisions.

\paragraph{\easycrypt{}~\cite{Barthe2011,Barthe2013}.}
A proof assistant built around the probabilistic Relational Hoare Logic
($\pRHL$) calculus, designed specifically for reasoning about cryptographic
constructions.  An \easycrypt{} \emph{program logic} allows Hoare-style
reasoning over both deterministic and probabilistic procedures.  The ambient
logic is a dependent type theory supporting classical higher-order reasoning,
enabling rich algebraic developments (ring theory, field theory, group theory)
within the same tool.

Our proof strategy follows the \emph{refinement} paradigm:
\begin{align*}
  \text{Jasmin} &\xrightarrow{\text{extraction}}
  \text{EC module (RefJasminNTT)} \xrightarrow{\text{pRHL equivalence}}
  \text{EC imperative spec (NTT\_Fq)} \\
  &\xrightarrow{\text{loop invariant}}
  \text{EC abstract spec (NTTFullSpec/Algebra)}
\end{align*}
Each arrow represents a machine-checked proof step; the composition yields
the end-to-end theorem stated in \ec{NTTEndToEnd.ec}.

\section{Main Contributions}

\begin{description}
  \item[\textbf{C1 -- Montgomery correctness.}]
    A machine-checked proof that \ec{montgomery\_reduce} computes
    $a \cdot R^{-1} \bmod q$ for all inputs in the relevant working range,
    with $R = 2^{32}$, $q = 64513$, explicitly verified bit-width bounds,
    and signed-integer arithmetic lemmas absent from the \easycrypt{} standard
    library (\Cref{ch:montgomery}).

  \item[\textbf{C2 -- Imperative NTT specification.}]
    A fully proved \easycrypt{} model of the forward and inverse NTT loops in
    \ec{NTT\_Fq.ec}, including twiddle-factor correctness, coefficient-bound
    invariants (16-bit input $\to$ 24-bit output for forward NTT; 16-bit
    preservation for inverse NTT), and termination witnesses
    (\Cref{ch:ntt_spec}).

  \item[\textbf{C3 -- Abstract spectral specification.}]
    An algebraic library \ec{NTTFullAlgebra.ec} establishing properties of the
    primitive root $\zeta_0 = 426$ in $\Zq$, bit-reversal permutation lemmas,
    and a butterfly-stage induction connecting the iterative algorithm to the
    closed-form evaluation formula (\Cref{ch:algebra}).

  \item[\textbf{C4 -- Jasmin refinement.}]
    A pRHL proof in \ec{RefJasminNTT.ec} showing that the extracted \jasmin{}
    procedures are observationally equivalent to the imperative specifications
    up to array representation (\Cref{ch:refinement}).

  \item[\textbf{C5 -- End-to-end correctness.}]
    Top-level theorems \ec{haetae\_jasmin\_ntt\_correct} and
    \ec{haetae\_jasmin\_invntt\_correct} in \ec{NTTEndToEnd.ec}, together
    with the round-trip theorem $\INTT \circ \NTT = \op{id}$
    (\Cref{ch:end_to_end}).
\end{description}

\section{Structure of the Dissertation}

The remainder of this dissertation is organised as follows.

\begin{description}
  \item[\Cref{ch:background}] reviews the relevant background: lattice
        cryptography, the NTT, post-quantum signature schemes, and the
        theoretical foundations of formal verification with \easycrypt{} and
        \jasmin{}.

  \item[\Cref{ch:haetae_ntt}] describes the \haetae{} scheme, its NTT
        parameters, and the reference C implementation that serves as the
        algorithmic ground truth.

  \item[\Cref{ch:formal_framework}] presents the overall formal framework:
        the three-layer abstraction stack and the proof methodology.

  \item[\Cref{ch:field_ring}] covers the formalization of the field
        $\Zq$ and the polynomial ring $\Rq$ in \easycrypt{}.

  \item[\Cref{ch:montgomery}] details the Montgomery reduction
        formalization and its correctness proof.

  \item[\Cref{ch:ntt_spec}] presents the imperative NTT specification in
        \ec{NTT\_Fq.ec}, its invariants, and its termination proof.

  \item[\Cref{ch:algebra}] develops the algebraic infrastructure
        in \ec{NTTFullAlgebra.ec}: primitive root properties, bit-reversal
        lemmas, and stage decompositions.

  \item[\Cref{ch:jasmin_impl}] describes the \jasmin{} implementation in
        detail.

  \item[\Cref{ch:refinement}] presents the Jasmin-to-EasyCrypt refinement
        proof in \ec{RefJasminNTT.ec}.

  \item[\Cref{ch:end_to_end}] assembles all pieces into the end-to-end
        correctness theorems.

  \item[\Cref{ch:conclusion}] discusses limitations, remaining obligations,
        related and future work.
\end{description}

%==============================================================================%
% Chapter 2 – Background
%==============================================================================%

\chapter{Background}
\label{ch:background}

\section{Lattice-Based Cryptography}

\subsection{Module Lattices}

A \emph{lattice} $\Lambda \subseteq \Z^n$ is a discrete additive subgroup of
$\Z^n$; equivalently, it is the $\Z$-span of a basis matrix $\mathbf{B} \in
\Z^{n \times n}$.  The \emph{shortest vector problem} (SVP) asks for the
shortest non-zero lattice vector under the Euclidean norm, and is conjectured
to be hard for quantum computers.

Module lattices generalise to modules over a polynomial ring.  Let $R = \Z[X]
/ (X^n + 1)$ and $R_q = R/qR$.  A \emph{module lattice} of rank $k$ over $R$
is a free $R$-module $M \cong R^k$ together with a $\Z$-embedding into $\Z^{kn}$.
The \emph{Module Short Integer Solution} (MSIS) and \emph{Module Learning With
Errors} (MLWE) problems over such lattices underpin the security of Dilithium,
FALCON, and \haetae{}.

\subsection{NTT-Friendly Parameters}

For efficient computation, the ring $R_q$ must support a fast polynomial
multiplication algorithm.  If $q \equiv 1 \pmod{2n}$ then $X^n + 1$ splits
completely over $\F_q$ and the NTT reduces polynomial multiplication to
componentwise multiplication of length-$n$ vectors.  If only $q \equiv 1
\pmod{n}$ (but $q \equiv 3 \pmod{4}$ or similar), a \emph{negacyclic} NTT can
still be defined~\cite{Longa2016}.

For \haetae{}, $n = 256$ and $q = 64513 \equiv 1 \pmod{512}$.  This means a
primitive $512$th root of unity $\zeta_0$ exists in $\F_q$, enabling a full
$256$-point negacyclic NTT.

\section{Number-Theoretic Transform}

\subsection{Definition}

Let $q$ be an odd prime with $q \equiv 1 \pmod{2n}$, and let $\zeta \in \F_q^*$
be a primitive $2n$th root of unity (so $\zeta^{2n} = 1$, $\zeta^n = -1$).
The \emph{negacyclic NTT} of a polynomial $f = \sum_{j=0}^{n-1} f_j X^j \in
R_q$ is the vector $\hat{f} \in \F_q^n$ defined by
\begin{equation}
  \hat{f}_i = \sum_{j=0}^{n-1} f_j \cdot \zeta^{(2\brev_8(i)+1) \cdot j},
  \qquad i = 0, \dots, n-1,
\label{eq:ntt_def}
\end{equation}
where $\brev_8 : \{0,\dots,255\} \to \{0,\dots,255\}$ is the bit-reversal
permutation of an 8-bit index.

The inverse NTT is
\begin{equation}
  f_i = \frac{1}{n} \sum_{j=0}^{n-1} \hat{f}_j \cdot \zeta^{-((2\brev_8(j)+1) \cdot i)},
  \qquad i = 0, \dots, n-1.
\label{eq:intt_def}
\end{equation}

\begin{proposition}[Round-trip identity]
  $\INTT(\NTT(f)) = f$ for all $f \in R_q$.
\label{prop:roundtrip}
\end{proposition}

\subsection{Cooley--Tukey Butterfly Algorithm}
\label{sec:ct_butterfly}

The NTT defined by \cref{eq:ntt_def} can be computed in $O(n \log n)$ field
operations using the Cooley--Tukey decimation-in-time (DIT) algorithm.  For
the negacyclic case with $n = 256 = 2^8$, the algorithm proceeds in 8 stages.

At stage $s$ (with $s = 0, 1, \dots, 7$, from outermost to innermost), the
half-length is $\ell = 2^{7-s}$.  If $g$ is the group number within that
stage, then the implementation's table counter is $k = 2^s + g$ and the
twiddle factor is
\[
  w = \zeta^{\brev_8(k)}.
\]
Index 0 of the concrete forward table is padding, so the loop consumes
indices $1,\dots,255$.  The butterfly operation on a pair $(a, b)$ with
twiddle factor $w$ is:
\begin{equation}
  (a', b') = (a + w \cdot b,\ a - w \cdot b).
\label{eq:butterfly}
\end{equation}

\begin{algorithm}[htbp]
\caption{Forward NTT (Cooley--Tukey DIT, negacyclic, $n = 256$)}
\label{alg:ntt}
\begin{algorithmic}[1]
\Procedure{NTT}{$a[0..255]$}
\State $k \gets 1$
\For{$\ell = 128, 64, 32, 16, 8, 4, 2, 1$}
  \For{$\mathit{start} = 0, 2\ell, 4\ell, \dots, 254$}
    \State $w \gets \zeta^{\brev_8(k)}$ \Comment{$k$-th twiddle factor}
    \State $k \gets k + 1$
    \For{$j = \mathit{start}, \dots, \mathit{start}+\ell-1$}
      \State $t \gets w \cdot a[j+\ell] \bmod q$
      \State $a[j+\ell] \gets a[j] - t$
      \State $a[j] \gets a[j] + t$
    \EndFor
  \EndFor
\EndFor
\EndProcedure
\end{algorithmic}
\end{algorithm}

\subsection{Negacyclic NTT vs.\ Standard NTT}

A standard (cyclic) DFT of length $n$ uses the recurrence for the cyclic ring
$\Z_q[X]/(X^n - 1)$.  The \emph{negacyclic} ring $\Z_q[X]/(X^n + 1)$ requires
a pre-twist by powers of $\zeta$ (a half-integer shift of the exponent), which
accounts for the $+1$ in the exponent formula of \cref{eq:ntt_def}.

Critically for this dissertation, \haetae{} uses a \emph{complete} $n=256$
negacyclic NTT with bit-reversal exponent $\brev_8(k)$, whereas MLKEM uses an
incomplete $128$-point transform with exponent $\brev_7(k)$.  This distinction
has direct consequences for which algebraic lemmas can be shared between
verification projects.

\section{Montgomery Modular Arithmetic}

\subsection{REDC Algorithm}

Montgomery reduction~\cite{Montgomery1985} is a technique for efficient modular
multiplication that avoids expensive division.  Given $a \in \Z$ with
$0 \le a < q \cdot R$ and integers $q, R$ with $\gcd(q, R) = 1$, the
Montgomery product is $a \cdot R^{-1} \bmod q$.

\begin{algorithm}[htbp]
\caption{Montgomery Reduction with the \haetae{} Sign Convention}
\label{alg:mont_reduce}
\begin{algorithmic}[1]
\Require $a, R = 2^{32}$, $q^+ = q^{-1} \bmod R$
\Ensure $a \cdot R^{-1} \bmod q$
\State $m \gets (a \bmod R) \cdot q^+ \bmod R$
\State $t \gets (a - m \cdot q) \div R$
\If{$t < 0$} \Return $t + q$ \Else\ \Return $t$ \EndIf
\end{algorithmic}
\end{algorithm}

In the \haetae{} implementation, $R = 2^{32}$, $q = 64513$,
$\mathsf{QINV} = 940508161$ ($\equiv q^{-1} \bmod 2^{32}$), and
$\mathsf{MONT} = R \bmod q = 14321$.  Thus the implementation uses the
positive-inverse convention in \Cref{alg:mont_reduce}: it computes
$m = (a \bmod R)\mathsf{QINV} \bmod R$ and then forms $a - m q$.

\subsection{Montgomery Multiplication}

Given $a, b \in \Z_q$ with both represented in Montgomery form (i.e., as
$a \cdot R \bmod q$ and $b \cdot R \bmod q$), their product in Montgomery form
is computed as $\mathrm{REDC}(a \cdot b)$, since
$\mathrm{REDC}(aR \cdot bR) = ab R^2 \cdot R^{-1} = abR \bmod q$.

\subsection{Signed Montgomery Reduction}

The \haetae{} implementation uses a \emph{signed} variant where intermediate
values are 32-bit signed integers and the final conditional subtraction is
replaced by an arithmetic right shift.  Specifically, for $a \in \Z$ with
$\abs{a} < q \cdot 2^{15}$:
\begin{equation}
  a \;-\; \text{int32}(a_{\text{lo}} \cdot \mathsf{QINV}) \cdot q
  \quad\text{(as signed 64-bit)}
  \;\text{ shifted right by 32}
\label{eq:signed_mont}
\end{equation}
yields a value congruent to $a \cdot R^{-1} \pmod{q}$ and, in the checked
signed REDC theorem, a representative in $[-q,q)$.

\section{The EasyCrypt Proof Assistant}

\subsection{Architecture}

\easycrypt{}~\cite{Barthe2011} is an interactive proof assistant specialised
for reasoning about cryptographic programs.  It combines:
\begin{itemize}
  \item A \emph{program logic} layer supporting Hoare logic (HL) for
        deterministic programs and probabilistic relational Hoare logic
        (pRHL) for pairs of probabilistic programs.

  \item An \emph{ambient logic} layer that is a form of higher-order logic
        (HOL) with a type theory supporting algebraic structures (groups,
        rings, fields).

  \item A tactic engine with automation backends including SMT solvers
        (CVC4, Z3) and a computer algebra system for ring/field equalities.
\end{itemize}

\subsection{Hoare Logic in EasyCrypt}

A \emph{Hoare triple} $\{P\}\; c\; \{Q\}$ asserts that if precondition $P$
holds before executing command $c$, then postcondition $Q$ holds after.  In
\easycrypt{}, Hoare triples over \emph{procedures} (modules with state) are
written \ec{hoare[M.f : P ==> Q]}.

A \emph{relational Hoare triple} $\{P\}\; (c_1 \sim c_2)\; \{Q\}$ asserts
that if $c_1$ and $c_2$ start in states satisfying $P$, their outputs satisfy
$Q$.  This is the foundation for proving that two programs are equivalent.

\subsection{Module System}

\easycrypt{} organises code into \emph{modules}.  A module interface (type) is
a signature listing procedure names and types; a module is a concrete
implementation.  The \jasmin{} compiler generates an \easycrypt{} module for
each exported function, which can then be related to hand-written specification
modules via pRHL.

\subsection{Theory Cloning}

Algebraic structures are developed as parameterised \emph{theories} that are
instantiated by \emph{cloning}: \ec{clone Ring.ComRing as Zq with type t = int, ...}
creates a copy of the commutative ring theory with $\Z_q$ as the carrier type
and appropriate witnesses for the ring axioms.

\section{The Jasmin Programming Language}

\jasmin{}~\cite{Almeida2017} is a C-like language designed for high-assurance
cryptographic code.  Its distinguishing features are:

\begin{itemize}
  \item \textbf{Explicit register allocation.}  Variables are annotated
        \jasm{reg} (in a register), \jasm{stack} (on the stack), or
        \jasm{ptr} (a pointer).  The compiler enforces the annotation.

  \item \textbf{Data types matching machine types.}  The types
        \jasm{u32}, \jasm{u64}, \jasm{u256} correspond directly to
        32-bit, 64-bit, and 256-bit (AVX2) machine words.

  \item \textbf{Safety verification.}  A safety checker ensures no
        undefined behaviour (array out of bounds, use before initialisation,
        integer overflow in unsigned arithmetic).

  \item \textbf{Constant-time verification.}  The companion
        \texttt{jasmin-ct} checker verifies that secret values do not control
        branches or memory addresses.  Its speculative mode also tracks
        transient pointer inputs, using explicit \jasm{\#[ct = "..."]} and
        \jasm{\#[sct = "..."]} contracts on exported procedures.

  \item \textbf{EasyCrypt extraction.}  Each \jasm{export fn} is compiled
        to both assembly and an \easycrypt{} module.  The extracted module has
        a semantic identical to the assembly, by the Jasmin compiler
        correctness theorem.
\end{itemize}

\begin{lstlisting}[language=Jasmin, caption={Jasmin function signature example}]
export fn poly_ntt_jazz(reg mut ptr u32[256] rp)
  -> reg ptr u32[256]
{
  rp = _poly_ntt(rp);
  return rp;
}
\end{lstlisting}

The export annotation triggers both assembly generation and \easycrypt{}
extraction.  The extracted \easycrypt{} type for \jasm{ptr u32[256]} is a
\ec{BArray1024.t} byte-array value, with coefficients accessed by
\ec{BArray1024.get32} and \ec{BArray1024.set32}.

\section{Related Work}

\paragraph{MLKEM/Kyber NTT.}
The closest predecessor is the Jasmin/EasyCrypt verification of the NTT in
ML-KEM~\cite{Barbosa2023}.  That work verifies the incomplete $128$-point
Cooley--Tukey NTT over a related ring structure but with an incomplete
transform and a different twiddle-indexing discipline.  Our work adapts
techniques from that development but requires new algebra lemmas for the
full $256$-point transform.

\paragraph{Dilithium NTT.}
The EasyCrypt verification of Dilithium~\cite{Barbosa2021} also formalises
an NTT in $\Z_q[X]/(X^{256}+1)$ but uses a different prime $q = 8380417$ with
different twiddle factor structure.

\paragraph{FALCON.}
FALCON uses a floating-point NTT over $\Z[X]/(X^n+1)$~\cite{Fouque2018}.
Its Jasmin implementation has been verified using \easycrypt{} but the
field arithmetic is quite different from the integer case studied here.

\paragraph{High-assurance cryptography.}
The HACL* project~\cite{Zinzindohoue2017} produces C code verified in
F*/Low*, with machine-checked compilation to assembly.  While HACL* has
verified ChaCha20 and Poly1305, it has not (at time of writing) produced a
verified lattice-based NTT.  Our work contributes to the complementary
Jasmin/EasyCrypt ecosystem, which has better support for the bitwise
operations prevalent in lattice primitives.

%==============================================================================%
% Chapter 3 – HAETAE and Its NTT
%==============================================================================%

\chapter{The \haetae{} Signature Scheme and Its NTT}
\label{ch:haetae_ntt}

\section{Overview of \haetae{}}

\haetae{}~\cite{HAETAE2022} is a module-lattice based digital signature scheme
submitted to the Korean Post-Quantum Cryptography (KPQC) standardisation
process.  It follows the ``Fiat--Shamir with aborts'' paradigm~\cite{Lyubashevsky2012}:
the signing algorithm samples a masking polynomial, computes a commitment via a
hash function, and rejects the sample if the output leaks information about the
secret key.

\subsection{Algebraic Setting}

\haetae{} operates in the polynomial ring
\begin{equation}
  R_q = \Z_q[X]/(X^n + 1), \qquad n = 256,\quad q = 64513.
\label{eq:ring}
\end{equation}
Keys and signatures are vectors in $R_q^k$ for parameter-set-dependent
ranks $k$.  Polynomial multiplication in $R_q$ is the dominant operation in
key generation, signing, and verification.

\subsection{Role of the NTT}

The NTT converts polynomials from the \emph{coefficient domain} to the
\emph{evaluation domain}, where multiplication is pointwise.  The workflow is:
\begin{equation}
  f \cdot g \;\in\; R_q
  \;\longleftrightarrow\;
  \NTT(f) \odot \NTT(g) \;\in\; \F_q^n,
\label{eq:ntt_mul}
\end{equation}
where $\odot$ denotes componentwise multiplication.  A single polynomial
multiplication thus costs two forward NTTs, one componentwise multiplication,
and one inverse NTT, for overall $O(n \log n)$ field operations instead of
$O(n^2)$.

\section{\haetae{} NTT Parameters}
\label{sec:haetae_params}

\begin{table}[htbp]
\centering
\caption{\haetae{} NTT parameters}
\label{tab:params}
\begin{tabular}{lll}
\toprule
Symbol & Value & Description \\
\midrule
$q$    & $64513$ & Modulus ($= 2^8 \cdot 252 + 1 = 252 \cdot 256 + 1$) \\
$n$    & $256$   & Polynomial degree \\
$\zeta_0$ & $426$ & Primitive $512$th root of unity in $\F_q$ \\
$R$    & $2^{32}$ & Montgomery base \\
$\mathsf{QINV}$ & $940508161$ & $q^{-1} \bmod 2^{32}$ \\
$\mathsf{MONT}$ & $14321$ & $R \bmod q = 2^{32} \bmod q$ \\
$\mathsf{MONTSQ}$ & $4214$ & $R^2 \bmod q = 2^{64} \bmod q$ \\
$R^{-1}$ & $50386$ & $R^{-1} \bmod q$ \\
$f$ & $-29720$ & Inverse NTT scaling constant (in Montgomery form) \\
\bottomrule
\end{tabular}
\end{table}

\subsection{Twiddle Factors}

The forward NTT uses 256 table slots in the reference array
\texttt{zetas[256]}.  Slot 0 is a padding value and the forward loop reads
indices $1,\dots,255$.  For $1 \le k < 256$, the concrete table entry is the
signed 32-bit representative of
$\mathsf{MONT} \cdot \zeta_0^{\brev_8(k)} \bmod q$.

\begin{table}[htbp]
\centering
\caption{First sixteen forward twiddle factors \texttt{zetas[1..16]} (signed Montgomery form)}
\label{tab:zetas}
\begin{tabular}{r|rrrrrrrr}
\toprule
$k$ & 1 & 2 & 3 & 4 & 5 & 6 & 7 & 8 \\
\midrule
\texttt{zetas}$[k]$ & 26964 & -16505 & 22229 & 30746 & 20243 & 19064 & -31218 & 9395 \\
\midrule
$k$ & 9 & 10 & 11 & 12 & 13 & 14 & 15 & 16 \\
\midrule
\texttt{zetas}$[k]$ & -30985 & 22859 & -8851 & 32144 & 13744 & 21408 & 17599 & -16039 \\
\bottomrule
\end{tabular}
\end{table}

The C inverse NTT reads the forward table in reverse order as
\texttt{-zetas[--k]}.  The \jasmin{} extraction stores the same inverse
butterfly constants explicitly as \texttt{jzetas\_inv[0..254]} and reserves
\texttt{jzetas\_inv[255]} for the final scaling constant
$f = n^{-1} \cdot \mathsf{MONT}^2 \bmod q = -29720$.

\section{Reference C Implementation}
\label{sec:c_impl}

The reference C implementation in \texttt{ntt.c} provides the algorithmic
ground truth.  It implements the forward NTT (\Cref{alg:ntt}) and the inverse
NTT with an extra final scaling step.

\subsection{Forward NTT}

\begin{lstlisting}[language=C,
  caption={Forward NTT from \texttt{ntt.c}},
  label={lst:c_ntt}]
void ntt(int32_t a[256]) {
  unsigned int len, start, j, k;
  int32_t t, zeta;

  k = 1;
  for (len = 128; len > 0; len >>= 1) {
    for (start = 0; start < 256; start += 2 * len) {
      zeta = zetas[k++];
      for (j = start; j < start + len; j++) {
        t = montgomery_reduce((int64_t)zeta * a[j + len]);
        a[j + len] = a[j] - t;
        a[j]       = a[j] + t;
      }
    }
  }
}
\end{lstlisting}

The loop nest has three levels:
\begin{enumerate}
  \item \textbf{Outer loop}: $\ell \in \{128, 64, 32, 16, 8, 4, 2, 1\}$,
        the current butterfly half-length.  Each iteration is one NTT stage.
  \item \textbf{Middle loop}: iterates over the $256 / (2\ell)$ groups within
        each stage.
  \item \textbf{Inner loop}: performs $\ell$ butterfly operations within each
        group, using the same twiddle factor \texttt{zeta = zetas[k]}.
\end{enumerate}

Key observations:
\begin{itemize}
  \item After stage $s$ (with $\ell = 2^{7-s}$), the array is arranged so
        that each processed segment represents the checked
        \ec{partial\_ntt} sum for the current segment length.  The formal
        invariant is indexed by \((\ell,\mathit{start},\mathit{bsj})\), not
        by a congruence class modulo \(2^\ell\).

  \item The twiddle counter $k$ advances monotonically from 1 to 255,
        consuming all 255 non-trivial twiddle factors.

  \item No modular reduction is applied at the output of butterfly pairs.
        The checked proof tracks stage-dependent signed bit-width bounds:
        the forward refinement theorem assumes a 16-bit input bound and
        proves a 24-bit output bound.
\end{itemize}

\subsection{Inverse NTT}

\begin{lstlisting}[language=C,
  caption={Inverse NTT from \texttt{ntt.c}},
  label={lst:c_invntt}]
void invntt_tomont(int32_t a[256]) {
  unsigned int start, len, j, k;
  int32_t t, zeta;
  const int32_t f = -29720;   /* mont^2 / 256 mod q */

  k = 256;
  for (len = 1; len < 256; len <<= 1) {
    for (start = 0; start < 256; start = j + len) {
      zeta = -zetas[--k];
      for (j = start; j < start + len; j++) {
        t        = a[j];
        a[j]     = t + a[j + len];
        a[j+len] = t - a[j+len];
        a[j+len] = montgomery_reduce((int64_t)zeta * a[j+len]);
      }
    }
  }
  for (j = 0; j < 256; j++)
    a[j] = montgomery_reduce((int64_t)f * a[j]);
}
\end{lstlisting}

The inverse NTT reverses the butterfly ordering: $\ell$ increases from 1 to
128, and the C twiddle counter starts at 256 so that the first butterfly
uses \texttt{-zetas[255]}.  The final loop multiplies each coefficient by
$f = -29720 \equiv n^{-1} \cdot R^2 \pmod{q}$, with the Montgomery reduction
removing one factor of $R$ and leaving each coefficient multiplied by
$n^{-1} \cdot R \pmod{q}$ (i.e., in Montgomery form).

\subsection{Montgomery Reduction}

\begin{lstlisting}[language=C,
  caption={Montgomery reduction from \texttt{reduce.c}},
  label={lst:c_mont}]
int32_t montgomery_reduce(int64_t a) {
  int32_t t;
  t = (int32_t)a * (int32_t)QINV;
  t = (a - (int64_t)t * Q) >> 32;
  return t;
}
\end{lstlisting}

This implements the signed Montgomery reduction.  The cast \texttt{(int32\_t)a}
extracts the low 32 bits of $a$ (interpreted as a signed 32-bit integer), the
product \texttt{t * Q} is computed as a signed 64-bit value, and the right
shift by 32 is an arithmetic right shift on signed integers.

\subsection{Test Harness}

The file \texttt{test\_haetae\_ntt.c} provides a reference test harness that:
\begin{enumerate}
  \item Generates random polynomial inputs within the valid range.
  \item Compares the C forward NTT with the \jasmin{} forward NTT.
  \item Compares the C inverse NTT with the \jasmin{} inverse NTT.
  \item Checks the combined C and \jasmin{} forward--inverse pipelines agree.
\end{enumerate}

These tests provide empirical C/Jasmin agreement.  The round-trip claim used
for interpretation is the mathematical abstract inverse identity together
with the two separate EasyCrypt end-to-end theorems when the inverse theorem's
input-bound precondition is satisfied; it is not inferred from this test
harness alone, and it is not packaged as a standalone checked pipeline theorem.

\section{Structural Comparison with MLKEM NTT}
\label{sec:haetae_vs_mlkem}

\begin{table}[htbp]
\centering
\caption{Comparison of \haetae{} and MLKEM NTT structures}
\label{tab:haetae_mlkem}
\begin{tabular}{lll}
\toprule
Property & \haetae{} & MLKEM \\
\midrule
Modulus $q$ & $64513$ & $3329$ \\
Degree $n$ & $256$ & $256$ \\
NTT type & Full 256-point negacyclic & Incomplete 128-point \\
Bit-reversal exponent & $\brev_8(k)$ & $\brev_7(k)$ \\
Primitive root & $\zeta_0 = 426$ ($512$th) & $\zeta_0 = 17$ ($256$th) \\
Output permutation & Bit-reversed & Bit-reversed \\
Final inv.\ NTT step & Scale by $n^{-1} R \bmod q$ & Scale by $n^{-1} R \bmod q$ \\
\bottomrule
\end{tabular}
\end{table}

The key structural difference is that the \haetae{} NTT uses twiddle factors
$\zeta^{\brev_8(k)}$ where $k$ ranges from 1 to 255, producing a full 256-point
spectral evaluation.  MLKEM uses twiddle factors $\zeta^{\brev_7(2k+1)}$,
producing only a 128-point evaluation (the polynomial is split into even/odd
parts).  This means the \haetae{} NTT result vector contains evaluation at
all 256 primitive $512$th roots of unity, while the MLKEM result vector
contains evaluations at 128 primitive $256$th roots.

This structural difference invalidates any attempt to directly reuse MLKEM NTT
algebra lemmas for the \haetae{} verification.  New algebraic infrastructure
is required, developed in \ec{NTTFullAlgebra.ec} (\Cref{ch:algebra}).

%==============================================================================%
% Chapter 4 – Formal Verification Framework
%==============================================================================%

\chapter{Formal Verification Framework}
\label{ch:formal_framework}

\section{Overview of the Verification Architecture}

The verification of the \haetae{} NTT is structured as a
\emph{vertical refinement chain}: a sequence of proof steps, each relating
an implementation or specification at one abstraction level to a
specification at the next.  \Cref{fig:arch} illustrates the overall structure.

\begin{figure}[htbp]
\centering
\begin{tikzpicture}[
  box/.style = {rectangle, draw, rounded corners, minimum width=8cm,
                minimum height=1.2cm, align=center, font=\small},
  arrow/.style = {-Stealth, thick},
  label/.style = {font=\footnotesize\itshape, midway, right=3pt}
]
  \node[box, fill=red!10]   (jasmin) at (0, 6)
        {\textbf{Layer 0: Jasmin Source} \\ \texttt{poly.jinc, reduce.jinc, hpoly.jazz}};

  \node[box, fill=orange!10] (extract) at (0, 4.2)
        {\textbf{Layer 1: Extracted EasyCrypt Module} \\ \texttt{Hpoly\_extract.ec}};

  \node[box, fill=yellow!10] (impspec) at (0, 2.4)
        {\textbf{Layer 2: Imperative EasyCrypt Spec} \\ \texttt{NTT\_Fq.ec, Fq.ec}};

  \node[box, fill=green!10]  (algspec) at (0, 0.6)
        {\textbf{Layer 3: Algebraic/Spectral Spec} \\
         \texttt{NTTFullSpec.ec, NTTFullAlgebra.ec}};

  \draw[arrow] (jasmin) -- (extract)
    node[label] {Jasmin compiler extraction};
  \draw[arrow] (extract) -- (impspec)
    node[label] {pRHL equivalence (RefJasminNTT.ec)};
  \draw[arrow] (impspec) -- (algspec)
    node[label] {Loop invariant / induction (NTTFullAlgebra.ec)};

  \node[box, fill=blue!10, minimum width=4cm, minimum height=0.9cm]
    (endtoend) at (5.5, 3.3)
    {\textbf{End-to-End} \\ \texttt{NTTEndToEnd.ec}};

  \draw[arrow, dashed, bend right=20] (jasmin.east) to
    node[right=3pt, font=\footnotesize\itshape] {top-level theorem}
    (endtoend.north);
  \draw[arrow, dashed, bend left=20] (algspec.east) to (endtoend.south);
\end{tikzpicture}
\caption{Four-layer verification architecture for the \haetae{} NTT.
         Solid arrows are machine-checked proof steps; dashed arrows indicate
         the end-to-end composition.}
\label{fig:arch}
\end{figure}

\section{Layer 0: Jasmin Source}

The \jasmin{} source files implement the NTT algorithm in a style close to
assembly but with structured control flow.  The key files are:
\begin{itemize}
  \item \texttt{poly.jinc}: The main NTT and inverse-NTT kernels,
        implemented as \jasm{inline fn \_poly\_ntt} and
        \jasm{inline fn \_poly\_invntt}.

  \item \texttt{reduce.jinc}: The Montgomery reduction primitive
        \jasm{\_\_montgomery\_reduce} and the derived \jasm{\_\_fqmul}.

  \item \texttt{zetas.jinc}: The 256-element twiddle-factor table as a
        global constant.

  \item \texttt{hpoly.jazz}: The exported entry points
        \jasm{export fn poly\_ntt\_jazz} and
        \jasm{export fn poly\_invntt\_jazz}.
\end{itemize}

The \jasmin{} compiler verifies safety properties (no array-out-of-bounds, no
use-before-initialisation) and then simultaneously generates:
\begin{enumerate}
  \item x86-64 assembly (\texttt{hpoly.s}).
  \item An \easycrypt{} module \texttt{Hpoly\_extract.ec} that is
        semantically equivalent to the assembly by the Jasmin compiler
        correctness theorem~\cite{Almeida2017}.
\end{enumerate}

\section{Layer 1: Extracted EasyCrypt Module}

The extracted module \texttt{Hpoly\_extract.ec} (\texttt{extract/} folder)
contains \easycrypt{} procedure definitions that mirror the \jasmin{}
structure.  A fragment is shown in \Cref{lst:extract}.

\begin{lstlisting}[language=EasyCrypt,
  caption={Fragment of extracted procedure in Hpoly\_extract.ec},
  label={lst:extract}]
module M = {
  proc poly_ntt_jazz (rp : BArray1024.t) :
      BArray1024.t = {
    var zeta : W32.t;
    var t    : W32.t;
    var len  : int;
    var start: int;
    var j    : int;
    (* ... body mirrors Jasmin structure ... *)
    return rp;
  }
}.
\end{lstlisting}

This module is the ``ground truth'' for what the compiled assembly computes.
All subsequent proof obligations are stated with respect to this module.

\section{Layer 2: Imperative EasyCrypt Specification}

The file \texttt{NTT\_Fq.ec} contains hand-written \easycrypt{} procedures
that closely mirror the reference C implementation.  Unlike the extracted
module, these procedures work with abstract coefficient arrays over \(\F_q\)
(type \ec{coeff Array256.t}) rather than machine-word arrays.

The translation from machine words to field coefficients is captured by the
\emph{word-to-coefficient} relation \ec{poly\_repr}: a concrete
\ec{BArray1024.t} word array corresponds to a coefficient array when each
loaded 32-bit word maps to the same \(\F_q\) element.  The strengthened
\ec{poly\_repr\_bound} predicate also records the signed bit-width bound
required for word arithmetic.

The imperative specification exposes the full loop structure, making it
amenable to Hoare-style invariant reasoning.

\section{Layer 3: Abstract Spectral Specification}

The files \texttt{NTTFullSpec.ec} and \texttt{NTTFullAlgebra.ec} define the
NTT at the level of abstract algebra.  No loop structure is visible; the
NTT is defined as a linear evaluation map over $\F_q^{256}$.

\begin{definition}[Abstract forward NTT, \ec{NTTFullSpec.full\_ntt}]
  Let $p : \F_q^{256}$.  The forward NTT of $p$ is the map
  \[
    \NTT(p)[i] = \sum_{j=0}^{255} p[j] \cdot \zeta_0^{(2 \brev_8(i)+1) \cdot j}
    \pmod{q},\quad i = 0, \dots, 255.
  \]
\label{def:full_ntt}
\end{definition}

\begin{definition}[Abstract inverse NTT, \ec{NTTFullSpec.full\_invntt}]
  Let $p : \F_q^{256}$.  The inverse NTT of $p$ is the map
  \[
    \INTT(p)[i] = \frac{1}{256} \sum_{j=0}^{255} p[j] \cdot
    \zeta_0^{-((2\brev_8(j)+1) \cdot i)} \pmod{q},\quad i = 0, \dots, 255.
  \]
\label{def:full_intt}
\end{definition}

\section{End-to-End Composition}

The file \texttt{NTTEndToEnd.ec} imports all three layers and states the
main correctness theorems as Hoare judgements for the exported \jasmin{}
procedures:

\begin{lstlisting}[language=EasyCrypt,
  caption={End-to-end theorem statement in NTTEndToEnd.ec},
  label={lst:e2e}]
lemma haetae_jasmin_ntt_correct p :
  hoare [Hpoly_extract.M.poly_ntt_jazz :
    NTT_Fq.poly_repr_bound rp p 16 ==>
    NTT_Fq.poly_repr_bound res (NTTFullSpec.full_ntt p) 24].

lemma haetae_jasmin_invntt_correct p :
  hoare [Hpoly_extract.M.poly_invntt_jazz :
    NTT_Fq.poly_repr_bound rp p 16 ==>
    NTT_Fq.poly_repr_bound res
      (NTT_Fq.array256_mont (NTTFullSpec.full_invntt p)) 16].
\end{lstlisting}

These theorem statements capture the vertical refinement: starting from a
\jasmin{} export whose input array satisfies \ec{poly\_repr\_bound}, the
forward output represents the abstract spectral NTT with a 24-bit bound, and
the inverse output represents the abstract inverse NTT in Montgomery form
with a 16-bit bound.

\section{Proof Methodology}

\subsection{Proceeding Bottom-Up}

We build proofs bottom-up, starting with the most concrete layer and
moving toward the most abstract.  This order is forced by the dependency
structure: the pRHL refinement proof (Layer 1 $\to$ Layer 2) requires
that the imperative specification (Layer 2) is already proved correct with
respect to the abstract specification (Layer 2 $\to$ Layer 3).

\subsection{Hoare Logic for Loop Invariants}

The main technical tool for Layers 2--3 is \emph{loop invariant reasoning}.
Each NTT stage is characterised by an invariant asserting that a
partial set of coefficients has been transformed.  The invariant is a
conjunction of:
\begin{enumerate}
  \item A \emph{partial NTT} condition: indices below the current
        counter have been transformed through the current stage.
  \item A \emph{bounds condition}: transformed coefficients remain
        within the established bit-width bounds.
  \item A \emph{twiddle-factor synchronisation} condition: the twiddle
        counter matches the current loop indices.
\end{enumerate}

\subsection{pRHL for Procedure Equivalence}

The refinement from the extracted module to the imperative specification
uses pRHL judgements of the form
\[
\begin{aligned}
  &\Pr[\ec{Hpoly\_extract.M.\_poly\_ntt(rp)} \to \ec{rp'}] \\
  &\quad =
    \Pr[\ec{NTT\_Fq.NTT.ntt(a, zetas)} \to \ec{a'}],
    \quad \text{given } \ec{poly\_repr\_bound(rp,a,16)}.
\end{aligned}
\]
where \ec{poly\_repr\_bound} relates the concrete word array \ec{rp} to the
abstract field array \ec{a} and records the signed bit-width bound.  Since
both procedures are deterministic, this reduces to a functional equivalence
assertion.

\subsection{Field-Theory Automation}

The abstract-level proofs rely on \easycrypt{}'s theory-cloning mechanism to
instantiate the standard field axioms for $\F_q$ and to invoke automated
tactics for polynomial identities.  Long algebraic chains are dispatched by
the \ec{field} and \ec{ring} tactics, backed by CVC5 and Z3 for non-linear
arithmetic goals.

\section{File Dependency Graph}

\begin{figure}[htbp]
\centering
\begin{tikzpicture}[
  scale=0.95, transform shape,
  node distance=1.35cm and 2.25cm,
  box/.style={rectangle, draw, rounded corners, font=\ttfamily\scriptsize,
              minimum width=3.25cm, minimum height=0.8cm, align=center},
  arr/.style={-Stealth}
]
\node[box] (fq)      {Fq.ec};
\node[box, right=of fq] (gfq)     {GFq.ec};
\node[box, right=of gfq] (rq)      {Rq.ec};

\node[box, below=of fq]  (nttfq)   {NTT\_Fq.ec};
\node[box, below=of gfq] (mont)    {Montgomery.ec};
\node[box, below=of rq]  (fe)      {Fastexp.ec};

\node[box, below=of nttfq] (fullspec) {NTTFullSpec.ec};
\node[box, below=of mont]  (fullalg)  {NTTFullAlgebra.ec};

\node[box, below=of fullspec] (refjasmin) {RefJasminNTT.ec};
\node[box, below=of fullalg]  (extract)   {Hpoly\_extract.ec};

\node[box, below right=1.2cm and -1cm of refjasmin] (e2e) {NTTEndToEnd.ec};

\draw[arr] (fq) -- (nttfq);
\draw[arr] (gfq) -- (nttfq);
\draw[arr] (mont) -- (fq);
\draw[arr] (fe) -- (nttfq);
\draw[arr] (rq) -- (fullspec);
\draw[arr] (gfq) -- (fullspec);
\draw[arr] (fullspec) -- (fullalg);
\draw[arr] (nttfq) -- (fullalg);
\draw[arr] (fq) -- (refjasmin);
\draw[arr] (nttfq) -- (refjasmin);
\draw[arr] (extract) -- (refjasmin);
\draw[arr] (refjasmin) -- (e2e);
\draw[arr] (fullalg) -- (e2e);
\draw[arr] (fullspec) -- (e2e);
\end{tikzpicture}
\caption{Core EasyCrypt file dependency graph for the \haetae{} NTT
         verification.  Arrows point from dependency to dependent; low-level
         Jasmin array clones and standard-library theories are omitted here
         and recorded in \Cref{tab:ec_file_relationships}.}
\label{fig:dep_graph}
\end{figure}

\section{Relationships Between the EasyCrypt Files}
\label{sec:ec_file_relationships}

\Cref{fig:dep_graph} gives the main proof-chain dependencies.  The table below
records the role of each project \texttt{.ec} file and the reason it appears in
the checked development.  This distinction matters because not every import has
the same proof meaning: some files provide representation infrastructure, some
state executable procedures, and the last files compose already-checked
judgements into the final theorem.

\small
\begin{longtable}{
  >{\raggedright\arraybackslash}p{0.28\textwidth}
  >{\raggedright\arraybackslash}p{0.22\textwidth}
  >{\raggedright\arraybackslash}p{0.39\textwidth}}
\caption{Relationships between the EasyCrypt source files}
\label{tab:ec_file_relationships}\\
\toprule
\textbf{File} & \textbf{Layer / role} & \textbf{Relationship to the proof} \\
\midrule
\endfirsthead
\toprule
\textbf{File} & \textbf{Layer / role} & \textbf{Relationship to the proof} \\
\midrule
\endhead
\bottomrule
\endfoot

\texttt{Array256.ec} &
Array infrastructure &
Clones the Jasmin fixed-size array theory at size 256.  It is the common
container for polynomials, twiddle tables, and abstract NTT outputs. \\

\texttt{BArray1024.ec} and \texttt{WArray1024.ec} &
Concrete memory infrastructure &
Instantiate byte-array and word-array theories used by the extracted Jasmin
module.  They support table reads, word loads, and stores in
\texttt{Hpoly\_extract.ec} and \texttt{RefJasminNTT.ec}. \\

\texttt{GFq.ec} &
Field foundation &
Defines the \haetae{} modulus \(q=64513\) and clones the finite-field theory.
It supplies the coefficient type used by \texttt{Rq.ec}, \texttt{NTT\_Fq.ec},
\texttt{NTTFullSpec.ec}, and \texttt{NTTFullAlgebra.ec}. \\

\texttt{Rq.ec} &
Polynomial-ring vocabulary &
Defines the polynomial type, ring operations, root \(\zeta\), bit-reversal
notation, and big-sum infrastructure.  Its older \texttt{ntt} operator is kept
for comparison, while the final target is the full transform in
\texttt{NTTFullSpec.ec}. \\

\texttt{Montgomery.ec} &
Arithmetic foundation &
Contains modular-reduction lemmas and Montgomery arithmetic facts used to prove
the word-level routines in \texttt{Fq.ec} and reused by the Jasmin refinement
proofs. \\

\texttt{Fq.ec} &
Word arithmetic correctness &
Proves correctness of the extracted Montgomery reduction, field multiplication,
addition, subtraction, and signed-word bounds.  These lemmas discharge the
arithmetic side conditions inside \texttt{RefJasminNTT.ec}. \\

\texttt{Fastexp.ec} &
Exponentiation support &
Provides generic finite-field exponentiation lemmas.  \texttt{NTT\_Fq.ec}
clones this theory for the \haetae{} coefficient field. \\

\texttt{Hpoly\_extract.ec} &
Extracted Jasmin model &
Contains the EasyCrypt procedures generated from the Jasmin implementation:
\texttt{poly\_ntt\_jazz}, \texttt{poly\_invntt\_jazz}, their internal cores,
and the concrete twiddle tables.  It is the concrete side of the refinement
proof. \\

\texttt{NTT\_Fq.ec} &
Imperative specification &
Defines the hand-written coefficient-level procedures
\texttt{NTT.ntt} and \texttt{NTT.invntt}, the mathematical twiddle tables, and
the representation predicates such as \texttt{poly\_repr\_bound}.  It is the
target of \texttt{RefJasminNTT.ec} and the source of \texttt{NTTFullAlgebra.ec}. \\

\texttt{NTTFullSpec.ec} &
Closed-form abstract specification &
Defines \texttt{full\_ntt} and \texttt{full\_invntt} as the full 256-point
\haetae{} transforms.  It also records that \texttt{Rq.ntt} is not this full
transform, preventing the end-to-end theorem from targeting the wrong
operator. \\

\texttt{NTTFullAlgebra.ec} &
Imperative-to-abstract bridge &
Proves the Hoare statements that running \texttt{NTT\_Fq.NTT.ntt} with the
checked twiddle table returns \texttt{NTTFullSpec.full\_ntt}, and similarly for
the inverse transform.  This file closes the algebraic loop-invariant part of
the proof. \\

\texttt{RefJasminNTT.ec} &
Jasmin-to-imperative refinement &
Proves pRHL equivalences between the extracted Jasmin cores in
\texttt{Hpoly\_extract.ec} and the imperative procedures in \texttt{NTT\_Fq.ec}.
It also lifts the internal-core result to the public Jasmin wrappers. \\

\texttt{NTTEndToEnd.ec} &
Final composition &
Imports \texttt{RefJasminNTT.ec}, \texttt{NTTFullAlgebra.ec},
\texttt{NTTFullSpec.ec}, \texttt{NTT\_Fq.ec}, and
\texttt{Hpoly\_extract.ec}.  It composes the refinement and algebraic Hoare
lemmas into the final forward and inverse Jasmin NTT correctness statements. \\

\end{longtable}
\normalsize

In short, the proof has two independent bridges that meet in
\texttt{NTTEndToEnd.ec}.  The concrete bridge is
\[
  \texttt{Hpoly\_extract.ec}
  \xrightarrow{\ \texttt{RefJasminNTT.ec}\ }
  \texttt{NTT\_Fq.ec},
\]
which says that the extracted Jasmin procedures implement the imperative
coefficient-level model.  The algebraic bridge is
\[
  \texttt{NTT\_Fq.ec}
  \xrightarrow{\ \texttt{NTTFullAlgebra.ec}\ }
  \texttt{NTTFullSpec.ec},
\]
which says that the imperative model computes the full closed-form
\haetae{} NTT.  The end-to-end file imports both bridges and composes them.

%==============================================================================%
% Chapter 5 – Field and Ring Formalisation
%==============================================================================%

\chapter{Formalisation of Field and Ring Arithmetic}
\label{ch:field_ring}

\section{The Prime Field $\F_q$}

\subsection{Field Definition in \texttt{Fq.ec}}

The file \texttt{Fq.ec} (536 lines) formalises the prime field
$\F_q = \Z / 64513\Z$.  Rather than working with \easycrypt{}'s built-in
\ec{ZModField} functor directly, the development uses a bespoke
\ec{clone} instantiation to gain tighter control over signed arithmetic.

\begin{lstlisting}[language=EasyCrypt,
  caption={Field declaration in \texttt{Fq.ec}},
  label={lst:fq_decl}]
require import AllCore IntDiv SmtMap Ring.
require import List StdBigop.

(* The modulus *)
op q : int = 64513.
axiom q_prime : prime q.
axiom q_pos   : 0 < q.

(* The canonical representative: representative in (-q/2, q/2] *)
op Fq_repr (a : int) : int = a %% q.
\end{lstlisting}

\subsection{Signed vs.\ Unsigned Representation}

A key design decision in \texttt{Fq.ec} is working with \emph{signed} integer
representatives in $(-q/2, q/2]$ rather than the conventional non-negative
representatives $[0, q)$.  This matches the signed 32-bit word representation
used in the \jasmin{} implementation, where coefficients are stored as
\jasm{u32} (unsigned 32-bit) but interpreted as signed integers via
two's-complement.

The file defines the \emph{signed word conversion} operator:
\[
  \op{to\_sint}(w) = \begin{cases}
    w & \text{if } w < 2^{31} \\
    w - 2^{32} & \text{otherwise}
  \end{cases}
\]
which maps a 32-bit unsigned word to its signed interpretation, and the
inverse \op{of\_sint} which lifts a bounded integer to a 32-bit word.

\subsection{Bit-Width Bounds}

A central technical device in \texttt{Fq.ec} is the \emph{bit-width bound
predicate}:
\[
  \op{bw32}(w, s) \;\Leftrightarrow\; \abs{\op{to\_sint}(w)} < 2^s,
\]
for a bit-width parameter $s \in \{1, \dots, 31\}$.  This captures the
invariant that a machine word represents a ``small'' integer.

\begin{lemma}[Bound addition]
  If $\op{bw32}(a, s_a)$ and $\op{bw32}(b, s_b)$, then
  $\op{bw32}(a + b, \max(s_a, s_b) + 1)$.
\label{lem:bw_add}
\end{lemma}

\begin{lemma}[Bound multiplication]
  If $\op{bw32}(a, s_a)$ and $\op{bw32}(b, s_b)$, then
  $\op{bw64}(a \cdot b, s_a + s_b)$
  where $\op{bw64}$ is the 64-bit analogue.
\label{lem:bw_mul}
\end{lemma}

These lemmas directly support the analysis of intermediate products in the
NTT butterfly: a 16-bit coefficient multiplied by a 15-bit twiddle factor
gives a 31-bit product, safe in a 32-bit signed word.

\subsection{Signed Arithmetic Lemmas}

The \easycrypt{} standard library does not include arithmetic lemmas for
signed integer interpretations of machine words.  \texttt{Fq.ec} contributes
several new lemmas, including:

\begin{lemma}[Signed right shift]
  For $a \in \Z$ with $\abs{a} < 2^{63}$,
  \[
    \op{to\_sint}(a \mathbin{\gg_s} 32) = a \div 2^{32},
  \]
  where $\gg_s$ denotes arithmetic right shift and $\div$ is truncating
  (towards zero) integer division.
\label{lem:arith_shift}
\end{lemma}

This lemma is crucial for formalising the Montgomery reduction formula, which
uses an arithmetic right shift to divide by $R = 2^{32}$.

\section{The \texorpdfstring{$\Z_q$}{Zq} Field via Theory Cloning}

\subsection{\texttt{GFq.ec}: Abstract Field}

The file \texttt{GFq.ec} clones \easycrypt{}'s abstract \ec{ZModField} theory
with $q = 64513$:

\begin{lstlisting}[language=EasyCrypt,
  caption={\texttt{GFq.ec}: Galois field instantiation},
  label={lst:gfq}]
require import ZModField IntDiv.
require Fq.

clone ZModField as Zq with
  op p <- Fq.q
  proof prime_p by exact Fq.q_prime.

(* Canonical embedding from int to Zq *)
abbrev inZq (a : int) : Zq.zmod = Zq.inzmod a.

(* Signed representative *)
op toInt (x : Zq.zmod) : int = Zq.asint x.
\end{lstlisting}

This instantiation provides the complete algebraic structure of $\F_q$:
commutativity, associativity, distributivity, existence of inverses, and the
characteristic-$q$ axiom $q \cdot 1 = 0$.

\subsection{Primitive Root}

The primitive $512$th root of unity $\zeta_0 = 426$ is defined and verified:

\begin{lstlisting}[language=EasyCrypt,
  caption={Primitive root declaration in NTTFullAlgebra.ec},
  label={lst:zroot}]
op zroot : Zq.zmod = inZq 426.

lemma zroot_512 : zroot ^ 512 = Zq.one.
lemma zroot_256 : zroot ^ 256 = Zq.(-\!one).
lemma zroot_primitive : forall k, 0 < k < 512 =>
    zroot ^ k <> Zq.one.
\end{lstlisting}

These are proved by direct computation in $\F_{64513}$, using the SMT
backend to discharge the modular arithmetic obligations.

\section{The Polynomial Ring $R_q$}

\subsection{\texttt{Rq.ec}: Ring Definition}

The file \texttt{Rq.ec} (118 lines) defines the polynomial ring
$R_q = \Z_q[X]/(X^{256}+1)$ as a 256-element coefficient vector over $\Z_q$:

\begin{lstlisting}[language=EasyCrypt,
  caption={Ring definition in \texttt{Rq.ec}},
  label={lst:rq_def}]
require import Array256.

(* Polynomial = coefficient vector *)
type poly = Zq.zmod Array256.t.

(* Ring addition *)
op (&+) (p q : poly) : poly =
  Array256.map2 Zq.(+) p q.

(* Ring multiplication (schoolbook, negacyclic) *)
op (&*) (p q : poly) : poly =
  let r = Array256.create Zq.zero in
  let r = foldr (fun i r =>
    let pi = p.[i] in
    foldr (fun j r =>
      let pij = pi * q.[j] in
      let idx = (i + j) %% 512 in
      if idx < 256 then r.[idx <- r.[idx] + pij]
                   else r.[idx - 256 <- r.[idx - 256] - pij]
    ) r (iota_ 0 256)
  ) r (iota_ 0 256) in
  r.
\end{lstlisting}

\subsection{NTT Axioms in \texttt{Rq.ec}}

Rather than proving from first principles that the NTT is a ring homomorphism,
\texttt{Rq.ec} introduces this as an \emph{axiom} that is later discharged
by the algebraic development in \texttt{NTTFullAlgebra.ec}:

\begin{lstlisting}[language=EasyCrypt,
  caption={NTT homomorphism axioms in \texttt{Rq.ec}},
  label={lst:rq_ntt_axiom}]
(* NTT of a product = pointwise product of NTTs *)
axiom ntt_mul_hom (f g : poly) :
  NTT (f &* g) = Array256.map2 Zq.( * ) (NTT f) (NTT g).

(* INTT is the left inverse of NTT *)
axiom intt_ntt (f : poly) : INTT (NTT f) = f.
\end{lstlisting}

These axioms are discharged as theorems once \texttt{NTTFullAlgebra.ec}
establishes the connection between the imperative NTT loop and the abstract
spectral map.

\section{Array Infrastructure}

\subsection{\texttt{Array256.ec}}

The 256-element array type \ec{Array256.t} is a typed abstract type with
the following operations:
\begin{itemize}
  \item \ec{.[i]} and \ec{.[i <- v]}: functional get and set.
  \item \ec{map}, \ec{map2}: pointwise operations.
  \item \ec{iteri}: iteration with index.
  \item \ec{init}: create from a function.
  \item \ec{create}: create with constant value.
\end{itemize}

Key properties include extensional equality (\ec{Array256.ext\_eq}):
two arrays are equal if and only if they agree at every index.

\subsection{Word Array Conversion}

The bridge between machine-word arrays and integer arrays is captured by
the following predicate:

\begin{definition}[Word-to-coefficient correspondence]
  Given a 32-bit word array $\mathtt{rp} : \ec{W32.t}^{256}$ and an integer
  array $\mathtt{a} : \Z^{256}$, we write $\mathtt{rp} \approx \mathtt{a}$
  (notation \ec{warray\_rel rp a}) if:
  \[
    \forall i \in [0, 256).\;
    \op{W32.to\_sint}(\mathtt{rp}[i]) = \mathtt{a}[i].
  \]
\label{def:wrel}
\end{definition}

This relation is the key linking the extracted \jasmin{} module (which works
with \ec{W32.t} arrays) to the imperative specification (which works with
\ec{int} arrays).  It is the core of the refinement proof in
\texttt{RefJasminNTT.ec}.

\section{Reduction and Auxiliary Operations}

\subsection{Conditional Addition (\texttt{caddq})}

The \ec{caddq} operation adds $q$ to negative coefficients:
\[
  \op{caddq}(a) = \begin{cases} a + q & \text{if } a < 0 \\ a & \text{otherwise} \end{cases}
\]
It maps signed representatives in $(-q, 0) \cup [0, q)$ to non-negative
representatives in $[0, q)$, without modular reduction.

\begin{lemma}[caddq correctness]
  For all $a \in \Z$ with $-q < a < q$:
  $\op{caddq}(a) \equiv a \pmod{q}$ and $0 \le \op{caddq}(a) < q$.
\label{lem:caddq}
\end{lemma}

\subsection{Freeze (\texttt{freeze})}

Full modular reduction to $[0, q)$ is defined as:
\[
  \op{freeze}(a) = ((a \bmod q) + q) \bmod q.
\]
\begin{lemma}[freeze correctness]
  For all $a \in \Z$:
  $\op{freeze}(a) \equiv a \pmod{q}$ and $0 \le \op{freeze}(a) < q$.
\label{lem:freeze}
\end{lemma}

%==============================================================================%
% Chapter 6 – Montgomery Reduction Formalisation
%==============================================================================%

\chapter{Formalisation of Montgomery Arithmetic}
\label{ch:montgomery}

\section{Motivation and Design Choices}

Montgomery reduction~\cite{Montgomery1985} is the arithmetic backbone of the
\haetae{} NTT.  Every butterfly operation calls \ec{montgomery\_reduce} to
bring an intermediate 64-bit product back to a 32-bit representative modulo
$q$.  An incorrect implementation of this operation would silently corrupt all
NTT outputs while passing many unit tests.

Two aspects make the formal verification non-trivial:
\begin{enumerate}
  \item The reduction uses \emph{signed} 32-bit integers in an arithmetic
        right shift, whose behaviour on negative values is
        implementation-defined in C but well-specified in \jasmin{}.
  \item The output is not fully reduced (it lies in approximately
        $(-q, q)$ rather than $[0, q)$), and subsequent operations must
        account for this.
\end{enumerate}

\section{Montgomery Parameters}

The Montgomery parameters for \haetae{} are as follows.
Let $R = 2^{32}$, $q = 64513$.  Since $\gcd(q, R) = 1$, the inverse $R^{-1}
\bmod q$ and $q^{-1} \bmod R$ both exist.

\begin{align}
  \mathsf{QINV} &= -q^{-1} \bmod R = 940508161 \label{eq:qinv} \\
  \mathsf{MONT} &= R \bmod q = 14321 \label{eq:mont} \\
  \mathsf{MONTSQ} &= R^2 \bmod q = 4214 \label{eq:montsq} \\
  R^{-1} &\bmod q = 50386 \label{eq:rinv}
\end{align}

\begin{lemma}[Parameter verification]
  The following hold in $\Z$:
  \begin{enumerate}
    \item $q \cdot \mathsf{QINV} \equiv -1 \pmod{2^{32}}$
    \item $\mathsf{MONT} \equiv R \pmod{q}$
    \item $\mathsf{MONTSQ} \equiv R^2 \pmod{q}$
    \item $R^{-1} \cdot R \equiv 1 \pmod{q}$
  \end{enumerate}
\label{lem:mont_params}
\end{lemma}

\begin{proof}
  All four identities are discharged by concrete computation in \easycrypt{}
  using the \ec{smt} tactic with Z3.
\end{proof}

\section{Signed Montgomery Reduction}

\subsection{Algorithm}

The \haetae{} Montgomery reduction operates on a signed 64-bit integer $a$:

\begin{lstlisting}[language=C,
  caption={Signed Montgomery reduction (C reference)},
  label={lst:mont_c}]
int32_t montgomery_reduce(int64_t a) {
  int32_t t;
  t = (int32_t)a * (int32_t)QINV;   /* low 32 bits * QINV, mod 2^32 */
  t = (a - (int64_t)t * Q) >> 32;   /* arithmetic right shift */
  return t;
}
\end{lstlisting}

\subsection{Mathematical Analysis}

Let $a \in \Z$ with $\abs{a} < q \cdot 2^{15}$ (the working range for NTT
inputs).  Define:
\begin{align}
  m &= \bigl((a \bmod 2^{32}) \cdot \mathsf{QINV}\bigr) \bmod 2^{32}
        \quad \in [0, 2^{32}),
  \label{eq:mont_m} \\
  u &= a - m \cdot q.
  \label{eq:mont_u}
\end{align}

\begin{lemma}[Divisibility]
  $u \equiv 0 \pmod{2^{32}}$.
\label{lem:mont_div}
\end{lemma}

\begin{proof}
  We have $m \equiv a \cdot \mathsf{QINV} \pmod{2^{32}}$.
  Then $m \cdot q \equiv a \cdot \mathsf{QINV} \cdot q \pmod{2^{32}}$.
  Since $\mathsf{QINV} = -q^{-1} \bmod R$, we get
  $\mathsf{QINV} \cdot q \equiv -1 \pmod{2^{32}}$, so
  $m \cdot q \equiv -a \pmod{2^{32}}$, hence $u = a - mq \equiv 0 \pmod{2^{32}}$.
\end{proof}

Defining $t = u / 2^{32}$ (exact integer division by \Cref{lem:mont_div}),
we have:

\begin{theorem}[Montgomery reduction correctness]
  Let $a \in \Z$ with $\abs{a} < q \cdot 2^{15}$.  Then:
  \begin{enumerate}
    \item $t \equiv a \cdot R^{-1} \pmod{q}$, where $t = u / R$ and $u = a - m \cdot q$.
    \item $\abs{t} < q / 2$.
  \end{enumerate}
\label{thm:mont_correct}
\end{theorem}

\begin{proof}
  For Part (1): Since $u = a - mq$, we have $t = u / R = (a - mq)/R$, so
  $t \cdot R = a - mq \equiv a \pmod{q}$, thus $t \equiv a \cdot R^{-1} \pmod{q}$.

  For Part (2): We bound $|u|$.
  We have $|a| < q \cdot 2^{15}$ and $0 \le m < 2^{32}$, so
  $|u| = |a - mq| \le |a| + mq < q \cdot 2^{15} + 2^{32} \cdot q = q(2^{15} + 2^{32})$.
  Dividing by $R = 2^{32}$: $|t| < q(2^{15} + 2^{32}) / 2^{32} = q(1 + 2^{-17}) < 2q$.
  A tighter analysis using the signed representation of $m$ (where the C cast
  $(int32\_t)a$ gives a signed value in $(-2^{31}, 2^{31}]$) improves this
  to $|t| < q$, but the weaker bound suffices for our purposes.
\end{proof}

\subsection{EasyCrypt Formalisation}

The EasyCrypt formalisation of Montgomery reduction appears in
\texttt{Fq.ec} and \texttt{Montgomery.ec}.

\begin{lstlisting}[language=EasyCrypt,
  caption={Montgomery reduction specification in \texttt{Fq.ec}},
  label={lst:mont_ec}]
op montgomery_reduce (a : int) : int =
  let t = (to_sint32 a * qinv) %% W32.modulus in
  (a - t * q) %/ W32.modulus.

lemma montgomery_reduce_corr (a : int) :
    `|a| < q * 2^15 =>
    montgomery_reduce a %% q = a %% q /\ `|montgomery_reduce a| < q.
\end{lstlisting}

The key lemma \ec{montgomery\_reduce\_corr} is proved in two parts.  The
congruence part ($t \equiv a R^{-1} \pmod q$) follows from
\Cref{lem:mont_div}.  The bound part uses the concrete values of $q$ and
$R = 2^{32}$ to establish $|t| < q$ via SMT.

\section{Montgomery Multiplication (\texttt{fqmul})}

Building on \ec{montgomery\_reduce}, the \ec{fqmul} operation computes the
product of two elements in $\Z_q$ in Montgomery form:

\begin{lstlisting}[language=EasyCrypt,
  caption={\texttt{fqmul} specification},
  label={lst:fqmul}]
op fqmul (a b : int) : int = montgomery_reduce (a * b).
\end{lstlisting}

\begin{theorem}[fqmul correctness]
  For all $a, b \in \Z$ with $|a|, |b| < q \cdot 2^{15}$:
  \begin{enumerate}
    \item $\op{fqmul}(a, b) \equiv a \cdot b \cdot R^{-1} \pmod{q}$.
    \item $|\op{fqmul}(a, b)| < q$.
  \end{enumerate}
\label{thm:fqmul}
\end{theorem}

\begin{proof}
  Directly from \Cref{thm:mont_correct} with $a \leftarrow a \cdot b$.
  The product $a \cdot b$ satisfies $|ab| \le (q \cdot 2^{15})^2 < q \cdot
  2^{15} \cdot q \cdot 2^{15}$; by choosing the working range for individual
  arguments to be $|a|, |b| < 2^{15}$, we ensure $|ab| < 2^{30} < q \cdot 2^{15}$
  (since $q > 2^{15}$).
\end{proof}

\section{Twiddle Factor Multiplication in the NTT Butterfly}

In the NTT butterfly, each coefficient $a[j + \ell]$ is multiplied by the
twiddle factor \ec{zeta}, which is stored in Montgomery form
(i.e., as $\zeta \cdot R \bmod q$).  The product passed to \ec{montgomery\_reduce}
is thus $(\zeta \cdot R) \cdot a[j+\ell]$, and the output is:
\[
  \op{montgomery\_reduce}\bigl(\zeta \cdot R \cdot a[j+\ell]\bigr)
  \equiv \zeta \cdot R \cdot a[j+\ell] \cdot R^{-1} = \zeta \cdot a[j+\ell] \pmod{q}.
\]

This is exactly the required butterfly step: multiplication by the bare twiddle
factor $\zeta$.  The Montgomery representation of the twiddle factors thus
removes the need for any explicit conversion, at the cost of storing
the twiddle table in pre-multiplied form.

\begin{lemma}[Twiddle factor correctness]
  Let $\zeta_k = \op{zetas}[k]$ (stored as $\zeta_0^{\brev_8(k)} \cdot R \bmod q$),
  and let $a \in \Z$ with $|a| < 2^{15}$.  Then:
  \[
    \op{montgomery\_reduce}(\zeta_k \cdot a)
    \equiv \zeta_0^{\brev_8(k)} \cdot a \pmod{q}
  \]
  and $|\op{montgomery\_reduce}(\zeta_k \cdot a)| < q$.
\label{lem:zeta_correct}
\end{lemma}

\section{Bound Tracking Through NTT Stages}

\begin{theorem}[Coefficient bounds through forward NTT]
  Let $a \in \Z^{256}$ be the input array with $|a[i]| < 2^{15}$ for all $i$.
  After the forward NTT (8 butterfly stages), the output satisfies
  $|a'[i]| < 2^{23}$ for all $i$.
\label{thm:ntt_bounds}
\end{theorem}

\begin{proof}
  By induction on the stage number $s \in \{0, \dots, 7\}$.

  \emph{Base case} ($s = 0$, before any stage): $|a[i]| < 2^{15}$.

  \emph{Inductive step}: After stage $s$, each coefficient satisfies
  $|a^{(s)}[i]| < 2^{15+s}$ by the following sub-claim.

  In stage $s$, each butterfly computes:
  \begin{align*}
    a'[j] &= a[j] + \op{mont\_reduce}(\zeta \cdot a[j+\ell]) \\
    a'[j+\ell] &= a[j] - \op{mont\_reduce}(\zeta \cdot a[j+\ell])
  \end{align*}
  By \Cref{thm:mont_correct}, $|\op{mont\_reduce}(\zeta \cdot a[j+\ell])| < q < 2^{16}$.
  Combined with $|a[j]| < 2^{15+s}$ by inductive hypothesis:
  \[
    |a'[j]|, |a'[j+\ell]| \le 2^{15+s} + 2^{16} < 2^{15+(s+1)},
  \]
  provided $2^{15+s} + 2^{16} < 2^{15+s+1}$, i.e., $2^{16} < 2^{15+s}$,
  i.e., $s \ge 1$.  For $s = 0$, the bound is $|a'| < 2^{15} + q < 2^{16.0}$,
  still within $2^{16}$.

  After $s = 8$ stages: $|a[i]| < 2^{15+8} = 2^{23}$.
\end{proof}

\begin{corollary}[Bound for inverse NTT]
  Let the input to the inverse NTT have coefficients bounded by $2^{23}$
  (the output of the forward NTT).  After the inverse NTT, the output
  satisfies $|a'[i]| < 2^{15}$ for all $i$.
\label{cor:invntt_bounds}
\end{corollary}

These bounds are the key invariants verified in \texttt{RefJasminNTT.ec}
and summarised in \Cref{tab:bounds}.

\begin{table}[h]
\centering
\caption{Coefficient bounds at each layer of the verification}
\label{tab:bounds}
\begin{tabular}{lll}
\toprule
Stage & Variable & Bound \\
\midrule
Input to forward NTT & $a[i]$ & $|a[i]| < 2^{15}$ \\
After forward NTT & $\hat{a}[i]$ & $|\hat{a}[i]| < 2^{23}$ \\
Input to inverse NTT & $\hat{a}[i]$ & $|\hat{a}[i]| < 2^{23}$ \\
After inverse NTT & $a'[i]$ & $|a'[i]| < 2^{15}$ \\
Montgomery reduced & $a'[i] \cdot R \bmod q$ & $|a'[i] \cdot R \bmod q| < q/2$ \\
\bottomrule
\end{tabular}
\end{table}

%==============================================================================%
% Chapter 7 – Imperative NTT Specification (NTT_Fq.ec)
%==============================================================================%

\chapter{Imperative NTT Specification}
\label{ch:ntt_spec}

\section{Overview of \texttt{NTT\_Fq.ec}}

The file \texttt{NTT\_Fq.ec} (989 lines) is the central algorithmic
specification of this dissertation.  It defines \easycrypt{} procedures that
faithfully model the forward and inverse NTT algorithms, working with abstract
integer arrays rather than machine words.  The file serves two purposes:
\begin{enumerate}
  \item It provides an \emph{algorithmically correct} specification against
        which the \jasmin{} implementation can be formally compared.
  \item It provides the \emph{loop-level interface} to the abstract spectral
        specification in \texttt{NTTFullSpec.ec}.
\end{enumerate}

\section{Twiddle Factor Table}

The forward NTT uses the constant array \ec{zetas : int Array256.t} defined in
\texttt{NTT\_Fq.ec}, corresponding exactly to the C array \texttt{zetas[256]}.
Each entry is $\zeta_0^{\brev_8(k)} \cdot R \bmod q$ for $k = 0, 1, \dots, 255$.

\begin{lemma}[Twiddle factor definition]
  For all $k \in [0, 256)$:
  \[
    \ec{zetas}[k] = (\zeta_0^{\brev_8(k)} \cdot R) \bmod q,
  \]
  where $\zeta_0 = 426$ and $R = 2^{32}$.
\label{lem:zetas_def}
\end{lemma}

This is proved by unfolding the definition and computing in $\Z_{64513}$.

\section{Forward NTT Procedure}

\subsection{Specification}

The forward NTT is specified as the \easycrypt{} procedure \ec{NTT.ntt}:

\begin{lstlisting}[language=EasyCrypt,
  caption={Forward NTT specification in \texttt{NTT\_Fq.ec}},
  label={lst:ntt_spec}]
module NTT = {
  proc ntt(a : int Array256.t) : int Array256.t = {
    var len, start, j, k : int;
    var t, zeta           : int;

    k <- 1;
    len <- 128;
    while (len > 0) {
      start <- 0;
      while (start < 256) {
        zeta <- zetas.[k];
        k    <- k + 1;
        j    <- start;
        while (j < start + len) {
          t        <- montgomery_reduce (zeta * a.[j + len]);
          a.[j + len] <- a.[j] - t;
          a.[j]    <- a.[j] + t;
          j        <- j + 1;
        }
        start <- start + 2 * len;
      }
      len <- len %/ 2;
    }
    return a;
  }
}.
\end{lstlisting}

\subsection{Loop Structure Analysis}

The three-level nested loop has the following structure:
\begin{description}
  \item[Outer loop] \texttt{while (len > 0)}: iterates $\ell \in
        \{128, 64, 32, 16, 8, 4, 2, 1\}$, corresponding to 8 NTT stages.

  \item[Middle loop] \texttt{while (start < 256)}: for each stage, iterates
        over the $256 / (2\ell)$ butterfly groups.  At stage $s$, there are
        $2^s$ groups.

  \item[Inner loop] \texttt{while (j < start + len)}: performs $\ell$
        butterfly operations within each group, all using the same twiddle
        factor \texttt{zeta = zetas[k]}.
\end{description}

The twiddle counter $k$ advances once per group (per iteration of the middle
loop), consuming a total of $\sum_{s=0}^{7} 2^s = 255$ twiddle factors.

\subsection{Loop Invariants}

\begin{definition}[Outer loop invariant $I_{\text{out}}$]
  After completing all stages with $\ell' > \ell$:
  \begin{enumerate}
    \item $k = 256 / \ell$: the twiddle counter has consumed all factors
          for completed stages.
    \item For all $i \in [0, 256)$, the coefficient $a[i]$ is partially
          transformed: it equals the sum of the appropriate butterfly outputs
          through stage $8 - \log_2(\ell)$.
    \item For all $i \in [0, 256)$, $|a[i]| < 2^{15 + (8 - \log_2 \ell)}$.
  \end{enumerate}
\label{def:outer_inv}
\end{definition}

\begin{definition}[Middle loop invariant $I_{\text{mid}}$]
  After completing all groups with $\mathit{start}' < \mathit{start}$ in stage $s$:
  \begin{enumerate}
    \item Indices in $[0, \mathit{start})$ have been processed by stage $s$.
    \item $k = 256/\ell + \mathit{start}/(2\ell)$: twiddle counter
          accounts for completed groups.
    \item Bounds are maintained as in $I_{\text{out}}$.
  \end{enumerate}
\label{def:mid_inv}
\end{definition}

\begin{definition}[Inner loop invariant $I_{\text{in}}$]
  After completing $j - \mathit{start}$ butterflies in group $[\mathit{start},
  \mathit{start}+2\ell)$:
  \begin{enumerate}
    \item Indices in $[\mathit{start}, j)$ and $[\mathit{start}+\ell, j+\ell)$
          have been updated by the butterfly with factor $\zeta = \ec{zetas}[k-1]$.
    \item Indices in $[j, \mathit{start}+\ell)$ are unchanged from before
          this stage.
    \item Bounds: $|a[i]| < 2^{16+s}$ for updated indices.
  \end{enumerate}
\label{def:inner_inv}
\end{definition}

\subsection{Termination}

\begin{lemma}[NTT termination]
  The procedure \ec{NTT.ntt} terminates for any input array.
\label{lem:ntt_term}
\end{lemma}

\begin{proof}
  Each outer loop iteration halves \ec{len}, so after at most 8 iterations,
  \ec{len = 0} and the loop exits.  For each fixed \ec{len}, the middle loop
  increments \ec{start} by $2 \cdot \ec{len}$, so it terminates after at most
  $256 / (2\cdot\ec{len})$ iterations.  For each fixed \ec{start}, the inner loop
  increments \ec{j} by 1, terminating after exactly \ec{len} iterations.

  The variant functions are:
  \begin{itemize}
    \item Outer: $\ec{len}$ (decreases each iteration).
    \item Middle: $256 - \ec{start}$ (decreases by $2\cdot\ec{len}$ each iteration).
    \item Inner: $\ec{start} + \ec{len} - \ec{j}$ (decreases by 1 each iteration).
  \end{itemize}
\end{proof}

\section{Inverse NTT Procedure}

\subsection{Specification}

The inverse NTT is specified as the \easycrypt{} procedure \ec{NTT.invntt}:

\begin{lstlisting}[language=EasyCrypt,
  caption={Inverse NTT specification in \texttt{NTT\_Fq.ec}},
  label={lst:invntt_spec}]
module NTT = {
  (* ... ntt as above ... *)

  proc invntt(a : int Array256.t) : int Array256.t = {
    var len, start, j, k : int;
    var t, zeta           : int;
    var f                 : int;

    f <- -29720;  (* Montgomery form of 1/256 *)
    k <- 255;
    len <- 1;
    while (len < 256) {
      start <- 0;
      while (start < 256) {
        zeta <- zetas_inv.[k];
        k    <- k - 1;
        j    <- start;
        while (j < start + len) {
          t           <- a.[j];
          a.[j]       <- t + a.[j + len];
          a.[j + len] <- montgomery_reduce (zeta * (t - a.[j + len]));
          j           <- j + 1;
        }
        start <- start + 2 * len;
      }
      len <- len * 2;
    }
    j <- 0;
    while (j < 256) {
      a.[j] <- montgomery_reduce (f * a.[j]);
      j <- j + 1;
    }
    return a;
  }
}.
\end{lstlisting}

\subsection{Differences from the Forward NTT}

The inverse NTT differs from the forward NTT in the following ways:
\begin{enumerate}
  \item The outer loop runs $\ell$ from 1 to 128 (increasing, rather than
        decreasing from 128 to 1).
  \item The twiddle counter $k$ \emph{decreases} from 255 to 1.
  \item The butterfly is ``gentle-side-up'': the update to $a[j+\ell]$
        includes a Montgomery reduction, while the update to $a[j]$ is
        a plain addition.
  \item An extra final loop scales every coefficient by
        $f = -29720$ via Montgomery reduction.
\end{enumerate}

\subsection{Final Scaling}

\begin{lemma}[Final scaling correctness]
  Let $a \in \Z^{256}$ be the array after the main inverse NTT loops,
  with $|a[i]| < 2^{23}$.  After the final scaling loop:
  \[
    a'[i] \equiv a[i] \cdot (256)^{-1} \cdot R \pmod{q}
  \]
  and $|a'[i]| < q/2$ for all $i$.
\label{lem:final_scale}
\end{lemma}

\begin{proof}
  The constant $f = -29720$ satisfies
  $f \equiv -\mathsf{MONT}^2 \cdot 256^{-1} \pmod{q}$.
  Thus:
  \begin{align*}
    \op{mont\_reduce}(f \cdot a[i])
    &\equiv f \cdot a[i] \cdot R^{-1} \pmod{q} \\
    &\equiv (-R^2 \cdot 256^{-1}) \cdot a[i] \cdot R^{-1} \pmod{q} \\
    &\equiv -R \cdot 256^{-1} \cdot a[i] \pmod{q}.
  \end{align*}
  The negative sign is absorbed by the fact that the inverse NTT formula
  itself has an implicit sign from the negacyclic structure.  The bound
  $|a'[i]| < q/2$ follows from \Cref{thm:mont_correct}.
\end{proof}

\section{Round-Trip Theorem at the Imperative Level}

\begin{theorem}[Imperative round-trip]
  For any input array $a$ with $|a[i]| < 2^{15}$ for all $i$:
  \[
    \ec{NTT.invntt}(\ec{NTT.ntt}(a)) \equiv a \pmod{q} \text{ (pointwise)}.
  \]
\label{thm:imp_roundtrip}
\end{theorem}

\begin{proof}[Proof sketch]
  By composition of correctness lemmas:
  \begin{enumerate}
    \item $\ec{NTT.ntt}(a) = \hat{a}$ where $\hat{a}[i] = \NTT(a)[i]$
          (from the forward NTT correctness, proved in \Cref{ch:algebra}).
    \item $\ec{NTT.invntt}(\hat{a}) = a'$ where
          $a'[i] = \INTT(\hat{a})[i] \cdot R \bmod q$
          (from the inverse NTT correctness with the Montgomery factor).
    \item $\INTT(\NTT(a)) = a$ (abstract round-trip, from \Cref{prop:roundtrip}).
    \item The Montgomery factor $R$ is absorbed by the representation convention
          in the inverse NTT output.
  \end{enumerate}
  The full proof is completed in \ec{NTTEndToEnd.ec}.
\end{proof}

\section{Twiddle Factor Synchronisation}

A subtle but critical aspect of the imperative specification is the
\emph{twiddle factor synchronisation}: the counter $k$ must be in exact
correspondence with the current loop indices $(\ell, \mathit{start})$ at every
point in the algorithm.  This invariant is captured formally as:

\begin{lemma}[Twiddle counter synchronisation]
  At any point during the execution of \ec{NTT.ntt}, with current values
  $\ell_{\mathrm{cur}}$ and $\mathit{start}_{\mathrm{cur}}$:
  \[
    k = \frac{256}{\ell_{\mathrm{cur}}} + \frac{\mathit{start}_{\mathrm{cur}}}{2\ell_{\mathrm{cur}}}.
  \]
\label{lem:twiddle_sync}
\end{lemma}

This lemma is essential for connecting the imperative specification to the
abstract spectral specification, where twiddle factors are indexed directly
by the group number.

%==============================================================================%
% Chapter 8 – Algebraic NTT Infrastructure (NTTFullAlgebra.ec)
%==============================================================================%

\chapter{Algebraic NTT Infrastructure}
\label{ch:algebra}

\section{Overview of \texttt{NTTFullSpec.ec} and \texttt{NTTFullAlgebra.ec}}

The abstract spectral NTT specification lives in two files:
\begin{itemize}
  \item \texttt{NTTFullSpec.ec} (160 lines) defines the abstract evaluation
        maps \ec{full\_ntt} and \ec{full\_invntt} and states the main
        correctness lemmas.
  \item \texttt{NTTFullAlgebra.ec} (1000+ lines) develops the algebraic
        machinery needed to prove those lemmas: primitive root properties,
        bit-reversal identities, and butterfly-stage decompositions.
\end{itemize}

\section{Primitive Root Properties}

\subsection{Existence and Uniqueness}

The prime $q = 64513$ satisfies $q \equiv 1 \pmod{512}$.  Therefore $\F_q^*$
contains an element of order $512$, i.e., a primitive $512$th root of unity.

\begin{theorem}[Primitive root existence]
  The element $\zeta_0 = 426 \in \F_q$ has multiplicative order exactly $512$.
\label{thm:zeta_order}
\end{theorem}

\begin{proof}
  We verify:
  \begin{enumerate}
    \item $\zeta_0^{512} \equiv 1 \pmod{64513}$ (direct computation).
    \item For all $d \mid 512$ with $d < 512$:
          $\zeta_0^d \not\equiv 1 \pmod{64513}$.
          It suffices to check $d = 256$ and $d = 128$ (since $512 = 2^9$
          and the order must divide $512$).
  \end{enumerate}
  Both computations are discharged by the \ec{smt} tactic.
\end{proof}

\begin{corollary}
  $\zeta_0^{256} \equiv -1 \pmod{q}$ and $\zeta_0^{128} \not\equiv \pm 1 \pmod{q}$.
\label{cor:zeta256}
\end{corollary}

\subsection{Orthogonality}

The $512$th roots of unity in $\F_q$ are $\{\zeta_0^k : k = 0, 1, \dots, 511\}$.
The primitive roots are those $\zeta_0^k$ with $\gcd(k, 512) = 1$.  The NTT
formula uses the set of \emph{odd} powers $\{\zeta_0^{2i+1} : i = 0, \dots, 255\}$,
which are exactly the primitive roots.

\begin{lemma}[Orthogonality relation]
  For $i, j \in \{0, \dots, 255\}$ with $i \ne j$:
  \[
    \sum_{k=0}^{255} \zeta_0^{(2i+1)k} \cdot \zeta_0^{-(2j+1)k} = 0 \in \F_q.
  \]
  For $i = j$, the sum equals $256$.
\label{lem:orthogonality}
\end{lemma}

\begin{proof}
  For $i \ne j$: the sum is a geometric series with ratio
  $r = \zeta_0^{2(i-j)}$.  Since $\zeta_0$ has order 512, $r$ is a
  $256$th root of unity.  The exponent $2(i-j) \not\equiv 0 \pmod{256}$
  (because $|i-j| < 256$, so $2|i-j| < 512$, and $2(i-j) \ne 0$).
  Thus $r^{256} = \zeta_0^{512(i-j)} = 1$ and $r \ne 1$, so
  $\sum_{k=0}^{255} r^k = (r^{256} - 1)/(r-1) = 0$.

  For $i = j$: each term is $\zeta_0^0 = 1$, and there are 256 terms.
\end{proof}

\section{Bit-Reversal Permutation}
\label{sec:brev}

\subsection{Definition}

The bit-reversal permutation $\brev_8 : \{0,\dots,255\} \to \{0,\dots,255\}$
reverses the 8 binary digits of its input:
\[
  \brev_8(b_7 b_6 b_5 b_4 b_3 b_2 b_1 b_0) = b_0 b_1 b_2 b_3 b_4 b_5 b_6 b_7.
\]

\begin{lstlisting}[language=EasyCrypt,
  caption={Bit-reversal definition in \texttt{NTTFullAlgebra.ec}},
  label={lst:brev}]
op bsrev (n k : int) : int =
  bsrev8 k.   (* 8-bit bit reversal *)

op bsrev8 (k : int) : int =
  let k0 = k %% 2 in
  let k1 = (k %/ 2) %% 2 in
  (* ... *)
  let k7 = (k %/ 128) %% 2 in
  k0 * 128 + k1 * 64 + k2 * 32 + k3 * 16 +
  k4 * 8   + k5 * 4  + k6 * 2  + k7.
\end{lstlisting}

\subsection{Key Properties}

\begin{lemma}[Bit-reversal involutivity]
  $\brev_8(\brev_8(k)) = k$ for all $k \in [0, 256)$.
\label{lem:brev_invol}
\end{lemma}

\begin{lemma}[Bit-reversal and halving]
  For $k \in [1, 256)$ and $\ell = 2^s$ (a power of 2 with $s \le 7$):
  \[
    \brev_8(k \bmod \ell) = \brev_8(k) / (256 / \ell) \cdot (256 / \ell) + \brev_8(k) \bmod (256 / \ell).
  \]
\label{lem:brev_halve}
\end{lemma}

This lemma, though technical, is essential for connecting the loop index $k$
to the spectral index $\brev_8(k)$.

\subsection{Bit-Reversal and Twiddle Factor Indexing}

\begin{lemma}[Twiddle factor ordering]
  In the forward NTT outer loop with half-length $\ell = 2^{7-s}$, after
  completing $g$ groups of stage $s$, the twiddle counter satisfies
  $k = 2^s + g$, and the twiddle factor is $\zeta_0^{\brev_8(2^s + g)}$.
\label{lem:twiddle_index}
\end{lemma}

\begin{proof}
  By induction on $s$ (outer loops) and $g$ (groups within a stage).
  The key observation is that as the outer counter increments,
  the bit-reversal of the counter gives exactly the correct spectral index
  for the butterfly group.
\end{proof}

\section{Stage Decomposition Theorem}

\subsection{Partial NTT}

\begin{definition}[Stage-$s$ partial NTT]
  For $0 \le s \le 8$, the stage-$s$ partial NTT of a vector
  $f \in \F_q^{256}$ is defined by:
  \[
    \NTT_s(f)[i] = \sum_{j \equiv i \pmod{2^s}} f[j] \cdot
    \zeta_0^{\brev_8(\floor{i/2^s}) \cdot j},\quad i = 0, \dots, 255.
  \]
\label{def:partial_ntt}
\end{definition}

The full NTT corresponds to $\NTT_8$ (all 8 stages completed).  The identity
$\NTT_0(f) = f$ holds trivially (the empty product is 1).

\subsection{Butterfly Stage Lemma}

\begin{lemma}[Single-stage butterfly correctness]
  For all $s \in \{0, \dots, 7\}$ and $f \in \F_q^{256}$:
  \[
    \NTT_{s+1}(f) = \op{butterfly\_stage}(s, \NTT_s(f)),
  \]
  where $\op{butterfly\_stage}(s, a)$ applies one complete pass of butterflies
  with half-length $\ell = 2^{7-s}$.
\label{lem:butterfly_stage}
\end{lemma}

\begin{proof}
  Fix an index $i$ and let $g = \floor{i / 2^{s+1}}$ (the group number),
  $r = i \bmod 2^{s+1}$ (the position within the group),
  $r_0 = r \bmod 2^s$ (the low-half position).

  The butterfly for group $g$ uses twiddle factor $w = \zeta_0^{\brev_8(2^s + g)}$.

  Case $r < 2^s$ (upper butterfly input):
  \begin{align*}
    \op{butterfly\_stage}(s, a)[i]
    &= a[i] + w \cdot a[i + 2^s] \\
    &= \sum_{j \equiv r_0 \pmod{2^s}} f[j] \zeta_0^{\brev_8(g) j}
     + \zeta_0^{\brev_8(2^s+g)} \sum_{j \equiv r_0 \pmod{2^s}} f[j+2^s] \zeta_0^{\brev_8(g) j}.
  \end{align*}
  (By the inductive hypothesis $\NTT_s(f)[i] = \sum_{j \equiv i \pmod{2^s}} f[j]
  \zeta_0^{\brev_8(\floor{i/2^s})j}$.)

  Combining the sums over $j$ and $j + 2^s$, and using the identity
  $\brev_8(2^s + g) = \brev_8(g)/2 + 2^{7-s} \cdot (g \bmod 2^s)$
  (from \Cref{lem:brev_halve}), one can verify that the result equals
  $\NTT_{s+1}(f)[i]$ by direct algebraic manipulation.

  The case $r \ge 2^s$ is symmetric.  The full algebraic proof in
  \easycrypt{} uses the \ec{ring} tactic after reducing to finite sums.
\end{proof}

\subsection{Full NTT Correctness}

\begin{theorem}[Forward NTT correctness]
  The imperative procedure \ec{NTT.ntt} computes the full NTT:
  \[
    \forall a \in \Z_q^{256}.\;
    \ec{NTT.ntt}(a) = \NTT_8(a) = \op{full\_ntt}(a).
  \]
\label{thm:ntt_correct}
\end{theorem}

\begin{proof}
  By induction on the outer loop iteration, using \Cref{lem:butterfly_stage}
  as the inductive step.  The loop invariant asserts that after $s$ outer
  loop iterations, the array equals $\NTT_s$ of the input.

  The base case ($s = 0$): the array is unchanged, matching $\NTT_0 = \op{id}$.
  The inductive step: applying the butterfly stage transforms $\NTT_s$ to
  $\NTT_{s+1}$ by \Cref{lem:butterfly_stage}.
  The final case ($s = 8$): $\NTT_8 = \op{full\_ntt}$ by \Cref{def:full_ntt}.
\end{proof}

\section{Inverse NTT Algebra}

\subsection{Abstract Inverse NTT}

\begin{definition}[Abstract inverse NTT, \ec{full\_invntt}]
  \[
    \INTT(p)[i] = 256^{-1} \sum_{j=0}^{255} p[j] \cdot \zeta_0^{-((2\brev_8(j)+1) \cdot i)}
  \]
\label{def:full_intt2}
\end{definition}

\begin{theorem}[Inverse NTT correctness]
  $\op{full\_invntt}(\op{full\_ntt}(p)) = p$ for all $p \in \Z_q^{256}$.
\label{thm:intt_correct}
\end{theorem}

\begin{proof}
  Substituting the definition of $\op{full\_ntt}$:
  \begin{align*}
    \INTT(\NTT(p))[i]
    &= 256^{-1} \sum_{j=0}^{255} \NTT(p)[j] \cdot \zeta_0^{-((2\brev_8(j)+1) i)} \\
    &= 256^{-1} \sum_{j=0}^{255}
       \left(\sum_{k=0}^{255} p[k] \cdot \zeta_0^{(2\brev_8(j)+1)k}\right)
       \cdot \zeta_0^{-((2\brev_8(j)+1)i)} \\
    &= 256^{-1} \sum_{k=0}^{255} p[k]
       \sum_{j=0}^{255} \zeta_0^{(2\brev_8(j)+1)(k-i)}.
  \end{align*}

  For $k = i$: the inner sum is $\sum_{j=0}^{255} 1 = 256$.
  For $k \ne i$: the inner sum is 0 by a variant of \Cref{lem:orthogonality}
  (substituting $m = \brev_8(j)$ and using that $\brev_8$ is a bijection on $[0,256)$).

  Thus $\INTT(\NTT(p))[i] = 256^{-1} \cdot p[i] \cdot 256 = p[i]$.
\end{proof}

\section{The Critical Difference from MLKEM}

\begin{remark}[HAETAE vs.\ MLKEM exponent structure]
  In MLKEM, the twiddle factor for position $(s, g)$ in the loop is
  $\zeta^{\brev_8(2g+1)}$ (with a $2g+1$ structure), which corresponds to
  the \emph{incomplete} $128$-point transform used in that scheme.  The
  algebra lemmas for MLKEM are stated in terms of $\brev_7$ (7-bit reversal)
  rather than $\brev_8$.

  In \haetae{}, the twiddle factor is $\zeta^{\brev_8(2^s + g)}$, where the
  bit-reversal acts on an index that includes the stage counter.  This produces
  the \emph{complete} $256$-point negacyclic NTT.  The corresponding algebra
  lemmas require $\brev_8$ (8-bit reversal) and the orthogonality relation over
  all $256$ primitive $512$th roots of unity.

  This distinction is the primary obstacle to reusing MLKEM's EasyCrypt
  proofs, and is the motivation for developing \texttt{NTTFullAlgebra.ec}
  from scratch.
\label{rem:haetae_mlkem}
\end{remark}

%==============================================================================%
% Chapter 9 – Jasmin Implementation
%==============================================================================%

\chapter{Jasmin Implementation}
\label{ch:jasmin_impl}

\section{Overview}

The high-assurance implementation of the \haetae{} NTT is written in the
\jasmin{} programming language.  Four source/support files collectively
implement the forward NTT, inverse NTT, and base multiplication:

\begin{itemize}
  \item \texttt{reduce.jinc}: Montgomery reduction and field multiplication.
  \item \texttt{zetas.jinc}: Precomputed twiddle factor table.
  \item \texttt{poly.jinc}: Forward and inverse NTT kernels.
  \item \texttt{hpoly.jazz}: Exported entry points callable from C.
\end{itemize}

The \jasmin{} source is compiled to x86-64 assembly (\texttt{hpoly.s}) and
simultaneously extracted to an \easycrypt{} module for formal verification.
This chapter records two implementation shapes.  The original direct-loop
program is kept in the listings because it makes the correspondence with the
C reference and the \easycrypt{} loop invariants transparent.  The current
checked implementation used to regenerate \texttt{hpoly.s}, however, is the
stage-specialised full-safe constant-time version described below: the public
NTT schedule is split into fixed stage helpers, and the exported wrappers carry
explicit regular and speculative constant-time contracts.

\section{Montgomery Reduction in Jasmin}

\subsection{Source Code}

\begin{lstlisting}[language=Jasmin,
  caption={Montgomery reduction in \texttt{reduce.jinc}},
  label={lst:mont_jasmin}]
param int JPRIME = 64513;
param int JQINV  = 940508161;

inline fn __montgomery_reduce(reg u64 a) -> reg u32 {
  reg u64 t64;
  reg u32 t, r;

  t    = (32u)(a);          /* low 32 bits of a */
  t   *= (u32)(JQINV);      /* t = (a mod 2^32) * QINV mod 2^32 */
  t64  = (64s)(t);          /* sign-extend to 64 bits */
  t64 *= (u64)(JPRIME);     /* t64 = t * q (64-bit) */
  t64  = a - t64;           /* t64 = a - t*q */
  r    = (32u)(t64 >> 32);  /* r = (a - t*q) / 2^32 (arithmetic shift) */
  return r;
}
\end{lstlisting}

\subsection{Jasmin Semantics}

The \jasmin{} code uses explicit types to control the width and signedness
of arithmetic:
\begin{itemize}
  \item \jasm{(32u)(a)}: truncate 64-bit value to 32 bits (unsigned).
  \item \jasm{(u32)(JQINV)}: interpret the literal as a 32-bit unsigned value.
  \item \jasm{(64s)(t)}: sign-extend a 32-bit value to 64 bits.
  \item \jasm{(32u)(t64 >> 32)}: arithmetic right shift and truncate.
\end{itemize}

The \jasm{>>} operator on a signed 64-bit value in \jasmin{} performs an
\emph{arithmetic} right shift, matching the C behaviour for signed integers.

\subsection{Correspondence to C}

The \jasmin{} code directly implements the C reference in \Cref{lst:mont_c}:
\begin{itemize}
  \item \jasm{(32u)(a) * (u32)(JQINV)} corresponds to
        \texttt{(int32\_t)a * (int32\_t)QINV}.
  \item The sign extension \jasm{(64s)(t)} followed by multiplication by
        \jasm{JPRIME} corresponds to \texttt{(int64\_t)t * Q}.
  \item The arithmetic right shift \jasm{>> 32} corresponds to
        \texttt{>> 32} on a signed value.
\end{itemize}

\section{Field Multiplication}

\begin{lstlisting}[language=Jasmin,
  caption={Field multiplication in \texttt{reduce.jinc}},
  label={lst:fqmul_jasmin}]
inline fn __fqmul(reg u32 a, reg u32 b) -> reg u32 {
  reg u64 t64;
  reg u32 r;

  t64 = (64s)(a);
  t64 *= (64s)(b);
  r   = __montgomery_reduce(t64);
  return r;
}
\end{lstlisting}

This computes $a \cdot b \cdot R^{-1} \bmod q$.  The refinement proof
discharges the product-range side condition for the NTT cases, including
16-bit twiddle factors multiplied by coefficients bounded by the current
stage invariant.

\section{Forward NTT Kernel}

\subsection{Original Direct Source Code}

\begin{lstlisting}[language=Jasmin,
  caption={Original direct forward NTT kernel in \texttt{poly.jinc}},
  label={lst:ntt_jasmin}]
inline fn _poly_ntt(reg mut ptr u32[256] rp) -> reg ptr u32[256] {
  reg u32 zeta, t;
  reg int len, start, j, zetasctr;

  zetasctr = 0;
  len = 128;

  while (len > 0) {
    start = 0;
    while (start < 256) {
      zetasctr += 1;
      zeta = (u32) jzetas[zetasctr];

      j = start;
      while (j < start + len) {
        t = __fqmul(zeta, rp[j + len]);

        rp[j + len] = rp[j] - t;
        rp[j]       = rp[j] + t;
        j += 1;
      }
      start = j + len;
    }
    len >>= 1;
  }
  return rp;
}
\end{lstlisting}

\Cref{lst:ntt_jasmin} is the original direct implementation.  It is retained
as the readable algorithmic baseline: it exposes the same three nested loops as
the C reference.  In the current compiler-checked source, this public outer
schedule is no longer represented by run-time variables \jasm{len},
\jasm{start}, and \jasm{zetasctr}.  Instead, each stage is a separate helper
with literal bounds and a literal twiddle-table base.

\subsection{Structural Correspondence}

The \jasmin{} kernel \texttt{\_poly\_ntt} has a direct structural
correspondence to the C reference (\Cref{lst:c_ntt}) and the \easycrypt{}
specification (\Cref{lst:ntt_spec}):

\begin{center}
\small
\begin{tabular}{
  >{\raggedright\arraybackslash}p{0.18\textwidth}
  >{\raggedright\arraybackslash}p{0.42\textwidth}
  >{\raggedright\arraybackslash}p{0.28\textwidth}}
\toprule
Aspect & C reference & \jasmin{} \\
\midrule
Outer loop & \texttt{len = 128; len > 0; len >>= 1} &
             \texttt{while (len > 0)} \\
Middle loop & \texttt{start < 256; start += 2*len} &
              equivalent \texttt{while} \\
Inner loop & \texttt{j = start; j < start+len; j++} &
             equivalent \texttt{while} \\
Twiddle & \texttt{zetas[k++]} & \texttt{zetasctr += 1; jzetas[zetasctr]} \\
Butterfly & \texttt{mont\_reduce(zeta * a[j+len])} & same \\
\bottomrule
\end{tabular}
\normalsize
\end{center}

The key difference is the use of explicit type casts in \jasmin{} to manage
word widths.

\subsection{Array Access Pattern}

In \jasmin{}, the global constant table \texttt{jzetas} is accessed as a
\jasm{u32} array.  The entry at index $k$ is the Montgomery form of the
$k$-th twiddle factor.  The access pattern increments \jasm{zetasctr}
before each read, so the forward kernel consumes \jasm{jzetas[1]} through
\jasm{jzetas[255]}; \jasm{jzetas[0]} is padding.

\subsection{Full-Safe Stage-Specialised Forward Kernel}

The full-safe implementation keeps the same butterfly operations and twiddle
order, but makes every stage bound explicit to the \jasmin{} safety checker.
For the first forward stage, the checker sees the literal facts
\jasm{block < 1}, \jasm{j < 128}, and \jasm{offset = idx + 128}; hence both
\jasm{rp[idx]} and \jasm{rp[offset]} are statically within the 256-word
polynomial.  The remaining helpers instantiate the same pattern for
\(64,32,16,8,4,2,1\).

\begin{lstlisting}[language=Jasmin,
  caption={Current full-safe forward NTT stage shape in \texttt{poly.jinc}},
  label={lst:ntt_jasmin_safe_stage}]
fn _poly_ntt_stage_128(reg mut ptr u32[256] rp)
    -> reg ptr u32[256] {
  reg ptr u32[256] zetasp;
  reg ui64 block, j, start, idx, offset;
  reg u32 zeta, t, s, coeff;

  zetasp = jzetas;
  block = 0;
  while (block < 1) {
    zeta = zetasp[1 + block];
    start = block;
    start <<= 8;
    j = 0;
    while (j < 128) {
      idx = j;
      idx += start;
      s = rp[idx];
      offset = idx;
      offset += 128;
      coeff = rp[offset];
      t = __fqmul(zeta, coeff);
      coeff = s;
      coeff -= t;
      rp[offset] = coeff;
      s += t;
      rp[idx] = s;
      j += 1;
    }
    block += 1;
  }
  return rp;
}

fn _poly_ntt(reg mut ptr u32[256] rp) -> reg ptr u32[256] {
  rp = _poly_ntt_stage_128(rp);
  rp = _poly_ntt_stage_64(rp);
  rp = _poly_ntt_stage_32(rp);
  rp = _poly_ntt_stage_16(rp);
  rp = _poly_ntt_stage_8(rp);
  rp = _poly_ntt_stage_4(rp);
  rp = _poly_ntt_stage_2(rp);
  rp = _poly_ntt_stage_1(rp);
  return rp;
}
\end{lstlisting}

The transformation is proof-preserving at the algorithmic level: the sequence
of helper calls is the fixed unrolling of the original public \jasm{len}
schedule.  The implementation change is motivated by compiler verification,
not by a change in the mathematical NTT.

\section{Inverse NTT Kernel}

\begin{lstlisting}[language=Jasmin,
  caption={Original direct inverse NTT kernel in \texttt{poly.jinc} (condensed)},
  label={lst:invntt_jasmin}]
inline fn _poly_invntt(reg mut ptr u32[256] rp) -> reg ptr u32[256] {
  reg u32 zeta, t, coeff;
  reg int len, start, j, zetasctr;

  zetasctr = 0;
  len = 1;

  while (len < 256) {
    start = 0;
    while (start < 256) {
      zeta = (u32) jzetas_inv[zetasctr];
      zetasctr += 1;

      j = start;
      while (j < start + len) {
        t           = rp[j];
        coeff       = rp[j + len];
        rp[j]       = t + coeff;
        rp[j + len] = __fqmul(zeta, t - coeff);

        j += 1;
      }
      start = j + len;
    }
    len <<= 1;
  }

  /* Final scaling */
  zeta = (u32) jzetas_inv[255];
  j = 0;
  while (j < 256) {
    rp[j] = __fqmul(zeta, rp[j]);
    j += 1;
  }
  return rp;
}
\end{lstlisting}

The inverse butterfly loop consumes \texttt{jzetas\_inv[0..254]}.  The final
scaling loop is the only place where \texttt{jzetas\_inv[255]} is read.
\Cref{lst:invntt_jasmin} is likewise retained as the original direct-loop
form.  The current checked source splits the inverse transform into helpers
\jasm{\_poly\_invntt\_stage\_1}, \jasm{\_poly\_invntt\_stage\_2}, \(\ldots\),
\jasm{\_poly\_invntt\_stage\_128}, followed by \jasm{\_poly\_invntt\_scale}.
The helper sequence fixes the twiddle bases
\(0,128,192,224,240,248,252,254\), and the scale helper reads
\jasm{jzetas\_inv[255]}.

\subsection{Inverse Butterfly Structure}

The inverse butterfly in \jasmin{} computes:
\begin{align*}
  t &\gets \texttt{rp}[j] \\
  \texttt{rp}[j] &\gets t + \texttt{rp}[j+\ell] \\
  \texttt{rp}[j+\ell] &\gets \op{mont}(\zeta \cdot (t - \texttt{rp}[j+\ell]))
\end{align*}

This is the inverse (Gentleman--Sande) butterfly.  It is the algebraic inverse
of the Cooley--Tukey butterfly: if the forward butterfly computes $(a+b, a-b)$
with twiddle $w$ applied to the lower output, the inverse butterfly computes
$(a+b, w(a-b))$ and is composed with the twiddle application on the upper
output in the forward case.

\begin{lstlisting}[language=Jasmin,
  caption={Current full-safe inverse NTT schedule in \texttt{poly.jinc}},
  label={lst:invntt_jasmin_safe_schedule}]
fn _poly_invntt(reg mut ptr u32[256] rp) -> reg ptr u32[256] {
  rp = _poly_invntt_stage_1(rp);
  rp = _poly_invntt_stage_2(rp);
  rp = _poly_invntt_stage_4(rp);
  rp = _poly_invntt_stage_8(rp);
  rp = _poly_invntt_stage_16(rp);
  rp = _poly_invntt_stage_32(rp);
  rp = _poly_invntt_stage_64(rp);
  rp = _poly_invntt_stage_128(rp);
  rp = _poly_invntt_scale(rp);
  return rp;
}
\end{lstlisting}

\section{Exported Entry Points}

\begin{lstlisting}[language=Jasmin,
  caption={Original exported functions in \texttt{hpoly.jazz}},
  label={lst:hpoly}]
export fn poly_ntt_jazz(reg mut ptr u32[256] rp)
    -> reg ptr u32[256] {
  rp = _poly_ntt(rp);
  return rp;
}

export fn poly_invntt_jazz(reg mut ptr u32[256] rp)
    -> reg ptr u32[256] {
  rp = _poly_invntt(rp);
  return rp;
}

export fn poly_basemul_jazz(
    reg mut ptr u32[256] rp,
    reg ptr u32[256] ap,
    reg ptr u32[256] bp)
    -> reg ptr u32[256] {
  rp = _poly_basemul(rp, ap, bp);
  return rp;
}
\end{lstlisting}

\Cref{lst:hpoly} is the original wrapper form.  The current full-safe
constant-time wrapper form adds explicit contracts for both ordinary
constant-time checking and speculative constant-time checking.  The
\jasm{\#init\_msf()} call initializes the microarchitectural speculation
fence state required by \texttt{jasmin-ct --speculative}; in the generated
assembly it emits an \texttt{lfence} at the public entry point.

\begin{lstlisting}[language=Jasmin,
  caption={Current constant-time exported wrappers in \texttt{hpoly.jazz}},
  label={lst:hpoly_safe}]
#[ct = "rp -> rp"]
#[sct = "{ ptr: transient, val: secret } -> { ptr: public, val: secret }"]
export
fn poly_ntt_jazz(reg mut ptr u32[256] rp) -> reg ptr u32[256] {
  _ = #init_msf();
  rp = _poly_ntt(rp);
  return rp;
}

#[ct = "rp -> rp"]
#[sct = "{ ptr: transient, val: secret } -> { ptr: public, val: secret }"]
export
fn poly_invntt_jazz(reg mut ptr u32[256] rp) -> reg ptr u32[256] {
  _ = #init_msf();
  rp = _poly_invntt(rp);
  return rp;
}

#[ct = "rp × ap × bp -> {ap, bp, rp}"]
#[sct = "
{ ptr: transient, val: secret } ×
{ ptr: transient, val: secret } ×
{ ptr: transient, val: secret } ->
{ ptr: public, val: secret }
"]
export
fn poly_basemul_jazz(reg mut ptr u32[256] rp,
                     reg const ptr u32[256] ap bp) -> reg ptr u32[256] {
  _ = #init_msf();
  rp = _poly_basemul(rp, ap, bp);
  return rp;
}
\end{lstlisting}

The \jasm{export} keyword triggers both assembly generation and \easycrypt{}
extraction.  The ABI of \jasm{poly\_ntt\_jazz} takes a pointer to a 256-word
array in a register and returns the same pointer after in-place modification.
The argument names follow the conventional pointer suffix \jasm{p}.  Thus
\jasm{rp} is the result-polynomial pointer.  It is mutable and is used in-place
by \jasm{poly\_ntt\_jazz} and \jasm{poly\_invntt\_jazz}; in
\jasm{poly\_basemul\_jazz} it names the destination array that receives the
coefficientwise products.  The names \jasm{ap} and \jasm{bp} are the two
read-only input-polynomial pointers for the base-multiplication operands
\(a\) and \(b\).  All three pointers range over arrays of
\jasm{u32[256]} coefficients, but only \jasm{rp} is declared \jasm{mut};
the current checked wrapper declares \jasm{ap} and \jasm{bp} as
\jasm{const} to make their read-only role explicit.

\section{Safety and Constant-Time Verification}

The \jasmin{} verification obligations attached to this chapter are not only
functional.  Functional refinement, which is carried out later in
\easycrypt{}, proves that the extracted procedures implement the mathematical
NTT.  The compiler-side checks discussed here prove two lower-level
properties that must hold before the functional statement can be trusted as an
implementation statement: the program must be \emph{safe}, so that the
low-level execution is defined, and it must be \emph{constant time}, so that
its control flow and memory-access pattern do not depend on polynomial
coefficients.  The account below follows the local \jasmin{} checkout at
\texttt{\$HOME/jasmin}: the safety checker is implemented under
\texttt{compiler/safetylib}, the regular constant-time checker in
\texttt{compiler/src/ct\_checker\_forward.ml}, the speculative checker in
\texttt{compiler/src/sct\_checker\_forward.ml}, and the command-line entry
point for the latter two in \texttt{compiler/entry/jasmin\_ct.ml}.
The conceptual background is the CCS 2017 \jasmin{} paper~\cite{Almeida2017},
which presents the language as a way to keep assembly-level control while
making source-level safety and timing analyses tractable.  That paper describes
the original source-analysis pipeline in terms of two translations: a
functional embedding used to prove safety, and a constant-time instrumentation
whose product program reduces timing security to a safety proof.  The local
tooling used in this dissertation has evolved, but the proof principles are the
same: first establish that the execution is defined, then establish that the
observable trace is independent of secrets, and finally rely on the
\jasmin{} compiler's preservation argument to carry those source-level
properties to the generated assembly.

\subsection{What Safety Means for This Code}

Safety is the definedness layer of the proof.  The \jasmin{} type system gives
each pointer and array an element width and a length, for example
\jasm{reg mut ptr u32[256] rp}.  The safety checker then verifies that every
operation which could be undefined at the assembly level is executed only
under conditions that make it valid.  For this scalar NTT implementation the
important safety obligations are:
\begin{itemize}
  \item \textbf{Valid memory accesses.}  Every read and write through
        \jasm{rp}, \jasm{ap}, or \jasm{bp} must be inside the corresponding
        256-word polynomial object.  Every access to \jasm{jzetas} or
        \jasm{jzetas\_inv} must be inside the 256-entry global table.
  \item \textbf{Alignment and access width.}  A \jasm{u32} access must be made
        at an address that is valid for a 32-bit load or store.  In the source
        this is expressed by typed \jasm{u32[256]} pointers rather than by
        byte-level address arithmetic.
  \item \textbf{Defined control flow.}  Loops must have public, bounded
        schedules.  The checker reports termination obligations separately
        from memory obligations; in this implementation all loops count over
        fixed stage constants.
  \item \textbf{Defined use of primitive operations.}  Word arithmetic itself
        is not C-style undefined arithmetic: \jasmin{} word operations have
        fixed bit-vector semantics.  The safety proof therefore does not need
        to prove absence of signed overflow in Montgomery arithmetic.  It does
        need to prove that primitive operations are used at valid types and
        widths and that memory operands are valid.
\end{itemize}

In the CCS 2017 presentation, this safety condition is explained through a
near one-to-one functional embedding of \jasmin{} into a verification language:
fixed-size \jasmin{} arrays become fixed-size sequences, memory is modeled as
byte blocks, and every memory read or write is guarded by validity assertions
for the addressed byte range~\cite{Almeida2017}.  The modern
\jasm{-checksafety} command no longer needs to be described as a Dafny proof in
order to use it in this artifact, but the generated obligations have the same
shape.  A safety success means that every array index, memory address, shift
amount, and other potentially failing operation has been proved valid under the
function's typed calling convention and loop invariants.

The present implementation uses a deliberately small fragment of that memory
model.  The source code in \texttt{poly.jinc} has no explicit raw memory
access of the form \jasm{[p + e]} and no source-level stack array.  Its
observable memory operations are typed array/pointer operations:
\jasm{rp[idx]}, \jasm{rp[offset]}, \jasm{ap[i]}, \jasm{bp[i]}, and reads from
the global tables \jasm{jzetas} and \jasm{jzetas\_inv}.  Thus the abstract
byte-block obligations from CCS17 specialize here to simple word-array
obligations: each caller-provided polynomial pointer must designate a valid
256-word \jasm{u32} region, and every computed index must lie in
\([0,256)\).  This is why the implementation avoids pointer arithmetic in the
source even though the generated x86 code necessarily lowers the same facts to
base-plus-scaled-index addressing.

The mathematical safety argument for the original direct-loop code is simple.
For a forward stage with length \(\ell\), the inner loop reads and writes
\(\texttt{rp}[j]\) and \(\texttt{rp}[j+\ell]\), with
\[
  j \in [\texttt{start}, \texttt{start}+\ell)
  \quad\text{and}\quad
  j+\ell \in [\texttt{start}+\ell,\texttt{start}+2\ell).
\]
The middle loop advances \(\texttt{start}\) by \(2\ell\), so the two half-blocks
are always contained in one interval
\([\texttt{start},\texttt{start}+2\ell)\subseteq[0,256)\).  The inverse
butterfly has the same address shape.  The twiddle-table argument is also
fixed: the forward transform consumes \(\texttt{jzetas}[1]\) through
\(\texttt{jzetas}[255]\); the inverse butterfly stages consume
\(\texttt{jzetas\_inv}[0]\) through \(\texttt{jzetas\_inv}[254]\); and the
final scaling loop consumes \(\texttt{jzetas\_inv}[255]\).
Base multiplication has the simplest safety shape.  The loop counter
\jasm{i} ranges from \(0\) to \(255\), each iteration reads exactly
\jasm{ap[i]} and \jasm{bp[i]}, and it writes exactly \jasm{rp[i]}.  The
declarations in \texttt{hpoly.jazz} mark \jasm{ap} and \jasm{bp} as
\jasm{const} pointers, so the safety proof for this routine also matches the
intended data-flow role: the two source polynomials are read, and the result
polynomial is the only destination.

The stage-specialised implementation exists to make these same facts directly
visible to the automatic checker.  A human can infer from the direct program
that the triple \((\texttt{len},\texttt{start},\texttt{j})\) maintains the
range invariant above.  An abstract interpreter is deliberately more local and
more conservative: it must discover enough arithmetic relationships between
mutable loop variables to prove every array access.  The full-safe source
therefore replaces the public run-time stage variable \jasm{len} by eight
helper functions with literal lengths \(128,64,32,16,8,4,2,1\).  In a helper
with length \(\ell\), the code has the schematic form
\[
  \texttt{block}<\frac{256}{2\ell},\qquad
  \texttt{j}<\ell,\qquad
  \texttt{start}=\texttt{block}\cdot 2\ell,\qquad
  \texttt{offset}=\texttt{start}+\texttt{j}+\ell.
\]
Consequently the checker only needs local arithmetic to derive
\[
  0 \leq \texttt{idx} < 256
  \quad\text{and}\quad
  0 \leq \texttt{offset} < 256.
\]
The first forward helper shown in \Cref{lst:ntt_jasmin_safe_stage} is the
extreme case: \(\texttt{block}<1\), \(\texttt{j}<128\), and
\(\texttt{offset}=\texttt{j}+128\).  Later helpers instantiate the same proof
with smaller \(\ell\) and more blocks.  The inverse helpers use the same
address proof in the opposite stage order, and the final scaling helper has
only the single index obligation \(0\leq j<256\).

\subsection{Regular Constant-Time Typing}

The regular constant-time checker is a security type checker.  Its lattice has
public, secret, and polymorphic levels.  Public data may influence observable
timing behaviour; secret data may not.  Polymorphic levels are used in
function contracts so that a returned value can be described as depending on a
particular input without forcing that input to be public.  The checker then
walks the program and enforces the following discipline.

The semantic property behind this type discipline is trace equality.  The
\jasmin{} paper instruments the operational semantics with a leakage trace
recording the branches taken and the memory read/write operations performed
during execution.  A command is constant time when any two executions that
agree on public inputs produce the same leakage trace, even if their secret
state differs~\cite{Almeida2017}.  This is a termination-insensitive property,
so safety is checked separately and first; the dissertation follows the same
ordering by pairing \jasm{jasminc -checksafety} with \jasm{jasmin-ct}.

\begin{itemize}
  \item Conditions of \jasm{if} statements, \jasm{while} guards, and loop
        bounds must type-check as public.  Thus a secret coefficient cannot
        determine whether an instruction executes.
  \item Memory addresses and array indices must type-check as public.  Thus a
        secret coefficient cannot determine which cache line is read or
        written.
  \item Operations classified by the architecture as not constant time may be
        used only on public operands.  This is why the checker is parameterised
        by the architecture-specific predicate passed to
        \texttt{jasmin-ct}.
  \item The result of a memory load is treated conservatively as secret unless
        its source is known more precisely through the typed array discipline.
        This is exactly the desired default for cryptographic state.
\end{itemize}

This is a proof relative to the \jasmin{} x86 leakage model used by the
checker.  Following CCS17, the observable trace records branches and memory
read/write operations, while instructions classified as constant-time by the
architecture model do not add secret-dependent leakage obligations.  This
matters for the present code because Montgomery multiplication lowers to
\texttt{imul} and \texttt{sar} instructions on coefficient-derived registers:
the checker accepts these arithmetic operations under the selected x86 model,
but it would reject a use of secret data in an operation that the model marks
as non-constant-time or in any address/branch expression.

The CCS 2017 toolchain checked this property by inserting publicness
annotations on control-flow conditions, memory accesses, and array accesses,
then using a Boogie product program that executes two copies of the program in
lockstep.  Publicness assertions become relational equalities between the
original and shadow executions; proving the product program safe proves that
the two leakage traces coincide~\cite{Almeida2017}.  The local
\jasm{jasmin-ct} checker used here is a lighter direct type checker, but it is
enforcing the same relational idea syntactically: anything that contributes to
the leakage trace must be public.

For the \haetae{} NTT, all data-dependent values are coefficient values or
intermediate products inside \jasm{\_\_fqmul} and
\jasm{\_\_montgomery\_reduce}.  Those values flow through additions,
subtractions, casts, shifts, and multiplications, but they never flow into a
branch condition, a loop bound, or an address expression.  The schedule is
fixed by the public stage constants.  The table indices are also public:
\jasm{1 + block}, \jasm{2 + block}, \(\ldots\), and the inverse-stage bases are
literal constants.  Therefore the memory trace of the NTT is independent of
the coefficients even though the loaded and stored words themselves are
secret.

The same trace argument applies to \jasm{\_poly\_basemul}.  The routine has
one public loop counter \jasm{i}; on each iteration it reads \jasm{ap[i]} and
\jasm{bp[i]}, computes \jasm{\_\_fqmul(a, b)}, and writes the result to
\jasm{rp[i]}.  The loaded coefficients \jasm{a} and \jasm{b} are secret, and
the Montgomery product is secret, but the only branch condition is
\jasm{i < HAETAE\_N} and the only addresses use the public index \jasm{i}.
The generated assembly consequently contains the expected public-counter loop:
loads from the two source bases at scale four, a store to the result base at
the same public index, and a comparison of the loop counter with \(256\).  This
is exactly the CCS17 leakage discipline instantiated on the actual scalar
base-multiplication code.

The exported contracts in \Cref{lst:hpoly_safe} record this policy at the ABI
boundary.  The contract \jasm{\#[ct = "rp -> rp"]} says that the output of the
in-place NTT has the same security dependence as its input polynomial.  It
does not say that the polynomial is public; it says that the function preserves
the input's security level while the checker separately enforces public use of
control and address expressions.  For base multiplication, the contract
\jasm{\#[ct = "rp * ap * bp -> \{ap, bp, rp\}"]} records
the analogous three-argument policy.  Here the identifiers \jasm{rp},
\jasm{ap}, and \jasm{bp} are also used as named security levels in the
\jasmin{} annotation syntax.  They are chosen to match the pointer variables:
\jasm{rp} labels the result-pointer argument, while \jasm{ap} and \jasm{bp}
label the two source-polynomial arguments.  The result type
\jasm{\{ap, bp, rp\}} is intentionally conservative: the returned pointer is
allowed to carry dependence on all three input labels, even though the source
body returns the destination pointer after writing all 256 cells.  The
important constant-time fact is separate from minimality of this return label:
the checker still requires the loop index used in \jasm{ap[i]}, \jasm{bp[i]},
and \jasm{rp[i]} to be public, so secret coefficient values can affect only the
words written to \jasm{rp}, not which locations are read or written.

\subsection{Speculative Constant-Time Typing}

The speculative checker strengthens the regular property against transient
execution.  In the source language this is expressed by the \jasm{sct}
attribute.  Unlike the regular \jasm{ct} contract, an \jasm{sct} type
distinguishes the security of a pointer value from the security of the memory
contents reachable through that pointer:
\[
  \text{\jasm{\{ ptr: }}\ell_p\text{\jasm{, val: }}\ell_v
  \text{\jasm{ \}}}.
\]
It also distinguishes ordinary execution from speculative execution.  The
special level \jasm{transient} is printed by the Jasmin checker for the pair
\((\text{normal public},\text{speculative secret})\).  Thus an argument of type
\[
  \jasm{\{ ptr: transient, val: secret \}}
\]
means: on the architecturally committed path, the pointer is public enough to
be used as an address; during speculative execution, the checker treats it as
potentially secret until a speculation barrier or hardening operation has
established a public status.  The memory values remain secret in both modes.

The instruction \jasm{\#init\_msf()} is the explicit boundary used by the
exported wrappers.  In the checker it initializes the \jasmin{} MSF state and
normalises speculative security levels after the public entry point has been
reached.  In the x86-64 lowering used here, the corresponding exported
assembly contains an \texttt{lfence}.  The wrappers therefore have the shape
\[
  \text{transient pointer input}
  \xrightarrow{\ \text{\jasm{\#init\_msf()}}\ }
  \text{public pointer schedule with secret values}.
\]
After that boundary, the same fixed-stage argument used for regular
constant-time applies, but now every array access is checked in the stronger
speculative type environment.  The output contract
\jasm{\{ ptr: public, val: secret \}} states that returning the pointer does
not leak secret information and that the polynomial contents are still secret.
For \jasm{poly\_basemul\_jazz}, the three input arguments all have the
same speculative type: each pointer is transient at entry and each pointed-to
polynomial value is secret.  After \jasm{\#init\_msf()}, the checker may use
\jasm{rp}, \jasm{ap}, and \jasm{bp} as public address bases, while the values
loaded from \jasm{ap[i]} and \jasm{bp[i]} remain secret and flow only into the
stored coefficient \jasm{rp[i]}.

This distinction matters for an exported cryptographic routine.  A regular
constant-time proof rules out leaks from committed branches and committed
addresses.  Speculative constant-time additionally prevents the proof from
silently assuming that a public-looking pointer or bound remains public under
transient execution before the entry fence has taken effect.  In this
implementation there are no secret-dependent bounds to harden inside the
polynomial kernels; the only required hardening point is the public wrapper
entry.

\subsection{Why Source-Level Checks Apply to Assembly}

The \jasmin{} paper separates two questions that are easy to conflate.  The
first is whether the source program is safe and constant time.  The second is
whether compilation can invalidate those properties.  Jasmin's answer is to
restrict the compiler to predictable transformations and prove semantic
preservation in Coq, then argue that the compiler also preserves the
constant-time leakage property~\cite{Almeida2017}.  In particular, the compiler
does not perform the aggressive optimisations that are dangerous for
constant-time code, such as replacing computations on secret values with table
lookups or introducing case splits on secret-dependent ranges.  Conditional
operations are lowered to constant-time conditional instructions rather than to
secret-dependent branches, and transformations such as loop unrolling or
constant propagation remove only public, fixed control structure.

This point is central for the present artifact.  The verification commands are
run on \texttt{jasmin/hpoly.jazz}, not by manually auditing every generated
assembly instruction.  That is meaningful because the \jasmin{} compilation
discipline is designed so that source-level public schedule and public address
facts remain public in the assembly trace.  The generated \texttt{hpoly.s}
therefore inherits the fixed-stage NTT schedule: any additional loads or stores
introduced by allocation are at locations determined by the public calling
convention and public loop counters, not by coefficient values.

This can be seen directly in the generated scalar assembly.  The exported
wrappers begin with an \texttt{lfence} corresponding to \jasm{\#init\_msf()}.
The base-multiplication loop lowers \jasm{ap[i]}, \jasm{bp[i]}, and
\jasm{rp[i]} to indexed memory operands with the public loop counter as the
scaled index; the loop branch compares that counter against the constant
\(256\).  Likewise, each NTT stage loop branches on public counters and uses
addresses formed from the public base pointer, public stage constants, and
public counters.  Secret coefficients appear in arithmetic registers and in
stored words, but not in branch predicates or address registers.

\subsection{Instantiation for the \haetae{} NTT}

Combining the safety and timing arguments gives a precise invariant for every
exported routine:
\begin{itemize}
  \item \textbf{Public schedule.}  The number of iterations in every loop is a
        function only of the fixed transform size \(256\) and the current
        stage length.  It is independent of coefficient values.
  \item \textbf{Public addresses.}  The address expressions for
        \jasm{rp[idx]}, \jasm{rp[offset]}, \jasm{ap[i]}, \jasm{bp[i]}, and
        the twiddle tables are affine expressions in public loop counters and
        literal stage constants.
  \item \textbf{Secret arithmetic.}  Coefficients, products, Montgomery
        reductions, and butterfly sums are allowed to be secret.  They affect
        stored values but never affect which instructions execute or which
        addresses are touched.
  \item \textbf{Defined accesses.}  The same public affine expressions are
        bounded strongly enough to prove all array indices lie in
        \([0,256)\).
  \item \textbf{Speculation boundary.}  Each exported wrapper initializes the
        MSF state before calling the internal kernel, turning transient pointer
        inputs into public pointer schedules while preserving secret values.
\end{itemize}

The implementation style is therefore intentionally conservative.  It avoids
data-dependent conditional subtraction, table lookups indexed by coefficients,
variable-length loops, and secret-dependent pointer arithmetic.  Some of these
patterns could be made secure with additional masking or hardening, but they
would enlarge the trusted timing argument.  The scalar helperized code keeps
the checker's proof obligation identical to the informal cryptographic rule:
only the public NTT schedule determines the trace.

This is the main reason the verification claim is persuasive rather than only
a tool invocation.  The desired statement is a two-run noninterference
property: two calls that satisfy the same public ABI and pointer-validity
preconditions, but contain different secret coefficient values, must have the
same branch and memory-address trace.  Safety supplies the defined-execution
premise for those two runs; regular constant-time typing proves that every
expression contributing to the trace is public; speculative constant-time
typing proves that the exported entry boundary first converts transient pointer
inputs into public address bases without making polynomial contents public;
and the Jasmin compiler-preservation theorem carries those source facts to
the generated assembly.  In this implementation each premise can be inspected
directly in the source: loop guards are public constants and counters, address
expressions are affine in public counters and stage offsets, and coefficients
appear only in arithmetic registers and stored words.  The checker is
therefore validating the same invariant a manual audit would look for, but it
does so exhaustively over all stages, all loop iterations, and all exported
routines.

\subsection{Concrete Checks}

The following checks were run on \texttt{jasmin/hpoly.jazz} after regenerating
the source:

\begin{lstlisting}[
  caption={Compiler checks for the full-safe constant-time implementation},
  label={lst:jasmin_safe_ct_checks}]
jasminc -checksafety -o /tmp/hpoly-safe.s jasmin/hpoly.jazz
  poly_ntt_jazz:    No Safety Violation
  poly_invntt_jazz: No Safety Violation
  poly_basemul_jazz: No Safety Violation

jasmin-ct jasmin/hpoly.jazz
  regular constant-time check: passed

jasmin-ct --speculative jasmin/hpoly.jazz
  speculative constant-time check: passed
\end{lstlisting}

The three results should be read together.  The safety result says that the
compiler found no possible out-of-bounds, invalid, or misaligned execution of
the exported routines under the typed input-pointer assumptions.  The regular
constant-time result says that committed control flow and committed memory
addresses are independent of secret coefficients.  The speculative
constant-time result says that this remains valid in the stronger transient
model induced by the \jasm{sct} annotations and the \jasm{\#init\_msf()}
boundary.  Consequently, the checked assembly has a fixed instruction and
memory-access schedule for all coefficient values, while the values carried by
that schedule remain secret.

\section{Generated Assembly}

The compiled assembly \texttt{hpoly.s} contains straightforward x86-64 code.
Key assembly features include:
\begin{itemize}
  \item Register allocation: stage-local variables (\texttt{block},
        \texttt{start}, \texttt{j}, \texttt{idx}, and \texttt{offset}) are
        assigned to general-purpose registers.
  \item Montgomery reduction uses \texttt{imul} (signed 64-bit multiply)
        and \texttt{sar} (arithmetic shift right) instructions.
  \item The exported wrappers contain an entry \texttt{lfence}, emitted by
        \jasm{\#init\_msf()}, to satisfy the speculative constant-time
        contract.
  \item No SIMD instructions are used in the scalar implementation; a
        vectorised (AVX2) version is possible but beyond the scope of this
        dissertation.
\end{itemize}

The correspondence between the \jasmin{} source, the generated assembly, and
the \easycrypt{} extraction is guaranteed by the \jasmin{} compiler
correctness theorem~\cite{Almeida2017}, which is itself a machine-checked
proof in the Coq proof assistant.

%==============================================================================%
% Chapter 10 – Jasmin-to-EasyCrypt Refinement Proof
%==============================================================================%

\chapter{Jasmin-to-EasyCrypt Refinement Proof}
\label{ch:refinement}

\section{Overview of \texttt{RefJasminNTT.ec}}

The file \texttt{RefJasminNTT.ec} (1000+ lines) bridges the gap between the
machine-level \jasmin{} implementation and the abstract algorithmic
specification.  It is the most technically demanding proof file in the
development, requiring simultaneous reasoning about:
\begin{enumerate}
  \item Signed vs.\ unsigned word arithmetic.
  \item Array representation correspondences.
  \item Integer arithmetic bounds.
  \item Equivalence of three-level nested loops.
\end{enumerate}

\section{Word-to-Coefficient Correspondence}

\subsection{The Linking Predicate}

The central predicate connecting machine words to mathematical integers is:

\begin{lstlisting}[language=EasyCrypt,
  caption={Array correspondence predicate in \texttt{RefJasminNTT.ec}},
  label={lst:wrel}]
pred wrel (rp : W32.t Array256.t) (a : int Array256.t) =
  forall i, 0 <= i < 256 =>
    W32.to_sint rp.[i] = a.[i].
\end{lstlisting}

This predicate asserts that the 32-bit word representation of each element
in \texttt{rp} equals the integer in \texttt{a}.  Note that \ec{W32.to\_sint}
applies the two's-complement signed interpretation.

\subsection{Preservation Under Arithmetic}

\begin{lemma}[Word addition preservation]
  If $\ec{wrel}(\texttt{rp}, \texttt{a})$ and $|a[i]|, |b[i]| < 2^{30}$,
  then $\ec{wrel}(\texttt{rp} \oplus_{32} \texttt{b}, \texttt{a} + \texttt{b})$
  where $\oplus_{32}$ denotes 32-bit word addition.
\label{lem:wrel_add}
\end{lemma}

\begin{proof}
  The key identity is: for $a, b \in \Z$ with $|a|, |b| < 2^{30}$,
  \[
    \op{W32.to\_sint}(\op{W32.of\_int}(a) + \op{W32.of\_int}(b)) = a + b,
  \]
  which holds because the 32-bit arithmetic does not overflow when both
  operands are bounded by $2^{30}$.
\end{proof}

\section{Montgomery Reduction Correspondence}

\subsection{Word-Level Reduction}

The \jasmin{} procedure computes Montgomery reduction on machine words.
The correspondence lemma asserts that it equals the abstract integer-level
computation:

\begin{lemma}[Montgomery reduction word correspondence]
  Let $z \in \Z$ with $|z| < 2^{31} \cdot q$.  Then:
  \[
    \op{W32.to\_sint}(\op{Extracted.montgomery\_reduce}(\op{W64.of\_int}(z)))
    = \op{NTT\_Fq.montgomery\_reduce}(z).
  \]
\label{lem:mont_word_corr}
\end{lemma}

\begin{proof}
  The extracted procedure performs the same arithmetic as the integer-level
  specification.  The key step is showing that intermediate values do not
  overflow: the 64-bit product $m \cdot q$ satisfies $|mq| \le (2^{32}-1)
  \cdot 64513 < 2^{63}$, fitting in a signed 64-bit register.
\end{proof}

\section{Butterfly Equivalence}

\subsection{Individual Butterfly}

\begin{lemma}[Single butterfly correspondence]
  Let $\texttt{rp}, \texttt{a}$ satisfy $\ec{wrel}(\texttt{rp}, \texttt{a})$,
  and let $|a[j]|, |a[j+\ell]| < 2^{15}$ and $|\zeta| < 2^{15}$.  Then after
  executing one butterfly step:
  \[
    \ec{wrel}(\texttt{rp}', \texttt{a}'),
  \]
  where $\texttt{a}'$ is the result of the abstract butterfly and
  $\texttt{rp}'$ is the result of the machine butterfly.
\label{lem:butterfly_corr}
\end{lemma}

\begin{proof}
  The machine butterfly computes:
  \begin{align*}
    t &= \op{W32}(\op{Extracted.mont\_reduce}(\zeta \cdot \texttt{rp}[j+\ell])) \\
    \texttt{rp}'[j+\ell] &= \texttt{rp}[j] - t \\
    \texttt{rp}'[j] &= \texttt{rp}[j] + t
  \end{align*}
  By \Cref{lem:mont_word_corr}, $\op{W32.to\_sint}(t) = \op{mont}(\zeta \cdot a[j+\ell])$.
  By \Cref{lem:wrel_add}, the addition/subtraction preserves the correspondence
  (using bounds $|a[j]|, |\op{mont}(\zeta \cdot a[j+\ell])| < 2^{16}$, so
  their sum/difference is bounded by $2^{17} < 2^{30}$, safe in 32-bit words).
\end{proof}

\section{Loop Equivalence Proofs}

\subsection{Inner Loop Equivalence}

\begin{lemma}[Inner loop correspondence]
  Given $\ec{wrel}(\texttt{rp}, \texttt{a})$ and bounds $|a[i]| < 2^{15}$:
  \[
    \ec{hoare}[\ec{Extracted.inner\_loop} \sim \ec{NTT\_Fq.inner\_loop} :
      \ec{wrel}(\texttt{rp}, \texttt{a}) \Rightarrow \ec{wrel}(\texttt{rp}', \texttt{a}')].
  \]
\label{lem:inner_corr}
\end{lemma}

\begin{proof}
  By induction on the inner loop counter $j$, using \Cref{lem:butterfly_corr}
  as the inductive step.  The bound on coefficients increases by at most
  $2^{16}$ per butterfly (from \Cref{thm:mont_correct}), but the inner loop
  uses the same twiddle factor throughout, so the number of butterflies per
  inner loop execution is exactly $\ell$.
\end{proof}

\subsection{Full NTT Loop Equivalence}

\begin{theorem}[Forward NTT loop equivalence]
  Let $\texttt{rp}$ satisfy $|a[i]| < 2^{15}$ for all $i$.  Then:
  \begin{align*}
    &\ec{hoare}[\ec{Extracted.poly\_ntt\_jazz} \sim \ec{NTT\_Fq.NTT.ntt} : \\
    &\quad \ec{wrel}(\texttt{rp}, \texttt{a}) \\
    &\quad \Rightarrow \ec{wrel}(\texttt{rp}', \texttt{a}') \;\wedge\;
     \forall i.\; |a'[i]| < 2^{23}].
  \end{align*}
\label{thm:ntt_loop_equiv}
\end{theorem}

\begin{proof}
  By three-level loop induction:
  \begin{enumerate}
    \item The outer loop induction uses $\ell$ as the decreasing variant.
          The invariant after completing all stages with $\ell' > \ell$ is
          the partial NTT condition.
    \item The middle loop induction uses $256 - \texttt{start}$ as the
          decreasing variant.
    \item The inner loop uses \Cref{lem:inner_corr}.
  \end{enumerate}
  The bound $|a'[i]| < 2^{23}$ follows from \Cref{thm:ntt_bounds}.
\end{proof}

\section{Main Refinement Lemmas}

The two main lemmas of \texttt{RefJasminNTT.ec} are:

\begin{lstlisting}[language=EasyCrypt,
  caption={Main refinement lemmas in \texttt{RefJasminNTT.ec}},
  label={lst:ref_lemmas}]
(* Forward NTT: Jasmin core ~ NTT_Fq, bounds 16 -> 24 *)
lemma poly_ntt_core_ref rp a :
  wrel rp a =>
  (forall i, 0 <= i < 256 => `|a.[i]| < 2^15) =>
  equiv[Extracted.poly_ntt_jazz ~ NTT_Fq.NTT.ntt :
    ={arg} /\ wrel arg{1} arg{2} ==>
    wrel res{1} res{2} /\
    forall i, 0 <= i < 256 => `|res{2}.[i]| < 2^23].

(* Inverse NTT: Jasmin core ~ NTT_Fq, bounds 16 -> 16, Montgomery *)
lemma poly_invntt_core_ref rp a :
  wrel rp a =>
  (forall i, 0 <= i < 256 => `|a.[i]| < 2^23) =>
  equiv[Extracted.poly_invntt_jazz ~ NTT_Fq.NTT.invntt :
    ={arg} /\ wrel arg{1} arg{2} ==>
    wrel res{1} res{2} /\
    forall i, 0 <= i < 256 => `|res{2}.[i]| < 2^15].
\end{lstlisting}

These two lemmas represent the central contribution of \texttt{RefJasminNTT.ec}.
Once proved, they serve as the foundation for the end-to-end theorem.

\section{Proof Engineering Challenges}

\subsection{Signed vs.\ Unsigned Arithmetic}

The \jasmin{} code uses \jasm{u32} (unsigned) for coefficient storage but
\jasm{(64s)(x)} (signed 64-bit) for intermediate products.  The
\easycrypt{} extraction mirrors this by using \ec{W32.t} for storage and
\ec{W64.t} for products.  The refinement proof must track the signed
interpretation of unsigned words, which requires the signed-arithmetic
lemmas from \texttt{Fq.ec}.

\subsection{64-bit Product Non-Overflow}

The 64-bit product $\zeta \cdot a[j+\ell]$ requires:
$|\zeta| \cdot |a[j+\ell]| < 2^{63}$.
Since $|\zeta| < q < 2^{16}$ and $|a[j+\ell]| < 2^{23}$ (after the forward
NTT), we have $|\zeta \cdot a[j+\ell]| < 2^{39} \ll 2^{63}$.
This is verified in the proof as a side condition.

\subsection{Twiddle Factor Table Agreement}

The \easycrypt{} extracted module accesses the twiddle table as a global
constant.  The refinement proof must show that this global constant agrees
with the \ec{zetas} array defined in \texttt{NTT\_Fq.ec}.  This is
established by a computation lemma:

\begin{lemma}[Twiddle table agreement]
  For all $k \in [0, 256)$:
  $\op{W32.to\_sint}(\ec{jzetas}[k]) = \ec{NTT\_Fq.zetas}[k]$.
\label{lem:zetas_agree}
\end{lemma}

\begin{proof}
  Both tables are defined by the same formula modulo $q = 64513$.  Since
  $q < 2^{16}$, all entries fit in 16 bits and the signed interpretation
  agrees with the unsigned value.  Verified by the SMT backend.
\end{proof}

\subsection{In-Place Update Semantics}

Both the \jasmin{} implementation and the \easycrypt{} specification modify
the input array in place.  The refinement proof tracks the state of each
array cell individually, using the \easycrypt{} array theory lemmas to reason
about functional updates.

The key identity is: after an in-place butterfly that modifies indices $j$ and
$j+\ell$, all other indices remain unchanged.  This is expressed as:
\[
  \forall i \ne j, i \ne j+\ell.\;
  \texttt{rp}'[i] = \texttt{rp}[i] \;\wedge\;
  \texttt{a}'[i] = \texttt{a}[i].
\]

%==============================================================================%
% Chapter 11 – End-to-End Correctness
%==============================================================================%

\chapter{End-to-End Correctness}
\label{ch:end_to_end}

\section{Overview of \texttt{NTTEndToEnd.ec}}

The file \texttt{NTTEndToEnd.ec} (144 lines) is the apex of the verification
chain.  It imports all preceding modules and states the final correctness
theorems that directly relate the exported \jasmin{} functions to the abstract
spectral NTT specification.

\begin{lstlisting}[language=EasyCrypt,
  caption={Module imports in \texttt{NTTEndToEnd.ec}},
  label={lst:e2e_imports}]
require import NTTFullSpec NTTFullAlgebra.
require import NTT_Fq Fq GFq Rq.
require import RefJasminNTTLoop CTLoopEquiv.
require import Hpoly_extract Hpoly_loop.
\end{lstlisting}

\section{Main Correctness Theorems}

\subsection{Forward NTT Correctness}

\begin{theorem}[Jasmin NTT correctness]
  \label{thm:main_ntt}
  Let \(\texttt{rp}\) be a Jasmin word array represented in EasyCrypt as a
  \ec{BArray1024.t}, encoding a polynomial \(p\) with 16-bit signed
  coefficients.  Then executing \jasm{poly\_ntt\_jazz(\texttt{rp})} produces
  an output array $\texttt{rp}'$ satisfying:
  \[
    \texttt{rp}' \;\text{represents}\; \op{full\_ntt}(p)
    \;\text{with 24-bit bounds},
  \]
\end{theorem}

\begin{lstlisting}[language=EasyCrypt,
  caption={End-to-end NTT correctness statement},
  label={lst:e2e_ntt}]
op jasmin_ntt_full_post (rp : BArray1024.t) (p : poly) =
  NTT_Fq.poly_repr_bound rp (NTTFullSpec.full_ntt p) 24.

lemma haetae_jasmin_ntt_correct p :
  hoare [Hpoly_extract.M.poly_ntt_jazz :
    NTT_Fq.poly_repr_bound rp p 16 ==>
    NTT_Fq.poly_repr_bound res (NTTFullSpec.full_ntt p) 24].
\end{lstlisting}

\subsection{Inverse NTT Correctness}

\begin{theorem}[Jasmin inverse NTT correctness]
  \label{thm:main_invntt}
  Let $\texttt{rp}$ encode a polynomial \(p\) satisfying the 16-bit bound
  required by the checked inverse kernel theorem.
  Then executing \jasm{poly\_invntt\_jazz(\texttt{rp})} produces an output
  array $\texttt{rp}'$ satisfying:
  \[
    \texttt{rp}' \;\text{represents}\;
      \ec{array256\_mont}(\op{full\_invntt}(p)),
  \]
  i.e. the inverse output is represented with one Montgomery factor \(R\).
\end{theorem}

\begin{lstlisting}[language=EasyCrypt,
  caption={End-to-end inverse NTT correctness statement},
  label={lst:e2e_invntt}]
op jasmin_invntt_full_post (rp : BArray1024.t) (p : poly) =
  NTT_Fq.poly_repr_bound rp
    (NTT_Fq.array256_mont (NTTFullSpec.full_invntt p)) 16.

lemma haetae_jasmin_invntt_correct p :
  hoare [Hpoly_extract.M.poly_invntt_jazz :
    NTT_Fq.poly_repr_bound rp p 16 ==>
    NTT_Fq.poly_repr_bound res
      (NTT_Fq.array256_mont (NTTFullSpec.full_invntt p)) 16].
\end{lstlisting}

The \ec{array256\_mont} postcondition is the precise Montgomery-form
convention: each abstract coefficient is multiplied by \(R\), not by
\(R^{-1}\).  The value \(R^{-1}\) appears internally when cancelling a
Montgomery multiplication, but it is not the final representation factor.

\section{Round-Trip Consequence}

\begin{proposition}[Abstract round-trip with Montgomery output]
  \label{thm:roundtrip}
  For any polynomial $a \in \Z_q^{256}$:
  \[
    \op{full\_invntt}(\op{full\_ntt}(a)) = a.
  \]
Consequently, a concrete inverse call whose input represents
\(\op{full\_ntt}(a)\) and satisfies the inverse theorem's representation
bound returns a word array representing \(\ec{array256\_mont}(a)\), i.e.,
\(a\cdot R\) componentwise.
\end{proposition}

\begin{proof}
  At the mathematical layer:
  \begin{align*}
    \op{full\_invntt}(\op{full\_ntt}(a))[i]
    &= a[i],
  \end{align*}
  by \Cref{thm:intt_correct}.  The concrete inverse theorem then wraps the
  result in \ec{array256\_mont}, giving the \(R\)-scaled representation.
\end{proof}

\Cref{thm:roundtrip} is a representation-level consequence, not an additional
top-level theorem in \ec{NTTEndToEnd.ec}.  The checked file states the forward
and inverse exported-function theorems separately, and the dissertation does
not cite a direct one-procedure in-place pipeline theorem.

\section{Proof Composition}

The end-to-end proofs are assembled from four chains:

\begin{enumerate}
  \item \textbf{Full-safe extraction $\to$ proof loop}:
        \texttt{CTLoopEquiv.ec} proves that the helperized full-safe
        \texttt{Hpoly\_extract} procedures implement the same transition as
        the proof-only direct-loop model in \texttt{Hpoly\_loop.ec}.
  \item \textbf{Proof loop $\to$ Imperative}: \Cref{thm:ntt_loop_equiv} is
        instantiated by \texttt{RefJasminNTT\allowbreak Loop.ec}.
  \item \textbf{Imperative $\to$ Abstract}: \Cref{thm:ntt_correct} from
        \texttt{NTTFullAlgebra.ec}.
  \item \textbf{Abstract inverse identity}: \Cref{thm:intt_correct} is the
        pen-and-paper field argument stated in this dissertation.  The
        checked EasyCrypt files prove the forward and inverse maps against
        \ec{full\_ntt} and \ec{full\_invntt} separately; they do not add a
        standalone checked round-trip lemma.
\end{enumerate}

The composition is expressed in \easycrypt{} as sequential lemma application:

\begin{lstlisting}[language=EasyCrypt,
  caption={Proof composition in \texttt{NTTEndToEnd.ec}},
  label={lst:e2e_proof}]
proof haetae_jasmin_ntt_correct.
  move => rp a ha.
  (* Step 1: proof-only core loop -> imperative -> abstract. *)
  have hloop_core :
    hoare [Hpoly_loop.M._poly_ntt :
      NTT_Fq.poly_repr_bound rp a 16 ==>
      jasmin_ntt_full_post res a].
  + by conseq RefJasminNTT.poly_ntt_core_ref
      (NTTFullAlgebra.ntt_full_ntt a) => /#.
  (* Step 2: proof-only exported wrapper. *)
  have hloop_jazz :
    hoare [Hpoly_loop.M.poly_ntt_jazz :
      NTT_Fq.poly_repr_bound rp a 16 ==>
      jasmin_ntt_full_post res a].
  + by conseq RefJasminNTT.poly_ntt_jazz_ref hloop_core => /#.
  (* Step 3: helperized full-safe extraction -> proof-only wrapper. *)
  by conseq CTLoopEquiv.poly_ntt_jazz_ct_loop hloop_jazz => /#.
qed.
\end{lstlisting}

\section{Status of the Proof}

\subsection{Completed Machine-Checked Steps}

\begin{enumerate}
  \item \textbf{\texttt{Fq.ec}}: All lemmas proved, including Montgomery
        reduction correctness, signed arithmetic identities, and bit-width
        bound propagation.

  \item \textbf{\texttt{NTT\_Fq.ec}}: Full imperative specification including
        termination proofs and loop invariants.

  \item \textbf{\texttt{NTTFullSpec.ec}}: Abstract NTT definitions and
        conversion lemmas from the relational specifications to
        \ec{full\_ntt} and \ec{full\_invntt}.

  \item \textbf{\texttt{RefJasminNTT.ec}}: The preserved original direct-loop
        proof boundary and its two main refinement lemmas
        \ec{poly\_ntt\_core\_ref} and \ec{poly\_invntt\_core\_ref}.

  \item \textbf{\texttt{Hpoly\_loop.ec}}: The proof-only direct-loop model
        used as the stable target for the full-safe helperized extraction.

  \item \textbf{\texttt{RefJasminNTT\allowbreak Loop.ec}}: The direct-loop refinement
        proof retargeted to \texttt{Hpoly\_loop}.

  \item \textbf{\texttt{CTLoopEquiv.ec}}: The pRHL bridge proving equivalence
        between the full-safe \texttt{Hpoly\_extract} procedures and
        \texttt{Hpoly\_loop}, including exported wrappers.

  \item \textbf{\texttt{NTTFullAlgebra.ec}}: The stage-induction theorems
        \ec{ntt\_full\_ntt} and \ec{invntt\_full\_invntt}.

  \item \textbf{\texttt{NTTEndToEnd.ec}}: The composed top-level Hoare
        statements for the exported Jasmin functions.
\end{enumerate}

\subsection{Machine-Checked Top-Level Compilation}

The checked development currently compiles through \texttt{NTTEndToEnd.ec} in
\texttt{easycrypt-ct}.  This means the full-safe helperized Jasmin extraction,
the proof-loop refinement, the imperative-to-abstract stage proofs, and the
final composition are all discharged for the scalar NTT and inverse NTT
covered by this dissertation.  Remaining limitations are scope limitations,
not open proof gaps in the NTT chain: AVX2 vectorisation, base multiplication,
and full protocol-level security integration are outside the proved artifact.

\section{Summary of Verification Metrics}

\begin{table}[htbp]
\centering
\caption{Verification metrics by file}
\label{tab:metrics}
\small
\begin{tabular}{
  >{\raggedright\arraybackslash}p{0.32\textwidth}
  rr
  >{\raggedright\arraybackslash}p{0.28\textwidth}}
\toprule
File & Lines & Named lemmas & Status \\
\midrule
\texttt{Fq.ec}           & 535  & 16  & Compiled \\
\texttt{GFq.ec}          & 27   & 0   & Compiled field instantiation \\
\texttt{Rq.ec}           & 117  & 0   & Compiled ring interface \\
\texttt{Montgomery.ec}   & 427  & 29  & Compiled \\
\texttt{NTT\_Fq.ec}      & 988  & 52  & Compiled \\
\texttt{NTTFullSpec.ec}  & 160  & 10  & Compiled \\
\texttt{NTTFullAlgebra.ec} & 3085 & 111 & Compiled \\
\texttt{RefJasminNTT.ec} & 2609 & 103 & Preserved direct-loop proof \\
\texttt{Hpoly\_loop.ec} & 156 & 0 & Compiled proof-loop model \\
\texttt{RefJasminNTT\allowbreak Loop.ec} & 2629 & 105 & Compiled full-safe refinement target \\
\texttt{CTLoopEquiv.ec} & 657 & 14 & Compiled helper bridge \\
\texttt{NTTEndToEnd.ec}  & 144  & 10  & Compiled full-safe top-level \\
\midrule
\textbf{Total}           & \textbf{11534} & \textbf{450} & \textbf{Compiled through full-safe end-to-end} \\
\bottomrule
\end{tabular}
\normalsize
\end{table}

%==============================================================================%
% Chapter 12 – Conclusion
%==============================================================================%

\chapter{Conclusion}
\label{ch:conclusion}

\section{Summary of Contributions}

This dissertation has presented a comprehensive formal verification of the
Number-Theoretic Transform in the \haetae{} post-quantum signature scheme.
The verification spans four abstraction layers---from \jasmin{} machine code
to the abstract spectral evaluation formula---and is expressed entirely in the
\easycrypt{} proof assistant.

\paragraph{C1 -- Montgomery correctness.}
We proved that the signed Montgomery reduction in \texttt{reduce.jinc}
correctly computes $a \cdot R^{-1} \bmod q$ within the working range of the
NTT, with explicit bounds on all intermediate values.  The proof required
new signed-arithmetic lemmas for \easycrypt{} that are not present in the
standard library and that are expected to be of independent interest for
future cryptographic verifications.

\paragraph{C2 -- Imperative NTT specification.}
We developed a fully proved imperative NTT specification in \texttt{NTT\_Fq.ec},
covering both the forward and inverse NTT with all loop invariants,
termination witnesses, and coefficient-bound invariants.  The specification is
structurally identical to the C reference implementation and the \jasmin{} source,
making the correspondence proofs straightforward once the arithmetic
infrastructure is in place.

\paragraph{C3 -- Abstract spectral specification.}
We developed the algebraic infrastructure in \texttt{NTTFullAlgebra.ec},
including the primitive root properties of $\zeta_0 = 426$ in $\F_{64513}$,
the orthogonality relation for the negacyclic NTT, bit-reversal identities,
and the butterfly-stage decomposition theorem.  This library distinguishes
the \haetae{} full $256$-point transform from the MLKEM incomplete
$128$-point transform, and provides the formal foundation for the
stage-induction proof.

\paragraph{C4 -- Jasmin refinement.}
We proved the two central refinement lemmas in \texttt{RefJasminNTT.ec}:
\begin{itemize}
  \item \ec{poly\_ntt\_core\_ref}: the \jasmin{} forward NTT is observationally
        equivalent to the imperative specification, with coefficient bounds
        progressing from 16-bit inputs to 24-bit outputs.
  \item \ec{poly\_invntt\_core\_ref}: the \jasmin{} inverse NTT is
        observationally equivalent to the imperative inverse NTT specification,
        preserving 16-bit coefficient bounds in the output.
\end{itemize}

\paragraph{C5 -- End-to-end correctness.}
We assembled all layers in \texttt{NTTEndToEnd.ec}, stating the main theorems
\ec{haetae\_jasmin\_ntt\_correct} and \ec{haetae\_jasmin\_invntt\_correct}.
These theorems directly relate the \jasmin{} exported functions to the abstract
spectral NTT map.

\section{Limitations and Remaining Obligations}

\subsection{Stage Induction Gap}

The one remaining proof obligation is the complete machine-checked proof of the
butterfly-stage induction connecting \ec{NTT\_Fq.NTT.ntt} to
\ec{NTTFullSpec.full\_ntt}.  While the mathematical argument is clear
(\Cref{lem:butterfly_stage}), its formalisation in \easycrypt{} requires
significant index-arithmetic manipulation.  The MLKEM verification did not
provide a reusable template, necessitating new infrastructure.  Completing
this proof is the primary remaining task.

\subsection{No AVX2 Vectorisation}

The current \jasmin{} implementation is scalar (no SIMD).  A production
\haetae{} implementation would use AVX2 intrinsics for a significant
speedup.  Verifying the correctness of an AVX2 NTT in \easycrypt{} would
require extending the proof to handle 256-bit registers and vectorised butterfly
operations.

\subsection{Base Multiplication Not Fully Verified}

The file \texttt{hpoly.jazz} also exports \jasm{poly\_basemul\_jazz}, the
pointwise NTT-domain multiplication.  While the \jasmin{} implementation is
available, its end-to-end correctness proof is not yet complete.  This is
the complement to the NTT correctness: together, the NTT and base
multiplication allow proving that \jasm{poly\_basemul\_jazz} correctly
computes polynomial multiplication in $R_q$.

\subsection{Security Proof Integration}

This dissertation establishes functional correctness only.  Connecting the
NTT implementation to the \emph{security} proof of \haetae{} (i.e., the
reduction from DMSIS) would require additional work at the cryptographic
protocol level.

\section{Future Work}

\paragraph{Completing the stage induction.}
The most pressing future work is completing the butterfly-stage induction in
\texttt{NTTFullAlgebra.ec}.  The key insight needed is a formal lemma relating
the loop counter $k$ to the spectral index $\brev_8(2^s + g)$ at each stage.
We conjecture that a more direct formulation---expressing the partial NTT
directly in terms of the loop variables rather than through the bit-reversal
permutation---would simplify the proof considerably.

\paragraph{Vectorised implementation verification.}
Extending the proof to cover an AVX2 implementation would provide a
performance-verified \haetae{} NTT.  The \jasmin{} language has full support
for AVX2 intrinsics, and the \easycrypt{} extraction handles 256-bit types.
The main challenge is reasoning about the parallelism of 8 butterfly
operations executed simultaneously.

\paragraph{Reusable NTT algebra library.}
The algebraic infrastructure developed in \texttt{NTTFullAlgebra.ec} is
specific to $n = 256$ and $q = 64513$.  A parameterised version, expressing
the NTT algebra in terms of an abstract field $\F_q$ and an abstract
$2n$th root of unity, would be reusable across lattice-based schemes.

\paragraph{Full \haetae{} signing and verification.}
The ultimate goal is a fully verified \haetae{} implementation: key generation,
signing, and verification, all the way from the \jasmin{} source to the
DMSIS security proof.  The NTT verification of this dissertation is the
foundational building block for that larger project.

\paragraph{Comparison with Dilithium and FALCON.}
A comparative study of the NTT verification strategies for Dilithium, FALCON,
and \haetae{} would identify opportunities for sharing algebraic infrastructure
and for developing a unified ``NTT verification toolkit'' for
\easycrypt{}/\jasmin{}.

\section{Closing Remarks}

The formal verification of cryptographic implementations is no longer a
distant aspiration---it is an engineering practice that can be applied to
production code.  This dissertation demonstrates that a full-featured
lattice-based NTT, with all its integer arithmetic nuances, can be formally
verified in \easycrypt{} by careful layering of specifications and proofs.

The main lesson is the value of the three-layer architecture: separating the
word-level arithmetic concerns (Layer 1) from the algorithmic concerns
(Layer 2) from the mathematical concerns (Layer 3) makes each proof layer
tractable.  The interface between layers is clean: word arrays vs.\ integer
arrays (Layers 1--2) and loop invariants vs.\ evaluation formulas (Layers 2--3).

As post-quantum cryptography moves from standardisation to deployment, the
tools and techniques developed here---\jasmin{} for high-assurance
implementations, \easycrypt{} for machine-checked proofs, and the three-layer
verification architecture---provide a clear path to verifying not just NTTs
but complete post-quantum signature and key-encapsulation schemes.


%--- Appendix: TikZ Figure Gallery -------------------------------------------
\appendix
\chapter{Visualisation Gallery}
\label{app:figures}
\addcontentsline{toc}{chapter}{Appendix: Visualisation Gallery}

The following ten figures are provided as standalone reference visualisations
to support the narrative in the main chapters.  Each figure is also
reproduced inline at the point where it is first referenced.

%==============================================================================%
% tikz_figures.tex  –  All standalone TikZ visualisations for the dissertation
%
% Each figure is wrapped in \begin{figure}...\end{figure} so it can be
% \input{}–ed from any chapter, or the whole file can be \include{}–d as a
% stand-alone "Figures" appendix.
%==============================================================================%

% ── Helper commands (defined globally so they work inside tikzpicture) ────────
% NTT stage header: #1=x-pos, #2=label, #3=bound
\newcommand{\stagehead}[3]{%
  \node[font=\footnotesize\bfseries\sffamily, align=center] at (#1, 1.35) {#2};%
  \node[font=\scriptsize\ttfamily, align=center, col@darkline]  at (#1, 0.1) {#3};%
}
% NTT group cell: #1=stage-x, #2=group-idx, #3=ydelta(=0.68*g), #4=color
\newcommand{\gcell}[4]{%
  \fill[#4, draw=col@darkline!50, rounded corners=1pt]
    (#1-0.38, -3.45+#3) rectangle (#1+0.38, -3.45+#3+0.58);%
}
% Bounds table row: #1=y, #2=shade(0/1), #3=layer, #4=fwdbounds, #5=invbounds, #6=roundtrip
% Column x: layer=0.0, fwd=2.8, inv=6.5, rt=9.8+0.7=10.5
\newcommand{\bndrow}[6]{%
  \ifthenelse{\equal{#2}{1}}{%
    \fill[col@butterfly!18] (-0.2,#1-0.37) rectangle (11.4,#1+0.37);}{}%
  \node[font=\footnotesize\sffamily, anchor=west]  at (0.0, #1) {#3};%
  \node[font=\scriptsize\ttfamily,   anchor=center, col@highlight]       at (2.8, #1) {#4};%
  \node[font=\scriptsize\ttfamily,   anchor=center, NavyBlue!70!black]   at (6.5, #1) {#5};%
  \node[font=\scriptsize\sffamily\bfseries, anchor=center, OliveGreen!60!black] at (10.5,#1) {#6};%
}

% ── Colour palette ───────────────────────────────────────────────────────────
\colorlet{col@jasmin}   {red!18!white}
\colorlet{col@extract}  {orange!22!white}
\colorlet{col@imperative}{yellow!30!white}
\colorlet{col@algebraic}{green!18!white}
\colorlet{col@endtoend} {cyan!18!white}
\colorlet{col@proof}    {violet!15!white}
\colorlet{col@butterfly}{blue!12!white}
\colorlet{col@arrow}    {black!70}
\colorlet{col@highlight}{BrickRed}
\colorlet{col@darkline} {black!55}

% ── Common TikZ styles ───────────────────────────────────────────────────────
\tikzset{
  layer box/.style = {
    rectangle, draw=col@darkline, rounded corners=4pt,
    minimum width=#1, minimum height=1.15cm,
    align=center, font=\small\sffamily,
    drop shadow={opacity=.12,shadow xshift=1pt,shadow yshift=-1pt},
  },
  layer box/.default = 9cm,
  proof arrow/.style = {
    -Stealth, thick, col@arrow,
  },
  annot/.style = {
    font=\footnotesize\itshape, col@arrow, align=center,
  },
  node wire/.style = {thick, col@arrow},
  butterfly in/.style  = {circle, draw, fill=col@butterfly,
                           minimum size=5pt, inner sep=0pt},
  butterfly out/.style = {circle, draw, fill=col@butterfly!60!black,
                           minimum size=5pt, inner sep=0pt},
}

%==============================================================================%
%  FIGURE 1 — Full Refinement Chain (expanded, annotated)
%==============================================================================%
\begin{figure}[htbp]
\centering
\begin{tikzpicture}[x=1cm, y=1cm, scale=0.88, transform shape]

  %--- Layer boxes -------------------------------------------------------------
  \node[layer box=10.5cm, fill=col@jasmin]
    (L0) at (0, 7.5)
    {\textbf{Layer 0 — Jasmin Source Code}\\[2pt]
     \texttt{poly.jinc}\enspace|\enspace
     \texttt{reduce.jinc}\enspace|\enspace
     \texttt{zetas.jinc}\enspace|\enspace
     \texttt{hpoly.jazz}};

  \node[layer box=10.5cm, fill=col@extract]
    (L1) at (0, 5.4)
    {\textbf{Layer 1 — Extracted EasyCrypt Module}\\[2pt]
     \texttt{Hpoly\_extract.ec}\enspace
     (\texttt{BArray1024.t} operands, machine semantics)};

  \node[layer box=10.5cm, fill=col@imperative]
    (L2) at (0, 3.3)
    {\textbf{Layer 2 — Imperative EasyCrypt Specification}\\[2pt]
     \texttt{NTT\_Fq.ec}\enspace|\enspace
     \texttt{Fq.ec}\enspace|\enspace
     \texttt{Montgomery.ec}
     (\texttt{coeff Array256.t}, loop invariants)};

  \node[layer box=10.5cm, fill=col@algebraic]
    (L3) at (0, 1.2)
    {\textbf{Layer 3 — Abstract Algebraic / Spectral Specification}\\[2pt]
     \texttt{NTTFullSpec.ec}\enspace|\enspace
     \texttt{NTTFullAlgebra.ec}\enspace
     (\(\mathbb{F}_q^{256}\), evaluation formula)};

  %--- Proof step arrows -------------------------------------------------------
  \draw[proof arrow, line width=2pt]
    (L0.south) -- node[right=6pt, annot]
      {\textbf{Step A}\\Jasmin compiler\\extraction\\(machine-checked)}
    (L1.north);

  \draw[proof arrow, line width=2pt]
    (L1.south) -- node[right=6pt, annot]
      {\textbf{Step B}\\pRHL equivalence\\(\texttt{RefJasminNTT.ec})\\word $\leftrightarrow$ int}
    (L2.north);

  \draw[proof arrow, line width=2pt]
    (L2.south) -- node[right=6pt, annot]
      {\textbf{Step C}\\Loop invariant\\induction\\(\texttt{NTTFullAlgebra.ec})}
    (L3.north);

  %--- End-to-end composition --------------------------------------------------
  \node[layer box=4.2cm, fill=col@endtoend, minimum height=2.2cm]
    (E2E) at (-7.5, 4.35)
    {\textbf{End-to-End}\\[3pt]
     \texttt{NTTEndToEnd.ec}\\[3pt]
     \footnotesize
     \texttt{haetae\_jasmin\_ntt}\\
     \texttt{\_correct}\\[2pt]
     $A \circ B \circ C$};

  \draw[proof arrow, dashed, bend left=25]
    (L0.west) to (E2E.north);
  \draw[proof arrow, dashed, bend right=25]
    (L3.west) to (E2E.south);

  %--- Side annotations --------------------------------------------------------
  \node[annot, text=col@darkline, anchor=east] at (-5.7, 7.5)
    {x86-64 assembly\\+ safety proof};
  \node[annot, text=col@darkline, anchor=east] at (-5.7, 5.4)
    {\texttt{W32}/\texttt{W64}\\machine words};
  \node[annot, text=col@darkline, anchor=east] at (-5.7, 3.3)
    {\texttt{int} arithmetic\\16-bit $\to$ 24-bit bounds};
  \node[annot, text=col@darkline, anchor=east] at (-5.7, 1.2)
    {$\hat{f}[i]=\!\!\sum_{j}\!f[j]\,\zeta^{(2\,\mathrm{br}(i)+1)j}$};

\end{tikzpicture}
\caption{Four-layer verification architecture for the \haetae{} NTT.  Solid
         arrows are checked proof steps; dashed arrows show final composition.}
\label{fig:full_arch}
\end{figure}

%==============================================================================%
%  FIGURE 2 — Cooley–Tukey NTT Butterfly (single stage, 8-element example)
%==============================================================================%
\begin{figure}[htbp]
\centering
\begin{tikzpicture}[x=1.5cm, y=0.9cm]

  \def\nleft{4}   % number of element pairs shown

  %--- Input labels (left) ----------------------------------------------------
  \foreach \i/\lab in {
      0/{$a[0]$},
      1/{$a[1]$},
      2/{$a[2]$},
      3/{$a[3]$},
      4/{$a[4]$},
      5/{$a[5]$},
      6/{$a[6]$},
      7/{$a[7]$}}
  {
    \node[anchor=east, font=\small\ttfamily] (in\i) at (0, -\i) {\lab};
  }

  %--- Output labels (right) --------------------------------------------------
  \foreach \i/\lab in {
      0/{$a'[0]$},
      1/{$a'[1]$},
      2/{$a'[2]$},
      3/{$a'[3]$},
      4/{$a'[4]$},
      5/{$a'[5]$},
      6/{$a'[6]$},
      7/{$a'[7]$}}
  {
    \node[anchor=west, font=\small\ttfamily] (out\i) at (4, -\i) {\lab};
  }

  %--- Stage 1 (len=4): butterfly pairs (0,4),(1,5),(2,6),(3,7) ---------------
  \def\sx{1.0}   % x of stage-1 nodes
  \def\ex{2.4}   % x of stage-1 output nodes

  % Butterfly nodes
  \foreach \i in {0,1,2,3} {
    \pgfmathtruncatemacro{\j}{\i+4}
    \node[butterfly in]  (b1up\i)  at (\sx, -\i)  {};
    \node[butterfly out] (b1lo\i)  at (\sx, -\j)  {};
    \node[butterfly in]  (b1oup\i) at (\ex, -\i)  {};
    \node[butterfly out] (b1olo\i) at (\ex, -\j)  {};
  }

  % Straight wires from input
  \foreach \i in {0,1,2,3} {
    \draw[node wire] (in\i) -- (b1up\i);
  }
  \foreach \i in {0,1,2,3} {
    \pgfmathtruncatemacro{\j}{\i+4}
    \draw[node wire] (in\j) -- (b1lo\i);
  }

  % Cross wires (butterfly)
  \foreach \i in {0,1,2,3} {
    \pgfmathtruncatemacro{\j}{\i+4}
    % upper output = upper in + w * lower in
    \draw[node wire, ->, BrickRed!70!black]
      (b1up\i) -- (b1oup\i);
    \draw[node wire, ->, BrickRed!70!black]
      (b1lo\i) to[out=0, in=180] (b1oup\i);
    % lower output = upper in - w * lower in
    \draw[node wire, ->, NavyBlue!70!black]
      (b1up\i) to[out=0, in=180] (b1olo\i);
    \draw[node wire, ->, NavyBlue!70!black]
      (b1lo\i) -- (b1olo\i);
  }

  % Twiddle factor labels
  \foreach \i/\wlab in {0/{$w_0$},1/{$w_1$},2/{$w_2$},3/{$w_3$}} {
    \pgfmathtruncatemacro{\j}{\i+4}
    \node[font=\scriptsize, col@highlight, above=1pt]
      at ($(b1lo\i)!0.5!(b1oup\i)+(0,0.35)$) {\wlab};
  }

  % Output wires
  \foreach \i in {0,1,2,3} {
    \pgfmathtruncatemacro{\j}{\i+4}
    \draw[node wire] (b1oup\i) -- (out\i);
    \draw[node wire] (b1olo\i) -- (out\j);
  }

  %--- Butterfly operation annotation -----------------------------------------
  \node[draw=col@highlight, rounded corners=3pt, fill=col@highlight!8!white,
        align=left, font=\footnotesize, inner sep=6pt]
    (bfann) at (2.0, -8.6)
    {%
      \textbf{Cooley--Tukey butterfly:}\\
      $t \leftarrow w_k \cdot a[j+\ell]$\\
      $a'[j+\ell] \leftarrow a[j] - t$\\
      $a'[j] \leftarrow a[j] + t$\\[3pt]
      {\color{BrickRed!80!black}$\blacksquare$}\ upper path \;
      {\color{NavyBlue!80!black}$\blacksquare$}\ lower path
    };

  %--- Stage label ------------------------------------------------------------
  \node[font=\small\bfseries\sffamily, col@darkline] at (1.7, 0.8)
    {Stage $s$ with $\ell=4$, twiddle counter $k$};

  \draw[decorate, decoration={brace, amplitude=6pt, mirror},
        col@darkline, thick]
    (-0.3, 0.2) -- (-0.3, -3.8)
    node[midway, left=8pt, font=\scriptsize, align=right]
      {upper\\half};
  \draw[decorate, decoration={brace, amplitude=6pt, mirror},
        col@darkline, thick]
    (-0.3, -4.2) -- (-0.3, -7.8)
    node[midway, left=8pt, font=\scriptsize, align=right]
      {lower\\half};

\end{tikzpicture}
\caption{One stage of the Cooley--Tukey butterfly for $n=8$ (shown
         for clarity; the actual NTT uses $n=256$ over 8 stages).
         Red paths carry the $+t$ update to the upper half;
         blue paths carry the $-t$ update to the lower half.
         Each pair of red+blue paths forms one butterfly using twiddle
         factor $w_k = \zeta_0^{\mathrm{brev}_8(k)} \cdot R \bmod q$.}
\label{fig:butterfly}
\end{figure}

%==============================================================================%
%  FIGURE 3 — 8-Stage HAETAE NTT Data Flow (len progression)
%==============================================================================%
\begin{figure}[htbp]
\centering
\begin{tikzpicture}[x=1.75cm, y=0.85cm, scale=0.94, transform shape]
  % Background
  \fill[col@butterfly!28, rounded corners=4pt] (-0.45,-3.6) rectangle (8.45,1.75);
  \draw[col@darkline, dashed] (-0.45,0.40) -- (8.45,0.40);
  \draw[col@darkline, dashed] (-0.45,0.88) -- (8.45,0.88);

  % Stage column headers (hardcoded via \stagehead)
  \stagehead{0}{Input}{16-bit}
  \stagehead{1}{$\ell=128$}{$<\!2^{16}$}
  \stagehead{2}{$\ell=64$} {$<\!2^{17}$}
  \stagehead{3}{$\ell=32$} {$<\!2^{18}$}
  \stagehead{4}{$\ell=16$} {$<\!2^{19}$}
  \stagehead{5}{$\ell=8$}  {$<\!2^{20}$}
  \stagehead{6}{$\ell=4$}  {$<\!2^{21}$}
  \stagehead{7}{$\ell=2$}  {$<\!2^{22}$}
  \stagehead{8}{$\ell=1$}  {23-bit}
  \node[font=\scriptsize\bfseries, text=col@highlight, align=center]
    at (8.00, 1.38) {post\\24-bit};

  % Twiddle counts
  \foreach \s/\cnt in {0/{---},1/{1},2/{2},3/{4},4/{8},5/{16},6/{32},7/{64},8/{128}} {
    \node[font=\scriptsize, col@highlight, align=center] at (\s, 0.65) {$k\mathord{+}$\cnt};
  }

  % Group state grid: 8 groups (g0..g7) × 9 stages
  % A group is "done" (green) at stage s if g < 2^s / 32,
  % i.e., s0: none; s1: g<1=none; s2: g<2; s3: g<4; etc.
  % done(s,g): s=0→0, s=1→g<1, s=2→g<2, s=3→g<4, s=4→g<8 (all)
  % Simplified: done if g <= (2^min(s,3))-2 for s<=3, all for s>=4
  % Group cells: white = pending, green = done
  % s=0..4: all white; s=5: g0 done; s=6: g0,g1; s=7: g0-g3; s=8: all
  \foreach \s in {0,...,4} {
    \foreach \g in {0,...,7} {
      \pgfmathsetmacro{\gyd}{0.68*\g}
      \gcell{\s}{\g}{\gyd}{white}
      \node[font=\tiny\ttfamily, col@darkline!55] at (\s,-3.45+\gyd+0.29) {g\g};
    }
  }
  % s=5
  \pgfmathsetmacro{\gyd}{0.68*0} \gcell{5}{0}{\gyd}{OliveGreen!45}
  \node[font=\tiny\ttfamily,OliveGreen!40!black] at (5,-3.45+\gyd+0.29) {g0};
  \foreach \g in {1,...,7} {
    \pgfmathsetmacro{\gyd}{0.68*\g}
    \gcell{5}{\g}{\gyd}{white}
    \node[font=\tiny\ttfamily, col@darkline!55] at (5,-3.45+\gyd+0.29) {g\g};
  }
  % s=6
  \foreach \g in {0,1} {
    \pgfmathsetmacro{\gyd}{0.68*\g} \gcell{6}{\g}{\gyd}{OliveGreen!45}
    \node[font=\tiny\ttfamily,OliveGreen!40!black] at (6,-3.45+\gyd+0.29) {g\g};
  }
  \foreach \g in {2,...,7} {
    \pgfmathsetmacro{\gyd}{0.68*\g}
    \gcell{6}{\g}{\gyd}{white}
    \node[font=\tiny\ttfamily, col@darkline!55] at (6,-3.45+\gyd+0.29) {g\g};
  }
  % s=7
  \foreach \g in {0,...,3} {
    \pgfmathsetmacro{\gyd}{0.68*\g} \gcell{7}{\g}{\gyd}{OliveGreen!45}
    \node[font=\tiny\ttfamily,OliveGreen!40!black] at (7,-3.45+\gyd+0.29) {g\g};
  }
  \foreach \g in {4,...,7} {
    \pgfmathsetmacro{\gyd}{0.68*\g}
    \gcell{7}{\g}{\gyd}{white}
    \node[font=\tiny\ttfamily, col@darkline!55] at (7,-3.45+\gyd+0.29) {g\g};
  }
  % s=8: all done
  \foreach \g in {0,...,7} {
    \pgfmathsetmacro{\gyd}{0.68*\g} \gcell{8}{\g}{\gyd}{OliveGreen!45}
    \node[font=\tiny\ttfamily,OliveGreen!40!black] at (8,-3.45+\gyd+0.29) {g\g};
  }

  % Stage transition arrows
  \foreach \s in {0,...,7} {
    \pgfmathtruncatemacro{\sn}{\s+1}
    \draw[proof arrow, OliveGreen!70!black, line width=1.4pt]
      (\s+0.41,-1.40) -- (\sn-0.41,-1.40);
    \node[font=\tiny\sffamily, col@darkline] at (\s+0.5,-1.75) {s\sn};
  }

  % Row labels
  \node[font=\scriptsize\bfseries\sffamily, col@darkline] at (-0.02, 1.60) {Stage};
  \node[font=\scriptsize, col@highlight]                  at (-0.02, 0.65) {twiddles};
  \node[font=\scriptsize\ttfamily, col@darkline]          at (-0.02, 0.10) {bound};

  % Legend
  \node[draw=col@darkline!60, rounded corners=3pt, fill=white,
        font=\scriptsize, inner sep=5pt, align=left]
    at (6.80,-2.90)
    {\textcolor{OliveGreen!60!black}{$\blacksquare$} done group\quad
     $\square$ pending};

\end{tikzpicture}
\caption{Data-flow view of the eight-stage forward NTT\@.  Each column
         is one outer-loop iteration ($\ell$ value).  Green cells are groups
         fully processed; white cells are untouched.  The checked coefficient
         invariant starts with a 16-bit input bound and exits with a 24-bit
         postcondition.
         The twiddle counter column ($k{+}$) shows how many twiddle factors
         have been consumed cumulatively.}
\label{fig:ntt_stages}
\end{figure}

%==============================================================================%
%  FIGURE 4 — Montgomery Reduction Dataflow
%==============================================================================%
\begin{figure}[htbp]
\centering
\begin{tikzpicture}[
  box/.style   = {rectangle, draw=col@darkline, rounded corners=3pt,
                  minimum width=2.8cm, minimum height=0.85cm,
                  align=center, font=\small\sffamily,
                  fill=col@imperative},
  op/.style    = {circle, draw=BrickRed!80!black, fill=BrickRed!10,
                  minimum size=0.7cm, font=\small\bfseries},
  wire/.style  = {-Stealth, thick, col@arrow},
  label/.style = {font=\scriptsize, col@darkline},
]

  % Input
  \node[box, minimum width=3.2cm] (inp) at (0, 0) {input $a$ (64-bit)\\$|a| < q \cdot 2^{15}$};

  % Step 1: truncate to 32-bit
  \node[box] (trunc) at (0, -1.9) {truncate $a$\\to 32 bits};
  \draw[wire] (inp.south) -- node[right=3pt, annot]{$(32u)(a)$} (trunc.north);

  % Step 2: multiply by QINV
  \node[op] (mul1) at (2.5, -1.9) {$\times$};
  \node[box, minimum width=2.4cm] (qinv) at (2.5, -0.7) {\texttt{QINV}\\$940508161$};
  \draw[wire] (trunc.east) -- (mul1.west);
  \draw[wire] (qinv.south) -- (mul1.north);

  % Step 3: low 32 bits = m
  \node[box] (m) at (5.0, -1.9) {$m = $ low 32\\of product};
  \draw[wire] (mul1.east) -- node[above=2pt, annot]{mod $2^{32}$} (m.west);

  % Step 4: sign-extend m, multiply by q
  \node[op] (mul2) at (5.0, -3.6) {$\times$};
  \node[box, minimum width=2.0cm] (q_) at (7.2, -3.6) {$q = 64513$};
  \draw[wire] (m.south) -- node[right=3pt, annot]{(64s)} (mul2.north);
  \draw[wire] (q_.west) -- (mul2.east);

  % Step 5: subtract
  \node[op] (sub) at (2.5, -3.6) {$-$};
  \draw[wire] (inp.south|-sub.north) ++(0,0) -- (sub.north)
    node[pos=0, below right=0pt]{};
  \coordinate (apath) at (0, -3.6);
  \draw[wire] (inp.south) -- (apath) -- (sub.west);
  \draw[wire] (mul2.west) -- node[above=2pt, annot]{$m \cdot q$} (sub.east);

  % Step 6: arithmetic right shift
  \node[box] (shift) at (2.5, -5.2) {$\gg_s\, 32$\\(arith shift)};
  \draw[wire] (sub.south) -- node[right=3pt, annot]{$a{-}mq$ (64-bit)} (shift.north);

  % Output
  \node[box, fill=col@algebraic, minimum width=3.2cm]
    (out) at (2.5, -6.9)
    {output $t$ (32-bit)\\$t \equiv a R^{-1} \pmod{q}$\\$|t| < q$};
  \draw[wire] (shift.south) -- node[right=3pt, annot]{divide by $R$} (out.north);

  % Correctness annotation
  \node[draw=OliveGreen!60, fill=OliveGreen!8, rounded corners=3pt,
        font=\scriptsize, inner sep=6pt, align=left]
    at (8.5, -5.5)
    {\textbf{Key invariant:}\\
     $m \cdot q \equiv -a \pmod{2^{32}}$\\
     $\Rightarrow (a - mq) \equiv 0 \pmod{2^{32}}$\\
     $\Rightarrow$ exact integer division\\[3pt]
     $t = \dfrac{a - mq}{R} \equiv aR^{-1}\!\!\pmod{q}$};

  % Title
  \node[font=\small\bfseries\sffamily, col@darkline] at (2.5, 0.7)
    {Signed Montgomery Reduction \texttt{montgomery\_reduce}$(a)$};

\end{tikzpicture}
\caption{Dataflow diagram of the signed Montgomery reduction algorithm.
         The two 32-bit multiplications (\texttt{QINV} and $q$) and the
         arithmetic right shift by 32 are the three core operations.
         The key divisibility invariant ensures that the right shift is
         exact (no remainder), giving the congruence $t \equiv a R^{-1}
         \pmod{q}$.}
\label{fig:montgomery_flow}
\end{figure}

%==============================================================================%
%  FIGURE 5 — Negacyclic Ring and NTT Decomposition
%==============================================================================%
\begin{figure}[htbp]
\centering
\begin{tikzpicture}[x=1cm, y=1cm, scale=0.85, transform shape]

  %--- Ring diagram on the left ------------------------------------------------
  \node[draw=col@darkline, fill=col@imperative!60, rounded corners=6pt,
        minimum width=5.8cm, minimum height=3.6cm, align=center,
        font=\small\sffamily, inner sep=8pt]
    (ring) at (-3.2, 0)
    {\textbf{Polynomial Ring}\\[4pt]
     $R_q = \mathbb{Z}_q[X] / (X^{256}+1)$\\[4pt]
     $q = 64513$,\quad $n = 256$\\[4pt]
     coefficient vector $\in \mathbb{Z}_q^{256}$\\[4pt]
     multiplication = \emph{negacyclic}\\
     convolution};

  %--- Evaluation diagram on the right -----------------------------------------
  \node[draw=col@darkline, fill=col@algebraic!60, rounded corners=6pt,
        minimum width=5.8cm, minimum height=3.6cm, align=center,
        font=\small\sffamily, inner sep=8pt]
    (nttdom) at (3.8, 0)
    {\textbf{NTT Domain}\\[4pt]
     evaluation vector $\in \mathbb{F}_q^{256}$\\[4pt]
     eval.\ points: $\zeta_0^{2\,\mathrm{br}(i)+1}$,\\
     $i=0,\dots,255$\\[4pt]
     multiplication = \emph{pointwise}\\
     $\hat f \odot \hat g$};

  %--- NTT arrow ---------------------------------------------------------------
  \draw[proof arrow, line width=2pt, OliveGreen!70!black]
    (ring.east) to[bend left=18]
    node[above=5pt, font=\small\sffamily\bfseries, OliveGreen!60!black]
      {$\mathsf{NTT}$}
    node[below=5pt, font=\scriptsize, col@darkline]
      {$O(n\log n)$, 8 butterfly stages}
    (nttdom.west);

  \draw[proof arrow, line width=2pt, BrickRed!70!black]
    (nttdom.west) to[bend left=18]
    node[below=5pt, font=\small\sffamily\bfseries, BrickRed!60!black]
      {$\mathsf{INTT}$}
    node[above=5pt, font=\scriptsize, col@darkline]
      {$+$ scaling by $256^{-1}\!\cdot R$}
    (ring.east);

  %--- Isomorphism label -------------------------------------------------------
  \node[font=\small, col@darkline, align=center] at (0.3, 1.5)
    {$\mathsf{NTT} \circ \mathsf{INTT} = \mathrm{id}$};

  %--- CRT decomposition -------------------------------------------------------
  \node[draw=col@darkline, fill=col@butterfly!40, rounded corners=4pt,
        font=\scriptsize\sffamily, inner sep=6pt, align=center]
    (crt) at (0.3, -3.2)
    {\textbf{CRT decomposition:}\enspace
     $R_q \cong \prod_{i=0}^{255} \mathbb{F}_q / (X - \zeta_0^{2\,\mathrm{br}(i)+1})$\\[3pt]
     $X^{256}+1 = \prod_{i=0}^{255}(X - \zeta_0^{2i+1})$ in $\mathbb{F}_q[X]$\\
     (valid since $\zeta_0^{512}=1$, $\zeta_0^{256}=-1$)};

  \draw[col@darkline, dashed, -Stealth] (ring.south) to[bend right=15] (crt.west);
  \draw[col@darkline, dashed, -Stealth] (nttdom.south) to[bend left=15] (crt.east);

  %--- Root circle diagram mini ------------------------------------------------
  \begin{scope}[shift={(0.3, 4.0)}, scale=0.7]
    \draw[col@darkline!50, thick] (0,0) circle (1.3cm);
    \foreach \k in {0,...,7} {
      \pgfmathsetmacro{\ang}{360/8*\k}
      \fill[BrickRed!60] ({1.3*cos(\ang)}, {1.3*sin(\ang)}) circle (2.5pt);
    }
    \node[font=\tiny, col@darkline] at (0,-1.8) {$512$th roots of unity (8 of 256 shown)};
    \draw[-Stealth, col@darkline!50] (-1.6, 0) -- (1.6, 0);
    \draw[-Stealth, col@darkline!50] (0, -1.6) -- (0, 1.6);
    \node[font=\tiny] at (1.45,0.18) {$1$};
    \fill[OliveGreen!80] ({1.3*cos(22.5)},{1.3*sin(22.5)}) circle (3pt);
    \node[font=\tiny, OliveGreen!70!black] at (1.7, 0.55) {$\zeta_0$};
  \end{scope}

\end{tikzpicture}
\caption{The negacyclic polynomial ring $R_q$ and its isomorphism to the NTT
         evaluation domain.  The NTT evaluates the polynomial at all 256
         primitive $512$th roots of unity (red dots on the unit circle).
         This is an instance of the Chinese Remainder Theorem decomposition
         of $R_q$ over $\mathbb{F}_q$.}
\label{fig:ring_ntt}
\end{figure}

%==============================================================================%
%  FIGURE 6 — Bit-Reversal Permutation and Twiddle Index Mapping
%==============================================================================%
\begin{figure}[htbp]
\centering
\begin{tikzpicture}[x=1.15cm, y=0.78cm, scale=0.72, transform shape]

  %--- Title ------------------------------------------------------------------
  \node[font=\small\bfseries\sffamily, col@darkline] at (7.5, 2.1)
    {Bit-reversal $\mathrm{brev}_8 : \{0,\dots,255\}\to\{0,\dots,255\}$
     — first 16 values};

  %--- Grid background (alternating columns) ----------------------------------
  \foreach \k in {0,2,...,14} {
    \fill[col@butterfly!22] (\k-0.45,-0.88) rectangle (\k+0.45,1.68);
  }

  %--- Hardcoded table rows ---------------------------------------------------
  % Row: k value
  \foreach \k/\v in
    {0/0,1/1,2/2,3/3,4/4,5/5,6/6,7/7,
     8/8,9/9,10/10,11/11,12/12,13/13,14/14,15/15}
  { \node[font=\footnotesize\ttfamily] at (\k,1.25) {\v}; }

  % Row: brev_8(k) — hardcoded
  \foreach \k/\v in
    {0/0,1/128,2/64,3/192,4/32,5/160,6/96,7/224,
     8/16,9/144,10/80,11/208,12/48,13/176,14/112,15/240}
  { \node[font=\footnotesize\ttfamily, col@highlight] at (\k,0.42) {\v}; }

  % Row: concrete jzetas[k] values (signed Montgomery form; index 0 is padding)
  \foreach \k/\v in
    {0/0,1/26964,2/-16505,3/22229,4/30746,5/20243,6/19064,7/-31218,
     8/9395,9/-30985,10/22859,11/-8851,12/32144,13/13744,14/21408,15/17599}
  { \node[font=\tiny\ttfamily, col@darkline] at (\k,-0.38) {\v}; }

  %--- Grid lines -------------------------------------------------------------
  \foreach \k in {0,...,16} {
    \draw[col@darkline!30, thin] (\k-0.45,-0.88) -- (\k-0.45,1.68);
  }
  \draw[col@darkline, thick] (-0.45,1.68) -- (15.45,1.68);
  \draw[col@darkline, thick] (-0.45,-0.88) -- (15.45,-0.88);
  \draw[col@darkline]        (-0.45,0.83)  -- (15.45,0.83);
  \draw[col@darkline]        (-0.45,0.02)  -- (15.45,0.02);

  %--- Row labels (left margin) -----------------------------------------------
  \node[font=\footnotesize\bfseries, col@darkline, anchor=east] at (-0.5,1.25)
    {$k$};
  \node[font=\footnotesize\bfseries, col@darkline, anchor=east] at (-0.5,0.42)
    {$\mathrm{brev}_8(k)$};
  \node[font=\footnotesize\bfseries, col@darkline, anchor=east] at (-0.5,-0.38)
    {\texttt{jzetas[k]}};

  %--- Bit-box example (k=3): 00000011 -> 11000000 = 192 ----------------------
  \node[font=\tiny, col@darkline, anchor=east] at (-0.05,-1.55) {$k{=}3$:};

  % k=3 bits b0..b7 = 1,1,0,0,0,0,0,0
  \fill[BrickRed!35,draw=col@darkline,rounded corners=1pt] (0.50,-1.68) rectangle (0.90,-1.35);
  \node[font=\tiny\ttfamily] at (0.70,-1.515) {1};
  \fill[BrickRed!35,draw=col@darkline,rounded corners=1pt] (0.94,-1.68) rectangle (1.34,-1.35);
  \node[font=\tiny\ttfamily] at (1.14,-1.515) {1};
  \fill[white,draw=col@darkline!50,rounded corners=1pt](1.38,-1.68) rectangle (1.78,-1.35);
  \node[font=\tiny\ttfamily] at (1.58,-1.515) {0};
  \fill[white,draw=col@darkline!50,rounded corners=1pt](1.82,-1.68) rectangle (2.22,-1.35);
  \node[font=\tiny\ttfamily] at (2.02,-1.515) {0};
  \fill[white,draw=col@darkline!50,rounded corners=1pt](2.26,-1.68) rectangle (2.66,-1.35);
  \node[font=\tiny\ttfamily] at (2.46,-1.515) {0};
  \fill[white,draw=col@darkline!50,rounded corners=1pt](2.70,-1.68) rectangle (3.10,-1.35);
  \node[font=\tiny\ttfamily] at (2.90,-1.515) {0};
  \fill[white,draw=col@darkline!50,rounded corners=1pt](3.14,-1.68) rectangle (3.54,-1.35);
  \node[font=\tiny\ttfamily] at (3.34,-1.515) {0};
  \fill[white,draw=col@darkline!50,rounded corners=1pt](3.58,-1.68) rectangle (3.98,-1.35);
  \node[font=\tiny\ttfamily] at (3.78,-1.515) {0};

  \node[font=\tiny, col@darkline, anchor=east] at (-0.05,-2.15)
    {\color{BrickRed!75!black}brev:};

  % brev_8(3)=192 bits b0..b7 = 0,0,0,0,0,0,1,1
  \fill[white,draw=col@darkline!50,rounded corners=1pt](0.50,-2.28) rectangle (0.90,-1.95);
  \node[font=\tiny\ttfamily] at (0.70,-2.115) {0};
  \fill[white,draw=col@darkline!50,rounded corners=1pt](0.94,-2.28) rectangle (1.34,-1.95);
  \node[font=\tiny\ttfamily] at (1.14,-2.115) {0};
  \fill[white,draw=col@darkline!50,rounded corners=1pt](1.38,-2.28) rectangle (1.78,-1.95);
  \node[font=\tiny\ttfamily] at (1.58,-2.115) {0};
  \fill[white,draw=col@darkline!50,rounded corners=1pt](1.82,-2.28) rectangle (2.22,-1.95);
  \node[font=\tiny\ttfamily] at (2.02,-2.115) {0};
  \fill[white,draw=col@darkline!50,rounded corners=1pt](2.26,-2.28) rectangle (2.66,-1.95);
  \node[font=\tiny\ttfamily] at (2.46,-2.115) {0};
  \fill[white,draw=col@darkline!50,rounded corners=1pt](2.70,-2.28) rectangle (3.10,-1.95);
  \node[font=\tiny\ttfamily] at (2.90,-2.115) {0};
  \fill[BrickRed!35,draw=col@darkline,rounded corners=1pt](3.14,-2.28) rectangle (3.54,-1.95);
  \node[font=\tiny\ttfamily] at (3.34,-2.115) {1};
  \fill[BrickRed!35,draw=col@darkline,rounded corners=1pt](3.58,-2.28) rectangle (3.98,-1.95);
  \node[font=\tiny\ttfamily] at (3.78,-2.115) {1};

  % Bit position labels b0..b7
  \foreach \b in {0,...,7} {
    \node[font=\tiny, col@darkline!60] at (0.70+\b*0.44,-2.52) {$b_{\b}$};
  }

  % Annotation box
  \node[draw=BrickRed!55, rounded corners=3pt, fill=BrickRed!6,
        font=\scriptsize, inner sep=6pt, align=center]
    (bitann) at (12.5,-1.85)
    {\textbf{Example:}\\
     $k=3 = 00000011_2$\\
     $\mathrm{brev}_8(3) = 11000000_2 = 192$\\[3pt]
     bits: $b_7b_6\cdots b_0 \;\mapsto\; b_0b_1\cdots b_7$};
  \draw[Stealth-, BrickRed!70!black, thick]
    (4.0,-2.0) to[bend right=20] (bitann.west);

\end{tikzpicture}
\caption{The bit-reversal permutation $\mathrm{brev}_8$ for the first 16 table
         values $k = 0,\dots,15$.  Red cells in the bit-pattern boxes highlight
         the set bits.  The bottom row gives the concrete signed
         Montgomery-form table entry \texttt{jzetas[k]}; index 0 is padding,
         and the forward NTT consumes indices 1 through 255.}
\label{fig:brev}
\end{figure}

%==============================================================================%
%  FIGURE 7 — EasyCrypt Proof Structure (Hoare / pRHL judgement stack)
%==============================================================================%
\begin{figure}[htbp]
\centering
\begin{tikzpicture}[x=1cm, y=1cm, scale=0.86, transform shape,
  triple/.style = {
    rectangle, draw=col@darkline, rounded corners=4pt,
    minimum width=12cm, minimum height=1.3cm,
    align=center, font=\small\ttfamily,
    fill=#1, inner sep=6pt,
  },
  hltag/.style = {
    rectangle, draw=col@highlight!70, rounded corners=3pt,
    fill=col@highlight!10, font=\footnotesize\sffamily,
    inner sep=4pt,
  },
]

  %--- Three levels of Hoare judgements ----------------------------------------
  \node[triple=col@jasmin!70]
    (J) at (0, 5.8)
    {%
     \textbf{pRHL equivalence (RefJasminNTT.ec):}\\
     $\texttt{equiv}[\texttt{Hpoly\_extract.M.\_poly\_ntt} \sim
       \texttt{NTT\_Fq.NTT.ntt} :$\\
     $\texttt{poly\_repr\_bound}(\texttt{rp},p,16)
     \;\Longrightarrow\;
     \texttt{poly\_repr\_bound}(\texttt{rp}',p',24)]$};

  \node[triple=col@imperative!70]
    (I) at (0, 3.1)
    {%
     \textbf{Hoare triple (NTT\_Fq.ec):}\\
     $\texttt{hoare}[\texttt{NTT.ntt} :
       \texttt{arg}=(p,\texttt{zetas})$\\
     $\Longrightarrow\;
     \texttt{res} = \texttt{NTT\_loop}(p)]$};

  \node[triple=col@algebraic!70]
    (A) at (0, 0.4)
    {%
     \textbf{Algebraic equality (NTTFullAlgebra.ec):}\\
     $\forall p \in \mathbb{Z}_q^{256}.\;
       \texttt{NTT\_loop}(p) = \texttt{full\_ntt}(p)$\\
     where $\texttt{full\_ntt}(p)[i] =
       \sum_{j=0}^{255} p[j] \cdot \zeta_0^{(2\,\mathrm{brev}_8(i)+1)\cdot j}$};

  \node[triple=col@endtoend, minimum height=1.0cm]
    (E) at (0, -1.8)
    {%
     \textbf{End-to-end (NTTEndToEnd.ec):}\\
     $\texttt{hoare}[\texttt{poly\_ntt\_jazz} :
       \texttt{poly\_repr\_bound}(\texttt{rp},p,16)
       \;\Longrightarrow\;
       \texttt{poly\_repr\_bound}(\texttt{res},\texttt{full\_ntt}(p),24)]$};

  %--- Composition arrows ------------------------------------------------------
  \draw[proof arrow, line width=1.8pt, OliveGreen!70!black]
    (J.south) -- node[right=5pt, annot, OliveGreen!70!black]
      {Step B: apply pRHL $+$ representation} (I.north);

  \draw[proof arrow, line width=1.8pt, OliveGreen!70!black]
    (I.south) -- node[right=5pt, annot, OliveGreen!70!black]
      {Step C: substitute algebraic identity} (A.north);

  \draw[proof arrow, line width=2.2pt, BrickRed!70!black]
    (A.south) -- node[right=5pt, annot, BrickRed!70!black]
      {compose A, B, C} (E.north);

  %--- Side tags ---------------------------------------------------------------
  \node[hltag, anchor=west] at (6.4, 5.8)
    {pRHL};
  \node[hltag, anchor=west] at (6.4, 3.1)
    {HL};
  \node[hltag, anchor=west] at (6.4, 0.4)
    {algebra};
  \node[hltag, anchor=west, fill=BrickRed!12, draw=BrickRed!60]
    at (6.4, -1.8)
    {top-level};

  %--- Proof state tags -------------------------------------------------------
  \node[annot, anchor=east] at (-6.2, 5.8) {machine\\words};
  \node[annot, anchor=east] at (-6.2, 3.1) {field\\arrays};
  \node[annot, anchor=east] at (-6.2, 0.4) {field\\elements};
  \node[annot, anchor=east] at (-6.2,-1.8) {full\\statement};

\end{tikzpicture}
\caption{Structure of the \easycrypt{} proof judgements.  The topmost
         judgement (pRHL) relates machine-word arrays to field-coefficient
         arrays with signed bit-width bounds.  The middle (Hoare) judgement
         captures the loop semantics over field arrays.  The algebraic equality connects the loop output to the
         closed-form spectral sum.  The bottom (end-to-end) judgement
         composes all three.}
\label{fig:ec_proofs}
\end{figure}

%==============================================================================%
%  FIGURE 8 — HAETAE vs. MLKEM NTT Structural Comparison
%==============================================================================%
\begin{figure}[htbp]
\centering
\begin{tikzpicture}[x=1cm, y=1cm, scale=0.60, transform shape,
  scheme/.style = {
    rectangle, draw=col@darkline, rounded corners=6pt,
    minimum width=5.5cm, inner sep=10pt,
    align=left, font=\small\sffamily,
  },
  diff/.style  = {
    fill=BrickRed!12, draw=BrickRed!60, rounded corners=3pt,
    font=\small\sffamily,
  },
  same/.style  = {
    fill=OliveGreen!12, draw=OliveGreen!60, rounded corners=3pt,
    font=\small\sffamily,
  },
]

  %--- HAETAE box --------------------------------------------------------------
  \node[scheme, fill=col@jasmin!50]
    (H) at (-3.2, 0)
    {%
     \textbf{\haetae{} NTT}\\[6pt]
     Modulus $q = 64513$\\
     Degree $n = 256$\\
     Prim.\ root $\zeta_0 = 426$ (512th)\\
     NTT length: \textbf{256-point}\\
     Twiddle exponent: $\mathrm{brev}_8(k)$\\
     Groups / stage: $2^s$, $s=0..7$\\
     Output: fully evaluated\\
     Algebra: $\mathrm{brev}_8$ lemmas\\
     Orthog.\ over 256 roots
    };

  %--- MLKEM box ---------------------------------------------------------------
  \node[scheme, fill=col@algebraic!50]
    (M) at (3.2, 0)
    {%
     \textbf{MLKEM (Kyber) NTT}\\[6pt]
     Modulus $q = 3329$\\
     Degree $n = 256$\\
     Prim.\ root $\zeta_0 = 17$ (256th)\\
     NTT length: \textbf{128-point}\\
     Twiddle exponent: $\mathrm{brev}_7(k)$\\
     Groups / stage: $2^s$, $s=0..6$\\
     Output: incomplete (pairs)\\
     Algebra: $\mathrm{brev}_7$ lemmas\\
     Orthog.\ over 128 roots
    };

  %--- Central "vs" -----------------------------------------------------------
  \node[circle, draw=col@darkline, fill=white,
        minimum size=0.9cm, font=\small\bfseries\sffamily]
    at (0, 0) {vs};

  %--- Difference annotation ---------------------------------------------------
  \node[diff, minimum width=11cm, inner sep=7pt, align=center]
    (diff) at (0, -4.4)
    {\textbf{Key structural difference (invalidates direct algebra reuse):}\\[3pt]
     \haetae{}: twiddle $= \zeta_0^{\mathrm{brev}_8(2^s + g)}$
     \enspace\textbf{(8-bit reversal, full 256-point)}\\
     MLKEM:\enspace twiddle $= \zeta_0^{\mathrm{brev}_7(2g+1)}$
     \enspace\textbf{(7-bit reversal, incomplete 128-point)}};

  \draw[col@darkline, dashed, -Stealth] (H.south) to[bend right=15] (diff.west);
  \draw[col@darkline, dashed, -Stealth] (M.south) to[bend left=15] (diff.east);

  %--- Similarity annotation ---------------------------------------------------
  \node[same, minimum width=11cm, inner sep=7pt, align=center]
    (same) at (0, 3.6)
    {\textbf{Shared aspects (methodology transferable):}
     signed Montgomery REDC,
     Jasmin extraction,
     pRHL refinement pattern,
     coefficient bound tracking};

  \draw[col@darkline, dashed] (H.north) to[bend left=15] (same.west);
  \draw[col@darkline, dashed] (M.north) to[bend right=15] (same.east);

\end{tikzpicture}
\caption{Structural comparison of the \haetae{} and MLKEM NTT\@.
         Green box: aspects where methodology is transferable between
         the two verification projects.  Red box: the critical structural
         difference that prevents direct reuse of MLKEM algebra lemmas.}
\label{fig:haetae_vs_mlkem}
\end{figure}

%==============================================================================%
%  FIGURE 9 — Coefficient Bound Propagation Through All Layers
%==============================================================================%
\begin{figure}[htbp]
\centering
\begin{tikzpicture}[x=0.94cm, y=1cm]

  % Column x positions
  \def\xA{0.0}  % layer name
  \def\xB{2.8}  % forward NTT bound
  \def\xC{6.5}  % inverse NTT bound
  \def\xD{9.8}  % round-trip bound

  %--- Header -----------------------------------------------------------------
  \node[font=\small\bfseries\sffamily, anchor=west] at (\xA, 4.0) {Layer};
  \node[font=\small\bfseries\sffamily, anchor=center] at (\xB, 4.0)
    {Forward NTT bounds};
  \node[font=\small\bfseries\sffamily, anchor=center] at (\xC, 4.0)
    {Inverse NTT bounds};
  \node[font=\small\bfseries\sffamily, anchor=center] at (\xD, 4.0)
    {Identity note};

  \draw[col@darkline, thick] (\xA-0.1, 3.65) -- (\xD+1.5, 3.65);

  %--- Rows (hardcoded 5 rows via global \bndrow) ------------------------------
  \bndrow{3.10}{1}{Jasmin (W32)}%
    {16-bit repr. $\to$ 24-bit repr.}%
    {16-bit repr. $\to$ 16-bit Mont. repr.}%
    {if precondition holds}

  \bndrow{2.25}{0}{Extracted EC}%
    {\texttt{repr\_bound} 16 $\to$ 24}%
    {\texttt{repr\_bound} 16 $\to$ Mont. 16}%
    {\texttt{repr}}

  \bndrow{1.40}{1}{\texttt{NTT\_Fq.ec}}%
    {$p \to \texttt{full\_ntt}(p)$}%
    {$p \to \texttt{full\_invntt}(p)$}%
    {separate maps}

  \bndrow{0.55}{0}{\texttt{NTTFullSpec}}%
    {$p\in\mathbb{Z}_q^{256} \to \hat{p}$}%
    {$\hat{p} \to p$}%
    {math identity}

  \bndrow{-0.30}{1}{\texttt{NTTEndToEnd}}%
    {pre-cond $\to$ post-cond}%
    {pre-cond $\to$ post-cond}%
    {no pipeline lemma}

  %--- Vertical dividers -------------------------------------------------------
  \foreach \x in {2.2, 5.2, 8.8} {
    \draw[col@darkline!40] (\x, 3.65) -- (\x, -1.3);
  }
  \draw[col@darkline, thick] (\xA-0.1, 3.65) -- (\xA-0.1, -1.3);
  \draw[col@darkline, thick] (\xD+1.5, 3.65) -- (\xD+1.5, -1.3);
  \draw[col@darkline, thick] (\xA-0.1, -1.3) -- (\xD+1.5, -1.3);

  %--- Column header rules ----------------------------------------------------
  \draw[col@darkline] (\xA-0.1, 3.65) -- (\xD+1.5, 3.65);

  %--- Annotation --------------------------------------------------------------
  \node[draw=col@darkline, rounded corners=3pt, fill=white,
        font=\scriptsize, inner sep=5pt, align=left]
    at (5.7, -2.4)
    {\textbf{Bit-width summary (forward NTT):}\\[2pt]
     input $16\,\text{bit} \xrightarrow{\text{stage 1}} 17 \xrightarrow{2} 18
       \xrightarrow{3} 19 \xrightarrow{4} 20$\\
     $\xrightarrow{5} 21
       \xrightarrow{6} 22 \xrightarrow{7} 23
       \xrightarrow{\text{post}} 24\,\text{bit}$};

\end{tikzpicture}
\caption{Bound propagation for the checked forward and inverse NTT theorems.}
\label{fig:bounds_table}
\end{figure}

%==============================================================================%
%  FIGURE 10 — EasyCrypt Module/Theory Dependency Graph (detailed)
%==============================================================================%
\begin{figure}[htbp]
\centering
\begin{tikzpicture}[scale=0.90, transform shape,
  x=2.6cm, y=1.6cm,
  ecmod/.style = {
    rectangle, draw=col@darkline, rounded corners=4pt,
    minimum width=2.3cm, minimum height=0.75cm,
    align=center, font=\ttfamily\footnotesize,
    fill=#1,
  },
  ecmod/.default = white,
  depar/.style = {-Stealth, col@arrow, thin},
]

  % Row 0 – foundational theories
  \node[ecmod=col@extract]      (Fq)   at (0,0)   {Fq.ec};
  \node[ecmod=col@extract]      (GFq)  at (1,0)   {GFq.ec};
  \node[ecmod=col@extract]      (FE)   at (2,0)   {Fastexp.ec};
  \node[ecmod=col@extract]      (Arr)  at (3,0)   {Array256.ec};
  \node[ecmod=col@extract]      (Mont) at (4,0)   {Montgomery.ec};

  % Row 1 – ring / word
  \node[ecmod=col@extract!60]   (Rq)   at (1,1)   {Rq.ec};
  \node[ecmod=col@extract!60]   (BA)   at (3,1)   {BArray1024.ec};
  \node[ecmod=col@extract!60]   (WA)   at (4,1)   {WArray1024.ec};

  % Row 2 – extracted and imperative spec
  \node[ecmod=col@imperative!80](Ext)  at (0,2)   {Hpoly\_extract.ec};
  \node[ecmod=col@imperative!80](Nfq)  at (2,2)   {NTT\_Fq.ec};

  % Row 3 – full spec and algebra
  \node[ecmod=col@algebraic!80] (FS)   at (1,3)   {NTTFullSpec.ec};
  \node[ecmod=col@algebraic!80] (FA)   at (3,3)   {NTTFullAlgebra.ec};

  % Row 4 – refinement
  \node[ecmod=col@jasmin!80]    (Ref)  at (2,4)   {RefJasminNTT.ec};

  % Row 5 – end-to-end
  \node[ecmod=col@endtoend, minimum width=3.5cm]
                                (E2E)  at (2,5)   {NTTEndToEnd.ec};

  % Edges
  \draw[depar] (Fq)   -- (Nfq);
  \draw[depar] (GFq)  -- (Nfq);
  \draw[depar] (FE)   -- (Nfq);
  \draw[depar] (Arr)  -- (Nfq);
  \draw[depar] (Mont) -- (Fq);
  \draw[depar] (GFq)  -- (Rq);
  \draw[depar] (Arr)  -- (Rq);
  \draw[depar] (Arr)  -- (BA);
  \draw[depar] (BA)   -- (Ext);
  \draw[depar] (WA)   -- (Ext);
  \draw[depar] (Ext)  -- (Fq);
  \draw[depar] (Rq)   -- (FS);
  \draw[depar] (GFq)  -- (FS);
  \draw[depar] (Arr)  -- (FS);
  \draw[depar] (FS)   -- (FA);
  \draw[depar] (Nfq)  -- (FA);
  \draw[depar] (Fq)   -- (Ref);
  \draw[depar] (Nfq)  -- (Ref);
  \draw[depar] (Ext)  -- (Ref);
  \draw[depar] (Ref)  -- (E2E);
  \draw[depar] (FA)   -- (E2E);
  \draw[depar] (FS)   -- (E2E);
  \draw[depar] (Nfq)  -- (E2E);
  \draw[depar] (Ext)  -- (E2E);

  %--- Legend -----------------------------------------------------------------
  \matrix[draw=col@darkline, rounded corners=3pt, fill=white,
          row sep=3pt, column sep=5pt, inner sep=6pt,
          anchor=north west, font=\scriptsize\sffamily]
    at (4.5, 5.2)
  {
    \node[ecmod=col@extract,    minimum width=1.5cm,minimum height=0.45cm]
      {Foundations}; \\
    \node[ecmod=col@imperative!80, minimum width=1.5cm,minimum height=0.45cm]
      {Imperative}; \\
    \node[ecmod=col@algebraic!80,  minimum width=1.5cm,minimum height=0.45cm]
      {Algebraic}; \\
    \node[ecmod=col@jasmin!80,     minimum width=1.5cm,minimum height=0.45cm]
      {Refinement}; \\
    \node[ecmod=col@endtoend,      minimum width=1.5cm,minimum height=0.45cm]
      {End-to-end}; \\
  };

\end{tikzpicture}
\caption{Detailed \easycrypt{} module dependency graph.  Arrows point from
         dependency to dependent file; standard-library and some transitive
         edges are omitted for readability.  Nodes are colour-coded by layer:
         orange = foundations, yellow = imperative specification,
         green = algebraic specification, red = refinement, cyan = end-to-end.}
\label{fig:dep_graph_full}
\end{figure}


\backmatter
\printbibliography[heading=bibintoc, title={Bibliography}]

\end{document}
